%
% ############################################################
% ##  REVIEW NOTE — TWO PASSAGES BELOW ARE KNOWN STALE       ##
% ##  Queued for revision AFTER the laboratory boredom       ##
% ##  section is drafted (author ruling 2026-07-29).         ##
% ############################################################
%
% WHY: both passages cite the laboratory keystone in its SUPERSEDED
%   form. The O-9 reprocessing (2026-07-16, N=8) withdrew it. P1's
%   "clinical EEG ~ boring" came from its ONE clean subject and does
%   NOT replicate: clinical DMN is LOW (engaging-like), and the
%   boring-like reading holds in 4/8 subjects = chance.
%   CORRECTED KEYSTONE: all-physiology-engaged vs self-report-bored
%   (clinical self-report 6.0/9). Only pupil dilation is significant
%   at N=8 (Boring < Clinical in all 8, p=0.008); every other surface
%   is directional only. Report convergence + pupil, never per-channel
%   significance.
%
% LAYER ASSIGNMENT (author-ruled 2026-07-29 — governs the rewording):
%   physiology AS A WHOLE = the OCCUPIED SURFACE (outward comportment);
%   the post-video RETROSPECTIVE self-report = the HOLLOW DEPTH.
%   Grounded in fcm-05 SS24-28: at the dinner party we are "totally
%   involved" and DENY boredom, recognizing it only afterward;
%   "behavioural engagement is not evidence against boredom; in this
%   form it is the very medium of it." NB the old mapping made a brain
%   measure the DEPTH channel, which fcm-05 C13 already forbade.
%
% THE TWO SITES (search the strings; line numbers drift):
%   [STALE-1]  M3.6  -- "the laboratory's physiological record can
%              disagree with its behavioral one" AND "set a brain
%              measure against a gaze measure". BOTH clauses are stale:
%              under the corrected layers, physiological and behavioural
%              are the SAME layer (both surface).
%   [STALE-2]  M4.1 P6 -- "a brain measure can read boring-like on the
%              very footage a gaze measure reads as engaged".
%
% WHAT SURVIVES UNCHANGED AT BOTH SITES: the structural point -- a
%   survey holds ONE channel; the laboratory holds TWO. Only the NAMING
%   of the two channels is wrong. These are small edits, not rewrites.
%
% *** THIS FILE IS GENERATED. *** It is flattened from the \input-based
%   DESKTOP-SECTION-ASSEMBLED-v1.tex; the canonical texts are
%   DESKTOP-M3-DRAFT-v1.tex and DESKTOP-M4-DRAFT-v1.tex, which carry the
%   same note in their headers. EDIT THOSE, then re-flatten -- which
%   overwrites this file and removes these markers. That is correct and
%   expected: the markers exist only until the revision is made.
%
% CANONICAL SOURCES WHEN THE TIME COMES:
%   corpus/index/Heidegger - The Fundamental Concepts of Metaphysics/
%     _synthesis/fcm-king-salvo-bridge.md   (S4, revised 2026-07-29)
%   tmp/Dissertation/Part_III/reanalysis/boredom-o9-physiological-findings.md
% ############################################################
%
% ============================================================
% Part III — Desktop Section — ASSEMBLED DRAFT v1 (2026-07-29)
% SELF-CONTAINED OVERLEAF VERSION (flattened from the \input-based
% assembled file 2026-07-29 — module files remain canonical; re-flatten
% after any module edit). Compile with XeLaTeX.
% Assembly pass per HANDOFF-DESKTOP-SECTION-DRAFTING-COMPLETE-2026-07-28.md
% §5 items 1–5 (user-approved 2026-07-29). Backup .backups/20260729T075642/.
%
% STRUCTURE: \input of the canonical per-module files (they remain the
%   source of record; edit THEM, not copies here). M0–M3 carry their own
%   \section commands; M4–M7 do not — the ASSEMBLY-PROVISIONAL heads below
%   supply them, subject to user approval and the arc decision (Open
%   decision #1). M4's outline title contained "Stalled Chain" (no-stall
%   doctrine) — a neutral head is used pending user ruling.
%
% CARRIED MARKERS (do not resolve here):
%   - M4.4 REVISION-PASS: blocked on Part II closing differential.
%   - M7.2 O-11 EXPANSION MARKER: section-scale close only.
%   - M7.1 ↔ III-5 three-state HARMONIZATION: blocked on arc decision.
%   - M2.3 not drafted (arc decision). M2.1 flag-1 procedure list OPEN.
%   - M0: class name/number (******) + Berkeley/Merced enrollment follow-up.
%
% ASSEMBLY-PASS RULINGS — ALL APPLIED 2026-07-29 (user-ordered):
%   (a) M0 "stall" x4 fixed (thesis -> "the world falls silent at the
%       passage"; close -> "the act was withheld"; frame/resolution ->
%       "fall(s) silent"). (b) M3.3 "priced/price" -> "weighed / the
%       effort it asked was too great". (c) M3.4 bare "incorporation" x2
%       -> "rhetorical incorporation" (third flagged hit was already
%       "Rhetorical incorporation" across a line break). (d) M3 title ->
%       "The Deployment Findings". (e) M5 \textit{Rhetoric} x2 ->
%       \textit{Rhetorica}. (f) M7.1 serial comma fixed.
%   M2.1 PROCEDURE LIST INTEGRATED (author-supplied 2026-07-29; flag 1
%   resolved — see M2 header for the confirm-at-review notes).
% ASSEMBLY TOUCHES MADE (banked conditionals): M5.1 first D&L cite →
%   ("Learning Affordances" 25) [short-title conditional fired: 2012
%   follow-up is in Works Cited]; M1 Parong & Mayer 2021 short title
%   normalized → "Cognitive and Affective Processes" (per M7 header).
%   Coding-methods appendix A-family total corrected 155 → 194 (verified
%   vs t3-stats). W26-R037 fidelity sweep: all rendered uses correct.
% ============================================================
\documentclass{article}
% >>> COMPILE WITH XeLaTeX <<<
\usepackage{graphicx}
\usepackage{xcolor}
\usepackage{amssymb}
\usepackage{newunicodechar}
\usepackage{fontspec}
\usepackage{polyglossia}
\setmainlanguage{english}
\setotherlanguage{greek}
\usepackage{geometry}
\geometry{margin=1in}
\usepackage{csquotes}
\usepackage{enumitem}
\usepackage{setspace}

\newfontfamily\greekfont{Gentium Plus}
\newcommand{\gk}[1]{{\greekfont #1}}

\newunicodechar{₀}{\ensuremath{_0}}
\newunicodechar{₁}{\ensuremath{_1}}
\newunicodechar{₂}{\ensuremath{_2}}
\newunicodechar{₃}{\ensuremath{_3}}
\newunicodechar{₄}{\ensuremath{_4}}
\newunicodechar{→}{\ensuremath{\rightarrow}}

\usepackage[colorlinks=true, linkcolor=blue, urlcolor=blue, citecolor=blue]{hyperref}

% Works Cited hanging indent
\newenvironment{workscited}
  {\section*{Works Cited}\begin{list}{}
    {\setlength{\itemindent}{-0.5in}\setlength{\leftmargin}{0.5in}
     \setlength{\itemsep}{0pt}\setlength{\parsep}{6pt}}}
  {\end{list}}

\title{Part III --- The Desktop Deployment \\
  \large ASSEMBLED SECTION DRAFT v1 (M0--M7, 2026-07-29) ---
  section title provisional pending arc decision}
\author{}
\date{2026-07-29}

\begin{document}
\maketitle

% ============================================================
% REMAINING USER-GATED ITEMS after the 2026-07-31 full readthrough
% (all eight subsections revised batch-by-batch; intro/thesis/close
% alignment pass EXECUTED; committed on user order 2026-07-31):
%   1. USER MANUAL REWRITES (markers in the module files):
%      - M2 coding paragraph ("Responses were anonymized...") + its
%        coding-passes footnote (DESKTOP-M2-DRAFT-v1.tex).
%      - M3.6 keyword-sweep footnote ("The counts come from a keyword
%        sweep...") (DESKTOP-M3-DRAFT-v1.tex).
%   2. FIGURE PENDING (rendered placeholder at M3.5): screenshot of
%      multiple students meeting inside the VLE.
%   3. WORKS CITED verification pins: Boellstorff 2023 final vol/no/
%      pages (****** markers kept per user order); Mark publisher/
%      subtitle; Tu venue; Parong & Mayer 2018 pagination; Aristotle/
%      Heidegger entries harmonize vs Part I at dissertation assembly.
%   4. BLOCKED (do not start): M4.4 revision pass <- Part II closing
%      differential; M2.3 + O-11 + three-state harmonization <- arc
%      decision <- lab boredom section.
% ============================================================
% ===== M0 Introduction (carries \section{From the Pilot to the Deployment})
% ========== BEGIN INLINED FILE: DESKTOP-M0-DRAFT-v2.tex ==========
% ============================================================
% Part III — Desktop Section — M0 (Introduction) — DRAFT v2 (2026-07-20)
% RESTRUCTURED to the author's six-move introduction structure (FCDP §10.1;
% memory feedback-user-intro-structure): hook > context > disruption >
% resolution > signpost > thesis. v1 preserved as DESKTOP-M0-DRAFT-v1.tex
% (major iteration on a copy).
% Move map: hook = NEW (headset shelf + two K&S verbatims) · context = M0.1 +
%   NEW frame passage + M0.2 · disruption = NEW (rhetorical question) ·
%   resolution = NEW (gap + frame-as-resolution, methodology folded) ·
%   signpost = M0.4 (GOLD STANDARD roadmap, untouched) · thesis = M0.3
%   (opener retrofit only).
% Conventions carried from v1 (MLA; ("Phenomenological Evaluation" NN);
%   first-person for K&S work; "I" for dissertation-argument claims; bans:
%   "record", sentence-initial "And", metatextual signposting, gambling
%   vocabulary, bare "incorporation" for our theory [always "rhetorical
%   incorporation"; bare form = Calleja's theory only]).
% Verified quote pages (PDF): indisputable-presence 390; not-a-single-student
%   389; "a hassle" (results wording) 388; hypothesis + unlikely 383;
%   restraint 390; no-collaboration 389; figures 388; reasons 389; scarcity
%   381; phenomenological framing 387; Heim AWS 182/185; Radianti 1/22.
% NO-STALL SWEEP APPLIED 2026-07-29 (user-ordered at assembly; doctrine
%   postdated M0 approval): "stall at any station" -> "fall silent at any
%   station" (frame passage); "lets it stall" -> "falls silent before it"
%   (resolution); thesis "a stall at the passage" -> "the world falls
%   silent at the passage" (M7.2's approved language); "completion
%   stalled" -> "the act was withheld" (thesis close).
% DOCTRINE (user 2026-07-31, readthrough batch 1): never phrase the
%   A3->A4 diagnosis as "the space provides no act" or "students get
%   stalled at A3" — a potentiality to actualize is always present
%   (watching is itself an actualization); the failure is in the WORTH of
%   the potentialities the world furnishes. Design's job = make them
%   worthwhile.
% REVISION-PASS MARKERS: (a) social-interior/Second Life passage — revise
%   against the new-media/presence sources once ingested + indexed (agent
%   running); (b) DONE 2026-07-31: roadmap re-verified vs final section,
%   fourth-part description updated (boredom three-form reading named);
%   (c) EXECUTED 2026-07-31: intro/thesis/close alignment pass complete —
%   thesis and M7.2 close now state and demonstrate the same three-part
%   claim in the same vocabulary.
% ============================================================

\section{From the Pilot to the Deployment}

% === (I) HOOK — [REVIEW: v3 this round; K&S verbatims added per user note.
%     Full-citation/convention footnote RELOCATED here from M0.1.]

In the spring of 2022, the science library at the University of California,
Irvine kept a shelf of Meta Quest 2 headsets on loan to the students of a
biomedical-engineering elective.\footnote{Part III draws on several of my
previous publications; the article referenced here---which is reproduced in
Appendix~\ref{app:vr-pilot}---is the following, and will be cited
parenthetically by abbreviated title and page number hereafter: King, Christine
E., and Dalton Salvo. ``Phenomenological Evaluation of an Undergraduate Clinical
Needs Finding Skills Through a Virtual Reality Clinical Immersion Platform.''
\textit{Biomedical Engineering Education}, vol.~4, 2024, pp.~381--97,
https://doi.org/10.1007/s43683-024-00139-5.} Any student could borrow one at no
cost and, through it, stand inside an operating room---a space few
undergraduates will ever be admitted to. It is, we would later conclude,
``indisputable that a majority of our students felt that VR increased their
sense of presence and immersion'' (``Phenomenological Evaluation'' 390). Yet
``[w]hile we had several Meta Quest 2's housed in the school's library, not a
single student went there to check one out'' (389). That untouched shelf is this
section's problem in miniature: a virtual world can offer presence and still
lose its students to the more convenient copy of itself. What follows examines
that problem at scale---three quarters of the same course, redeployed on
ordinary computers to 241 students across the University of California---and
asks what a virtual learning environment (VLE) must offer, beyond presence, to
be worth entering.
% [Hook follow-up APPROVED 2026-07-20.]

% === (II) CONTEXT — part 1: the pilot (M0.1, locked text; seam edits marked)

Clinical immersion for non-medical staff within operating rooms and other
restricted medical spaces is scarce by structure. The programs that admit
undergraduates to such clinical settings draw substantially more applicants than
they can serve, so that a majority of the students who seek a clinical-immersion
experience will not have one---an inequity the biomedical-engineering literature
documents against a clear demand.\footnote{Guilford, Kotche, and Schmedlen,
surveying clinical-immersion experiences across biomedical-engineering programs,
find that only 48\% of applicants to clinical-immersion courses and 24\% of
applicants to immersion programs are able to participate annually (113--22).
Moravec et al.\ report on one such program---a short-term clinical immersion for
undergraduate biomedical engineering students---and the best practices such
programs have developed under the same constraint of scarce placements and clear
demand (217--23).} The virtual clinical-immersion platform whose deployment this
section reads was built against that scarcity: filmed operations, physician and
engineer interviews, and explorable procedure rooms, delivered so that any
enrolled student could be emplaced where only a limited few traditionally have
access. In its first published test, a pilot cohort of twenty-two
% [SEAM EDIT: "Its first published test came in the spring of 2022, when a
%  pilot cohort..." — the hook now carries spring 2022; wording adjusted.]
met the same operating-room footage in four competing formats: 2D video on
YouTube, an interactive computer build,\footnote{The interactive computer build
was developed in Unreal Engine~4 for standard PCs: essentially the same
explorable virtual environment as the Meta Quest 2 build, but navigated by mouse
and keyboard on a flat monitor, and delivering the traditional 2D videos rather
than the stereoscopic footage (``Phenomenological Evaluation'' 383).} a Google
Cardboard head mounted display (HMD),\footnote{Google Cardboard is a minimal
head-mounted display: a folded cardboard or plastic viewer holding two lenses,
into which the user inserts a smartphone. The phone's screen, split into a
stereoscopic pair and viewed through the lenses, produces a three-dimensional
image, and the phone's own motion sensors track the rotation of the head. In
this course the Cardboard route delivered the operating-room footage through
YouTube's VR mode---video only, with none of the interactive environment of the
PC and Meta Quest 2 builds (``Phenomenological Evaluation'' 383). YouTube has
since discontinued this VR viewing mode, and the route the Cardboard format
relied on is no longer available. Given the nature of this setup, of all VR methods it is the most prone to producing
Alternate World Syndrome (AWS): motion sickness, physical disorientation, and
physical discomfort. In Michael Heim's definition, ``AWS is the relativity
sickness that comes from switching back-and-forth between the primary and
virtual worlds'' (\textit{Virtual Realism} 182); the VR experience ``can mix
images and expectations from alternate worlds so as to distort our perception of
the primary world, making us prone to errors in mismatched contexts'' (185).}
and a Meta Quest 2 virtual reality (VR) HMD. Interestingly, across twenty
post-course responses we discovered a two-part verdict that would not line up.
The highest level of presence\footnote{We do not stipulate a formal definition
of presence in the article; it uses the term in the sense standard to the
VR-pedagogy literature it draws on---the felt sense of being physically present
in the virtual rather than the actual environment---and treats it, alongside
embodiment, as one of the ``intensely subjective experiential phenomena'' a
phenomenological method is fitted to assess (``Phenomenological Evaluation''
382, 387). The construct is anchored to Makransky and Lilleholt (2018) and
Makransky, Terkildsen, and Mayer (2019) and measured throughout by
self-report.} went to the headset: 70\% reported that the Meta Quest
VR version elicited a greater felt sense of presence. Student preference, on the
other hand, went the other way: 40\% judged the traditional 2D YouTube
videos most educationally beneficial. The two VR HMD formats together drew only
that same share---20\% choosing the Google Cardboard and 20\% the Meta Quest
2---while 10\% selected the PC build and 20\% judged all versions equally
beneficial (``Phenomenological Evaluation''
388). The students' explanation as to why the 2D videos were perceived as being
more educationally beneficial was largely grounded in physical
discomfort\footnote{The discomfort belonged overwhelmingly to the Google
Cardboard rather than the Meta Quest 2: asked which version was \textit{least}
educationally beneficial, 40\% of students named the Cardboard against only 5\%
for the Quest (``Phenomenological Evaluation'' 388). Among the students who
judged the 2D videos most beneficial, 38.9\% cited physical
discomfort---disorientation, the weight of the headset over long viewing
periods, the difficulty of wearing glasses inside an HMD---a figure that pools
the two VR formats (389). Our own judgment in the article is that some of these
issues ``would have been less relevant'' had students used the Meta Quest 2
rather than the Cardboard, ``which is a relatively horrible experience by
comparison'' (390).} and convenience.\footnote{Among the students who judged the
2D videos most beneficial, 44.5\% named ease of use as the ground: convenience
and accessibility of access, and the quality-of-life controls the VR formats
lacked---variable playback speed, better volume and playback control, and the
ability to zoom in on the footage (``Phenomenological Evaluation'' 389). The
rest of the students' reasoning is broken down as follows: 11.1\% cited the
inability to take notes while wearing an HMD, and 5.6\% had never tried a VR
version at all (389).} The loaner headsets went unused because, as one student
described it, having to go to the library to check out the VR headset is ``a
hassle'' (388).
% [SEAM EDIT: "Despite making several Meta Quest HMDs available for checkout
%  through UCI's library, not a single student went to the library because..."
%  — the not-a-single-student fact now lives in the hook (389, verbatim);
%  this sentence trimmed to the 388 results-wording quote to avoid repetition.]

In our previous article, we stated our expectations at the start. We did not
expect virtual clinical immersion to rival physical presence in the operating
room: ``[i]t is unlikely that VR clinical immersion programs will be as
pedagogically effective as in-person observation'' (``Phenomenological
Evaluation'' 383). We hypothesized instead that the first-person perspective
enabled by VR, ``and the enhanced sense of presence and agency it generates, has
the potential to be more beneficial to students' development of such [clinical
needs-finding] skills than 2D videos as it more closely simulates the experience
of being physically present in the operating room'' (383). The results
substantiated half of that initial hypothesis---presence was at its highest in
the Meta Quest VR environment, but student preference did not follow suit---and
on the perceived pedagogical benefit, we were unable to render a verdict.
Because the course's commitment to open access let every student default to the
most convenient format,\footnote{Other than the occasional in-class use of the
VR headsets, students were not required to use any particular format and were
allowed to use whichever of the four versions of the video content they wanted:
YouTube, the PC version, Google Cardboard, or Meta Quest 2.} the four versions
could not compete on equal grounds: while a majority of students undeniably felt
more present in virtual reality, we could not ``conclusively claim that
interactive VR learning environments are more or less pedagogically efficacious
compared to 2D videos'' (390).\footnote{Our inability to render a verdict places
the pilot within a larger evidential pattern. Comparative studies that hold
content constant repeatedly find that added immersion raises felt presence
without raising learning: Makransky, Terkildsen, and Mayer, porting an identical
science simulation from desktop to headset, measured more presence and less
tested learning; Parong and Mayer's content-matched comparison found a slideshow
outperforming immersive VR on factual learning even as VR won every affective
measure; and Mayer, Makransky, and Parong's review of two decades of media
comparisons finds no reliable learning advantage for immersive VR as a medium.}
The pilot thus established two aspects and clearly drew a separation between
them: presence was achieved to a much higher extent in VR, but was not enough to
override the students' preference for convenience when having to view the
educational content at home on their own time. Whether the interactive virtual
world was more pedagogically effective---whether felt presence converts into
pedagogical worth---was left standing as a question our pilot's own conditions
could not answer.

% === (II) CONTEXT — part 2: the frame (NEW passage, approved 2026-07-20 with
%     "rhetorical incorporation" terminology fix)

Three concepts, built across the preceding parts, do the analytical work here.
The first is the actualization chain: engagement understood as the actualization
of desire, running from perception (A\textsubscript{1}) through the charged
image (\textit{phantasma}, A\textsubscript{2}) and committed judgment
(\textit{doxa}, A\textsubscript{3}) into completed action
(A\textsubscript{4})---a sequence a designed world can solicit, sustain, or
fall silent at any station. The chain also has a zero point
(A\textsubscript{0}): before any of its stations can fire, the student and the
world must reach one another at all---the threshold condition every later
moment presupposes. The second is rhetorical incorporation: the union a virtual world
achieves with its inhabitant, one in substance but dual in
being---world-disclosedness, the world opened to a student as a place they are
in, and co-disclosedness, the same world opened to students as a place they are
in together. The third is veri(dis)similitude: the requirement that a virtual
world be similar enough to what it renders to be entered as real, yet different
enough from its source material to justify its existence. A virtual learning
environment, in these terms, is not a delivery medium but a manufactured and
curated rhetorical ecology---an environment authored, in every element, toward a
single pedagogical end, whose design succeeds when both registers of disclosure
open and rhetorical incorporation is achieved.

% === (II) CONTEXT — part 3: the institutional resolution (M0.2, locked text)

After the pilot, we moved the course to a 2D-PC primary
deployment; while students cannot experience stereoscopic 3D video without a VR
headset, the videos in this interactive space are still 180- and 360-degree
presentations. Our primary reason for moving away from the 3D stereoscopic VR
development was due to logistics, feasibility, and necessary scale. This class
needed to be made available to undergraduate students across the entire
University of California (UC) system, with class sizes up to 120 students each
offering. Acquiring and distributing 120 Meta Quest headsets to each student
every quarter was simply not possible. Secondly, we chose to prioritize the
Windows PC-based version first as it was the largest student platform population
(compared to Apple laptops).\footnote{While we prioritized development for the
PC version, we are currently in the process of developing a version for Apple
computers, and development for VR headsets will commence once the Apple version
is operational.}

As just noted, this class---BME 179: Biomedical Engineering Design:
Addressing Unmet Clinical Needs, a one-unit online course---is offered
systemwide to undergraduate biomedical engineering (BME) students across
the UCs.\footnote{This class is
available to undergraduate BME students at every UC except the University of
California, San Francisco (UCSF), as UCSF is exclusively graduate-level health
sciences.} Students arrive with different course schedules, and the class is open
across two academic calendars---Berkeley and Merced run on semesters while
the rest of the universities are on quarters, though no students from the
two semester campuses have enrolled in the terms taught thus far---so a
fully asynchronous online format is not a convenience of choice, but the
condition of the class's systemwide existence, as a synchronous version
would be nearly impossible to logistically organize. An
in-person version was never possible at all, and cohort-wide scheduled
simultaneity was never in the design space: the deployment maximized the only
availability there was.
% AUTHOR FOLLOW-UP RESOLVED 2026-07-31 (batch-5 rulings): (1) class =
% BME 179: Biomedical Engineering Design: Addressing Unmet Clinical Needs
% (course syllabus, same across offerings; sources/Clinical Immersion
% Course Syllabus_Spring 2026.docx); (2) Berkeley/Merced: open to those
% campuses, but no semester-campus students have enrolled thus far —
% prose adjusted accordingly.

These constraints are more than administrative history; they define what kind of
thing the virtual learning environment is. The virtual world was built, in the
first place, because the actual clinical world admits so few---the operating
room's structural scarcity was the problem, and universal access was the
solution the virtual environment promised. A virtual world that students cannot
enter would therefore fail in a specific and instructive way: it would
reproduce, in software, the very exclusion it was built to overcome.
Availability is thus not a logistical afterthought but the primary pedagogical
property of this course and VLE: before a virtual world can show, teach, or
immerse, it must be enterable, and whatever a student cannot enter can disclose
nothing to them. In the vocabulary of the actualization chain, this is a
condition at A\textsubscript{0}, the chain's threshold: perception, the charged
image, committed judgment, and the act are all stations a student can only
reach inside a world they occupy, so a world that cannot be entered does not
engage its students poorly---it never begins to engage them at all.
Availability is the threshold's deployment face: the form A\textsubscript{0}
takes when the world must be downloaded, installed, and run before it exists
for its student. Nor is the scaling wall we hit peculiar to this course.
Immersive VR in higher education remains largely at the prototype
stage---experimental and development work rather than a regular part of teaching
(Radianti et al.\ 1, 22)---a pattern that will be discussed in greater detail
later. One further consequence of the deployment should be noted now and
developed later: a fully asynchronous format changes the time of the course.
There is no shared class hour; each student meets the virtual world on their own
schedule and at their own pace. The VLE itself is built for company---it is
multiplayer, with spatial proximity voice chat---and we clearly inform students
of this capability and encourage them to meet there in their assigned groups.
But with no scheduled hour to gather them, convening is left to each group's own
initiative, and the default is that a student enters the virtual hospital alone.
Whether students take up the invitation to meet---and what solitary viewing does
to their experience of the virtual world when they do not---are questions the
three terms of survey responses will answer.

% === (III) DISRUPTION (NEW; approved 2026-07-20 with largest-share
%     parenthetical added)

Set side by side, the pilot and the field leave a precise question standing.
70\% of our pilot cohort felt more present in the headset, yet 40\% judged
plain 2D video the most educationally beneficial format (the
largest share of the cohort to settle on any single format), and the loaner
headsets sat untouched in the library; the comparison literature, holding
content constant across media, finds the same shape at scale---added immersion
reliably raises presence and just as reliably fails to raise learning, while
presence itself proves attainable on an ordinary monitor. Presence, in short, is
neither scarce nor sufficient. What, then, must a virtual learning environment
add to presence before the world it renders earns its place in a curriculum?
What would a student gain inside the virtual world that the same footage, one
click away on YouTube, cannot give them---and what kind of analysis could even
name that difference?

% === (IV) RESOLUTION (NEW; approved 2026-07-20 with Second Life rephrase)
% REVISION-PASS MARKER: the social-interior passage (from "There is a material
% reason" through "authored for a curriculum") is to be revised against the
% new-media/presence sources once ingested + indexed (Boellstorff, Davis &
% Boellstorff, Champion, Spagnolli, Mantovani & Riva, Riva, et al.).

The VR-pedagogy tradition has measured presence with increasing precision while
conceding it cannot say what presence is for; the social-presence tradition has
known since Short, Williams, and Christie's founding work in 1976 that the felt
company of others carries satisfaction and learning, but it matured in distance
education---discussion forums, video sessions, text---around media that were
channels rather than places. The two research programs descend from the same
1976 origin and cite one another almost never, having built their houses in
different fields around different objects. There is a material reason the
divide has persisted: until recently, few deployed pedagogical virtual
environments had a social interior. Multiplayer capability and spatial
proximity voice chat---the features that make a virtual world a place where
students can be together---have been rarely implemented. Where educators
have taught inside socially capable virtual worlds in relatively recent
times---the Second Life classrooms documented by anthropologists and education
researchers offer one example---the world has often been a general-purpose
social platform adapted to teaching rather than an environment authored for a
curriculum.\footnote{On the education-research side, the fullest example in
this section's literature review is Merchant et al.'s structural-equation study
of college chemistry instruction inside Second Life (N\,=\,204), taken up in
detail below. On the anthropological side, Tom Boellstorff---an anthropologist
at UC Irvine---conducted multi-year ethnographic fieldwork inside Second Life
through the avatar Tom ``Bukowski,'' the fieldwork behind his monograph
\textit{Coming of Age in Second Life}; he recounts that work in ``Virtual
Worlds and Futures of Anthropology'' and carries the program forward to the
present platform generation in ``Toward Anthropologies of the Metaverse.''} A
desktop VLE of the present kind therefore sits at a junction
neither literature has occupied: an explorable virtual hospital authored for one
clinical curriculum, deployed for credit across a university system, its
students group-assigned and equipped with in-world proximity voice chat---at
once a
rendered world of the kind the VR-pedagogy literature measures and a social
space of the kind the social-presence literature theorizes. The VR-pedagogy
studies compare media and score presence but do not ask what being-together
inside the world does; the social-presence studies measure felt community in
forums and video sessions but not inside a navigable world; to our knowledge, no
analysis has read a single deployment with both sets of questions in hand.
Reading one requires a frame built for the junction, and that is what this
dissertation's rhetorical-ontological and phenomenological frame supplies. Where
the field's instruments report that presence rose or interest fell, they cannot
say how a wholly authored environment solicits a student's desire, carries it
toward action, or falls silent before it; the actualization chain names those
moments.
Where presence names a feeling, world-disclosedness and co-disclosedness name
the two registers in which a virtual world is opened to its inhabitants---as a
place I am in, and as a place we are in together---so that what a design
achieves or forfeits can be located, not merely scored. A VLE, on this reading,
is a manufactured and curated rhetorical ecology with a single pedagogical end,
and its effectiveness is a design question in the strictest sense: the
environment teaches when it is designed so that both registers open---when
rhetorical incorporation is achieved.

% === (V) SIGNPOST (M0.4, GOLD STANDARD — untouched; see FCDP §10.2)

The survey given to all students across the Winter 2025, Spring 2025, and
Winter 2026 quarters was designed, fielded, and collected by myself and
Prof.\ Christine King---the instructor of record---and is read here
through the rhetorical-ontological and phenomenological frame this dissertation
has been building. The data yield the findings, with every count and quoted
response coming from the survey corpus. The theoretical frame supplies the how
and the why: the actualization chain and the methodology of
Rhetorical-Ontological Diachronic Analysis (RODA) serving as a background
grammar for interpreting students' responses---which moment of engagement is in play when a
student reports presence, confusion, the flight to YouTube, or the desire for
more socially interactive course design. Accordingly, what follows moves through
seven parts. The first reviews the relevant literature in its two traditions:
the desktop route as a pedagogical modality in its own right, with presence
relocated from hardware to the design of the environment; the comparative
evidence that more presence does not deliver more learning; the rationale for
asynchronous deployment at scale; and the research on learning together in
shared virtual space---social presence and computer-supported collaborative
learning---whose near-total lack of contact with the VR-pedagogy literature is
itself a finding this section works across. The second part describes the VLE
and its deployment: the virtual hospital as a designed world, the cohort of 241
students across the three quarters, and the survey instrument together with the
standing limits of what it can show. The third presents the survey findings in
order: the threshold problems that kept some students out of the virtual world;
what students found, and could not do, once inside; the flight to YouTube; the
presence students attested on ordinary screens; the group co-watching capability
that was had, wanted, and largely unactivated; and the shift in student
skepticism from technical friction to the question of the platform's value. The
fourth part conducts the rhetorical-ontological and phenomenological analysis:
it locates the primary rupture point in the student experience, reads the
boredom the data attest through the dissertation's three-form account, takes
up the justification question the data raise, and shows how the VLE is
world-disclosed and co-disclosed to its students. The fifth turns diagnosis into design, ranking
the design levers by the moment of engagement each one works on, with the
activation of the WE given the fullest treatment. The sixth states the limits of
what this survey data can unconceal. The seventh names what the analysis hands
forward---to the boredom analysis, a separate section that follows this one, and
to the field.
% NOTE (user 2026-07-20): the boredom analysis is its OWN SECTION, not a
% subsection of this one. RE-VERIFIED 2026-07-31 (close-out pass): seven-part
% count matches M1–M7 as revised through the full readthrough; fourth-part
% description updated to name the three-form boredom reading.

% === (VI) THESIS (M0.3 relocated to final position; opener retrofit marked)

By reading the three quarters of survey data through this frame, the following
analysis will substantiate a claim with three connected parts.
% [SEAM EDIT: was "The claim this section will substantiate can be stated at
%  the outset, in three connected parts." — "at the outset" no longer true in
%  final position; retrofit per approved plan.]
First: the deployed desktop world succeeds where our pilot said screens could
succeed. Across three terms, students attest to presence on ordinary monitors at
rates matching what the pilot measured in VR headsets; a virtual world on a flat
screen can still place its student somewhere, and world-disclosedness does not
require a headset. But presence, I will show, is not the finish line. As
deployed, the virtual hospital gives students a great deal to see---filmed
operations, physician and engineer interviews, explorable procedure
rooms---and watching is nearly all it gives them to do. Watching is itself an
act: a student can stand in the virtual operating room and refuse to attend,
and the one who watches has actualized a potentiality the world furnished. The
shortfall lies in what else the world makes available: no task to carry out
inside it, no procedure to attempt, no problem whose solution it asks of the
student---no richer or more pedagogically effective potentiality for the
student's engagement to actualize. In the chain's terms, the deployment
carries engagement to committed judgment (A\textsubscript{3})---the world wins
the student's attention and their assent that it is worth taking
seriously---but the only completed act (A\textsubscript{4}) it affords that
commitment is one the student could have performed anywhere: watching. The
world falls silent at the passage from A\textsubscript{3} to
A\textsubscript{4} not because no act follows, but because the acts it affords
stop at watching. Second:
because every video housed in the virtual hospital also exists one click away on
YouTube, the desktop world carries a standing burden of justification. It must
be different enough from the video of itself to justify existing, and that
difference can only be made of what video cannot render: a place to be in, a
task to do, and others to be there with. Third: of those three materials, the
last is the one this deployment alone possesses. Watching together in the same
virtual space, within earshot of one another, is a capability no other route in
this program's history has offered---the published VR version had no real-time
collaborative features at all, and students could only enter it alone
(``Phenomenological Evaluation'' 389). Where place and task deepen
world-disclosedness, being there together opens a different register
altogether---co-disclosedness: the virtual world disclosed not to each student
alone but to a WE. I will argue that this capability is at once the desktop
world's deepest answer to ``why not just YouTube?'' and, as the data
demonstrates, its deepest unrealized asset. The scope of this claim is a matter
of development history rather than principle: nothing in a headset forbids
company; it is on the desktop route that company was built. What the analysis
owes, then, is evidence for each part: that the world was disclosed, that
the acts it furnished stopped at watching, and that co-disclosedness---the
one capability that could justify the world---went largely unused.
% [SEAM EDIT: "What the section owes" -> "What the analysis owes" (self-
%  reference trim). THESIS REVISION EXECUTED 2026-07-31 (user-ordered
%  close-out, NS#3): thesis verified aligned with the completed section —
%  part 1 = M3.4/M4.1 (70.7% presence; acts stop at watching, doctrine
%  form); part 2 = M4.2 (burden of justification; place/task/others);
%  part 3 = M4.3/M7.2 (co-disclosedness = deepest answer + unrealized
%  asset; one confirmed convening). M7.2 close rewritten same day to
%  demonstrate the claim in matching terms ("The vocabulary held").]

% ============================================================
% v2 LOCKED 2026-07-20 (hook follow-up approved; seams 1-3 approved; all six
% moves user-ruled). OUTSTANDING REVISION PASSES: (a) social-interior passage
% in the resolution — revise against new-media/presence sources once indexed
% (ingestion Phase 1 done, snapshot doc-text-20260720T145233Z; entry phases
% user-gated); (b) roadmap part-count + thesis — at first complete section
% draft; (c) RESOLVED 2026-07-31: BME 179 named in prose; Berkeley/Merced
% no-enrollee fact rendered.
% ============================================================

% ========== END INLINED FILE: DESKTOP-M0-DRAFT-v2.tex ==========

% ===== M1 (carries \section{The Two Literatures})
% ========== BEGIN INLINED FILE: DESKTOP-M1-DRAFT-v1.tex ==========
% ============================================================
% Part III — Desktop Section — M1 (The Two Literatures) — DRAFT v1 (rev 2, 2026-07-20)
% Mode: Stage 3 guided drafting, parallel session (module-by-module cadence)
% Blueprint: PART-III-DESKTOP-SECTION-OUTLINE-v1-2026-07-15.md, M1 as revised
%   2026-07-19 (M1.2+M1.3 merged; M1.5 deferred to the boredom-experiment
%   section; four movements total). Sectioning level provisional pending
%   Open decision #1 (arc placement).
% Conventions: MLA parenthetical cites (abbreviated titles where an author
%   set has multiple works; bare page for same-work repeat within a
%   paragraph); footnotes substantive only; first person for King & Salvo
%   (full MLA cite + convention footnote lives in the M0 draft);
%   NO abstract quotations; no sentence-initial "And"; "record" banned as a
%   name for the evidence (use "the survey data"); no metatextual signposting.
% REV 2 (2026-07-20, per user notes + cross-session verification): fidelity
%   DE-EMPHASIZED throughout M1.1 — the psychological-determination verdict
%   (DA-03), M0.3's locked claim (world-disclosedness does not require a
%   headset), and Dubovi's own Table 3 (only Involved/Control correlates with
%   learning) all demote representational fidelity to one tool among several;
%   interaction/task now leads wherever the pair is named. Dubovi quote kept
%   as expansion-and-disagreement per user; vague "half of that sentence"
%   framing replaced with plain statement of the two clauses.
% Quotes: PDF-verified per sources/vle-pedagogy-lit-quote-bank.md (ROUND 2
%   dispositions folded in). Pages empirical from headers/footers.
%   NB Dalgarno & Lee 2012: no printed footers in PDF; pages inferred from
%   the proceedings range (236-245) — bank §0.
% Batch 1 (2026-07-19, rev 2 2026-07-20): M1.1. Batch 2 (2026-07-20): merged
%   M1.2/3 below.
%
% WORKS CITED entries introduced by this module (for later assembly):
%   Dalgarno, Barney, and Mark J. W. Lee. "What Are the Learning Affordances
%     of 3-D Virtual Environments?" British Journal of Educational
%     Technology, vol. 41, no. 1, 2010, pp. 10-32.
%   ---. "Exploring the Relationship between Afforded Learning Tasks and
%     Learning Benefits in 3D Virtual Learning Environments." Future
%     Challenges, Sustainable Futures: Proceedings ascilite Wellington 2012,
%     ascilite, 2012, pp. 236-45.
%   Dubovi, Ilana, et al. "Now I Know How! The Learning Process of Medication
%     Administration among Nursing Students with Non-Immersive Desktop
%     Virtual Reality Simulation." Computers & Education, vol. 113, 2017,
%     pp. 16-27.
%   Fowler, Chris. "Virtual Reality and Learning: Where Is the Pedagogy?"
%     British Journal of Educational Technology, vol. 46, no. 2, 2015,
%     pp. 412-22.
%   Makransky, Guido, and Gustav Bøg Petersen. "Investigating the Process of
%     Learning with Desktop Virtual Reality: A Structural Equation Modeling
%     Approach." Computers & Education, vol. 134, 2019, pp. 15-30.
%   Merchant, Zahira, et al. "The Learner Characteristics, Features of
%     Desktop 3D Virtual Reality Environments, and College Chemistry
%     Instruction: A Structural Equation Modeling Analysis." Computers &
%     Education, vol. 59, 2012, pp. 551-68.
%   Mikropoulos, T. A., and A. Natsis. "Educational Virtual Environments: A
%     Ten-Year Review of Empirical Research (1999-2009)." Computers &
%     Education, vol. 56, 2011, pp. 769-80.  [initials kept until verified]
% ============================================================

\section{The Two Literatures}

% --- M1.1 Desktop 3D as a modality in its own right (rev 2, 2026-07-20)

The desktop route has a research literature of its own, and that literature's
founding model already refuses the reading a headset-first field invites: the
desktop build as a lesser headset. Dalgarno and Lee's model of learning in
three-dimensional virtual environments locates what is pedagogically
distinctive about the medium in what an environment lets its learner do and
see---learner interaction first, representational fidelity alongside it---and
treats presence and co-presence not as endowments of hardware but as
consequences of those characteristics: in the model's words, ``a consequence
of the representational fidelity and learner interactions facilitated by the
environment'' (``Learning Affordances'' 14). The affordances the model
derives are, accordingly, task-possibilities in Gibson's sense---things an
environment makes available to an agent to do---not display achievements.
The model arrived carrying its authors' own verdict on the field's evidence:
``[m]uch of what has been published about the educational uses of 3-D
technologies is largely `show-and-tell', presenting only anecdotal evidence
or personal impressions that cannot be usefully generalised beyond the local
context'' (23). When the same authors put the model to its own test two
years later---surveying fifty-three higher-education staff on which afforded
learning tasks had yielded which of the five predicted benefits---respondents
endorsed every benefit at nearly uniform strength regardless of task, none of
the hypothesized task-to-benefit links reached significance, and the authors
flagged the deflating explanation themselves: the uniform agreement ``could
be more a reflection of their enthusiasm and positive attitude towards the
use of 3D VLEs than an evidence-based assessment of the actual learning
benefits'' (``Afforded Tasks'' 243). The field's central model thus stands
named beside its own inconclusive self-test.

The model's sharpest critique asked not for more data but for pedagogy: a
decade of educational-VR research had run technology-first, deriving expected
benefits from technical affordances with almost no explicit account of how
learning occurs.\footnote{The ten-year review behind the charge found
explicit learning theory almost wholly absent from the literature it
surveyed: ``[a]ll the other reviewed articles do not refer explicitly to a
learning theory'' (Mikropoulos and Natsis 775, qtd.\ in Fowler 412). Fowler's
repair works outward from intended learning outcomes to the technology that
would serve them, rather than inward from the technology's affordances
(412--13, 416--20).} What Fowler's repair adds, and what this section will
hold onto, is a curve: ``the assumption that high levels of representational
fidelity and learner interaction will result in deeper learning is
questionable. It may be argued that there are optimum levels of these
characteristics that maximise learning. Going beyond the optimum increases
costs and results in a limited or even negative return with respect to
learning benefits'' (Fowler 415). Neither characteristic improves learning in
proportion to its own increase: each has an optimum level, and beyond it,
added fidelity or added interaction yields limited or even negative returns
for learning. That caution matters most for representational fidelity,
because fidelity is the characteristic that headset hardware exists to
advance---and the one a headset-first field kept mistaking for the medium's
pedagogical substance. Our own pilot had already produced a concrete instance
of that curve. We expected that students' familiarity with VR and video games would
incline them toward the virtual-reality versions; the correlation ran the
other way: of the students who judged the plain 2D videos most
educationally beneficial, 66.7\% were weekly gamers, against 33.3\% of those
who favored the VR versions---and our interviews traced the judgment to
builds that ``do not meet the production standards of contemporary VR and
traditional video games'' (``Phenomenological Evaluation'' 389). We stated
the result then at its proper strength---a possible correlation, awaiting
validation (389)---but its direction is the one the curve predicts once
experience is allowed to move the optimum: the students most practiced in
polished, feature-rich worlds hold the highest bar, and a build that
undershoots it loses them first.

Within the desktop medium itself, the process studies converge on where the
leverage lies, and it is not the display. Makransky and Petersen, modeling
199 medical undergraduates through a desktop genetics simulation, found two
routes from the medium's features to knowledge gains---``the affective and
the cognitive paths'' (``Investigating the Process'' 24): presence feeding
motivation on one, perceived usability feeding cognitive benefit on the
other, the two converging on self-efficacy as the sole direct predictor of
learning.\footnote{Self-efficacy $\rightarrow$ learning, $\beta = 0.579$;
presence's effect reaches learning only through motivation and then
self-efficacy, never directly (``Investigating the Process'' 23--24).} In
% FLAG 2026-07-31 (from M0 Second Life footnote verification): index unit
% ped-sec-16 records the SL treatment as three graded activities over roughly
% eight weeks, not a full "semester of Second Life chemistry" — user to rule
% on rewording at the subsection-2 readthrough.
Merchant and colleagues' semester of Second Life chemistry (N\,=\,204), the
gate is usability: ``the 3D virtual reality features can support the
development of learners' spatial orientation ability, self-efficacy, and
presence only when the learners' [sic] perceive the experience as meaningful
and the system easy to use'' (Merchant et al.\ 562).\footnote{The mediation
is nearly total ($\beta = 0.956$, VR features $\rightarrow$ usability;
Merchant et al.\ 560). The design is single-condition desktop---no headset
arm---so the study licenses no claim about media sufficiency in either
direction.} Dubovi, Levy, and Dagan's nursing simulation, finally, returned
the literature's flattest correction to the hardware account of presence: on
ordinary monitors, their students' presence scores met or exceeded the norms
the construct's founding instrument had recorded for immersive
VR,\footnote{Involved/Control 5.2\,$\pm$\,0.7, Natural 4.7\,$\pm$\,1.03,
Interface Quality 5.0\,$\pm$\,1.1, against Witmer and Singer's immersive-VR
norms of 5.2\,$\pm$\,0.8, 4.1\,$\pm$\,1.1, and 4.8\,$\pm$\,1.1 respectively
(Dubovi et al.\ 23).} and they drew a two-part conclusion whose parts deserve
different fates: ``presence is not a unique attribute of the environment, but
it is induced by the fidelity of representation and the high degree of
interaction and user control'' (Dubovi et al.\ 25).

The first clause is right, and the survey data I will analyze later on
confirm it at scale:
presence is not what headsets have and screens lack. The second clause gives
fidelity a role the evidence does not sustain---not the wider field's
evidence, and not this study's own. The one component of their presence
measure that correlated with learning was involvement and control
($r = 0.211$ overall, $0.231$ on the procedure's core subscale); the
realism-adjacent components tracked nothing (25).\footnote{Their own design
principle, stated two sentences before the fidelity clause, names control
alone: ``ensuring that participants gain sufficient control of the
environment is an important condition for learning'' (Dubovi et al.\ 25).
The fidelity clause is carried by citation to prior literature, not by the
study's own data.} Presence flourishes in worlds whose fidelity is schematic
or frankly poor---the cube-built landscapes of \textit{Minecraft}, the aging
textures of \textit{World of Warcraft}---worlds no fidelity-led account
explains. What induces presence is not the exactness of the copy but the
constitution of the artifact as a world: an environment built as an ecology,
one that answers action---a world the student moves and is moved by in turn.
Fidelity is one tool for building such a world. It is not the ground of it.

% --- merged M1.2/3 Presence on screens, and its limit (batch 2, 2026-07-20)
% Quotes PDF-verified per bank §2 (Lombard & Ditton by direct PDF read —
%   OCR-LOW protocol; JCMC unpaginated -> section-name cites). Mayer et al.
%   = 2022 version-of-record (print issue 2023). Parong & Mayer 2018 pages
%   = APA Online First; pin at assembly. Additional Works Cited entries:
%   Barrett, Robin C. A., et al. "Comparing Virtual Reality, Desktop-Based
%     3D, and 2D Versions of a Category Learning Experiment." PLOS ONE,
%     vol. 17, no. 10, 2022, e0275119.
%   Calleja, Gordon. In-Game: From Immersion to Incorporation. MIT Press,
%     2011.
%   Chow, Meyrick. "Determinants of Presence in 3D Virtual Worlds: A
%     Structural Equation Modelling Analysis." [venue pinned at assembly],
%     2016.
%   Cook, David A., et al. "Technology-Enhanced Simulation for Health
%     Professions Education: A Systematic Review and Meta-Analysis." JAMA,
%     vol. 306, no. 9, 2011, pp. 978-88.
%   Kononowicz, Andrzej A., et al. "Virtual Patient Simulations in Health
%     Professions Education: Systematic Review and Meta-Analysis by the
%     Digital Health Education Collaboration." Journal of Medical Internet
%     Research, vol. 21, no. 7, 2019, e14676.
%   Lombard, Matthew, and Theresa Ditton. "At the Heart of It All: The
%     Concept of Presence." Journal of Computer-Mediated Communication,
%     vol. 3, no. 2, 1997.
%   Makransky, Guido, Thomas S. Terkildsen, and Richard E. Mayer. "Adding
%     Immersive Virtual Reality to a Science Lab Simulation Causes More
%     Presence but Less Learning." Learning and Instruction, vol. 60, 2019,
%     pp. 225-36.
%   Mayer, Richard E., Guido Makransky, and Jocelyn Parong. "The Promise
%     and Pitfalls of Learning in Immersive Virtual Reality." International
%     Journal of Human-Computer Interaction, vol. 39, no. 11, 2022,
%     pp. 2229-38.
%   Parong, Jocelyn, and Richard E. Mayer. "Learning Science in Immersive
%     Virtual Reality." Journal of Educational Psychology, vol. 110, no. 6,
%     2018, pp. 785-97.
%   ---. "Cognitive and Affective Processes for Learning Science in
%     Immersive Virtual Reality." Journal of Computer Assisted Learning,
%     vol. 37, 2021 [pages pinned at assembly].
%   Witmer, Bob G., and Michael J. Singer. "Measuring Presence in Virtual
%     Environments: A Presence Questionnaire." Presence: Teleoperators and
%     Virtual Environments, vol. 7, no. 3, 1998, pp. 225-40.

Presence itself entered the research literature as a property of persons,
not of headsets. Lombard and Ditton's founding synthesis defines it as
``the perceptual illusion of nonmediation''---an illusion that ``occurs when
a person fails to perceive or acknowledge the existence of a medium in
his/her communication environment and responds as he/she would if the medium
were not there''---and the definition is deliberately medium-general, built
to hold across television, film, and virtual reality alike (``Presence
Explicated''). The consequence is drawn in the same section: ``[b]ecause it
is a perceptual illusion, presence is a property of a person,'' varying with
a medium's form, its content, and the person attending to it (``Presence
Explicated''). The founding instrument agrees about the location. Witmer and
Singer, constructing the field's first validated presence questionnaire,
reject the view that immersion is an objective property of the display
hardware: ``immersion, like involvement and presence, is something the
individual experiences'' (227). Their paper is equally plain about what had
not been shown: factors that increase presence were expected to increase
learning, ``but as yet we have no direct evidence to support our
contention'' (239). Two decades of determinants research kept resolving the
question user-side---in Chow's structural model of a desktop virtual world,
computer self-efficacy, the student's confidence with the machine, was the
strongest predictor of presence, ahead of any property of the display
(11--13).\footnote{Chow's model is scoped to self-presence rather than
spatial presence, and his own design implication follows the data: invest in
orientation and user confidence rather than in richer rendering alone
(12--13).} Across research traditions that rarely cite one another, the
convergence is the same: presence is achieved by persons and occasioned by
worlds; no hardware tier owns it.

The media-comparison studies of the past two decades tested the direct
question: does the added presence of a more immersive medium produce more
learning? Across content-matched designs, the answer was no.
Makransky, Terkildsen, and Mayer ported an
identical laboratory simulation from desktop monitor to head-mounted
display: presence rose sharply in the headset (d\,=\,1.30 in the first
lesson, 2.20 in the second), tested learning fell, and EEG measures located
the mechanism in cognitive overload (``Adding Immersive'' 232--33).
\footnote{The dissociation also holds inside a single desktop condition: in
Merchant and colleagues' chemistry model, usability generated presence
($\beta = 0.540$) while presence's path to learning was null
($\beta = 0.069$, $p = .367$; Merchant et al.\ 560).} Parong and Mayer ran
the media comparison with content matched line for line: a narrated
slideshow beat immersive VR on the knowledge posttest in both experiments
while VR won every affective measure taken---enjoyment, motivation,
engagement, interest (``Learning Science'').\footnote{The slideshow's advantage on the knowledge
test given after the lesson was d\,=\,0.92 in the first experiment and 1.12
in the second. The same
study carries the constructive result: a generative summarizing prompt
inside VR closed the learning gap without reducing enjoyment---the deficit
is remediable by instructional design, not intrinsic to the medium.} Their
mediation follow-up closed the causal loop: the medium's effect on retention
ran through self-reported extraneous load and high-arousal positive
affect---the features that make immersive VR appealing are the ones
implicated in its learning cost (``Cognitive and Affective
Processes'' 236--37).
% [Short title normalized at assembly 2026-07-29 to match M7.1's
%  "Cognitive and Affective Processes" (MLA first-noun-phrase form);
%  per M7 header's "normalize at assembly" instruction.]
Mayer, Makransky, and Parong's review of the program's two decades states
the upshot plainly: ``[t]here is not strong evidence that students learn
academic material more deeply when it is presented in IVR than when it is
presented with conventional media'' (2236). Barrett and colleagues sharpen
the point from the psychophysics side: on one category-learning task,
learning accuracy was statistically equivalent across VR, desktop-3D, and 2D
versions; what differed was response time, fixation behavior, and---in VR
alone---attrition from physical discomfort (12--14).\footnote{The
equivalence claim is accuracy-specific, and the 2D arm is small
(n\,=\,8)---a scope limit worth carrying with the citation.}

The health-professions simulation meta-analyses locate where screen-based
practice does reliably improve learning: in skills and clinical reasoning
rather than declarative knowledge.
Kononowicz and colleagues' meta-analysis of screen-based virtual patients
found nothing added on knowledge over traditional formats (SMD\,=\,0.11,
non-significant) and a moderate-to-large advantage on skills and clinical
reasoning (SMD\,=\,0.90) (6--7). The base rate behind that split is settled:
across 609 studies of technology-enhanced simulation, pooled effects against
no intervention are large on every outcome family, and the reviewers close
the question themselves---``we question the need for further studies
comparing simulation with no intervention'' (Cook et al.\
987).\footnote{The two meta-analyses do not overlap: Cook and colleagues
pool technology-enhanced simulation broadly and explicitly exclude the
screen-based virtual patients Kononowicz and colleagues analyze. The
virtual-patient meta-analysis of that earlier era is Cook, Erwin, and
Triola (2010), a different study, not drawn on here.} Simulation teaches;
the evidence thins exactly where a world is only watched. The knowledge/
skills split names where the medium's value lives: not in what a student
sees from inside it, but in what the student is given to do there.

Read together, these two decades leave a precise explanatory hole. Presence
belongs to the person; presence does not convert into learning on its own;
learning gains concentrate where the environment gives its student something
to do. The field has measured all three results and has no construct that
joins them---no name for the difference between being placed in a world and
being taken into it as an agent whose acts complete what the world offers.
Game studies built that construct: in Gordon Calleja's \textit{In-Game: From
Immersion to Incorporation}, incorporation names the double movement in which
the virtual environment is absorbed into the player's consciousness while the
player is embodied within it as an agent. The
education literature does not import it; the bridge between the two fields
remains unbuilt in both directions. Part II rebuilt the construct on this
dissertation's own ground---rhetorical incorporation, one in substance and
dual in being, world-disclosedness and co-disclosedness---and the desktop
deployment now gives it empirical work to do. The survey data show a world
that places its students: presence attested on ordinary monitors, exactly
where the field's own findings say it can live. What the data must then
answer is whether the world completes what it offers---and that is a
question about rhetorical incorporation, not about presence, which is why the frame
Parts I--II built is the instrument this section reads with.

% --- M1.4 The asynchronous/scale rationale (batch 3, 2026-07-20)
% Quotes PDF-verified per bank §3. Additional Works Cited entries:
%   Cossio, [first names pinned at assembly], et al. "Cybersickness and
%     Discomfort from Head-Mounted Displays Delivering Fully Immersive
%     Virtual Reality: A Systematic Review." Nurse Education in Practice,
%     vol. 85, 2025, 104376.
%   Makransky, Guido, et al. "Equivalence of Using a Desktop Virtual
%     Reality Science Simulation at Home and in Class." PLOS ONE, vol. 14,
%     no. 4, 2019, e0214944.
%   Radianti, Jaziar, et al. "A Systematic Review of Immersive Virtual
%     Reality Applications for Higher Education." Computers & Education,
%     vol. 147, 2020, 103778. [full subtitle pinned at assembly]

Whether an at-home, unsupervised deployment gives up anything a supervised
classroom would provide has been tested directly. Makransky, Mayer, Veitch,
and colleagues had 112 students complete the same desktop science simulation
either at home or in class: learning outcomes did not differ (d\,=\,0.14),
nor did intrinsic motivation (d\,=\,0.02) or self-efficacy (d\,=\,0.09),
and the null pattern survived adjustment for prior knowledge and goal
orientation (``Equivalence'' 8--9).\footnote{One interaction survived
adjustment: mastery orientation predicted self-efficacy more strongly in
the classroom group (``Equivalence'' 9--10). The sample is self-selected
(112 of 289 eligible students consented), a limit the authors flag
themselves (11--12).} Their interpretation extends Clark's old verdict on
media to physical context: the instructional method---the simulation
itself---drives the learning, and felt presence is proposed, though not
tested, as what renders the location unimportant (10). The scope of the
design deserves equal attention. Both groups worked with the identical
desktop simulation---a Labster-built virtual biology laboratory---and only
the setting changed: one group at home, unsupervised, the other in a
classroom with a teacher present. No group received a different medium.
The study therefore answers one question and leaves a second untouched.
It shows that moving a desktop simulation out of the classroom and into
the student's own home, on the student's own schedule, costs nothing
measurable in learning, motivation, or self-efficacy---the fact an
asynchronous online course needs established. It does not show that an
interactive 3D simulation teaches better than a flat video of the same
content, because no video condition existed in the design. That second
question is the one our own deployment poses: every term, part of the
class entered the interactive desktop world while another part watched
the same operating-room footage on YouTube instead, and the survey data
speak directly to what separated those two experiences.

The field's own deployment history tells the complementary story. Radianti
and colleagues, reviewing 38 immersive-VR applications in higher education,
found that ``the majority of articles indicated that VR for education is
still in its experimental stages---prototyping and testing with students,''
with evaluation skewed toward usability rather than learning and the
maturity level itself ``a barrier for its adoption in regular teaching
activities'' (22).\footnote{The review's inclusion criterion is
head-mounted immersive VR only---desktop VR is excluded by design---so the
prototype-stage verdict describes precisely the hardware route, not the
desktop one (24).} A virtual world in routine use every term, across a
university system, is what that literature documents almost never
happening; the desktop build is what makes the routine part attainable.
The headset route also carries a physical cost the desktop route does not
levy. Cossio and colleagues' review of 25 HMD studies reports a mixed
rather than uniform picture---twelve studies with inconsistent or very low
cybersickness scores, five with high scores, oculomotor disturbance the
most common symptom---but where head-mounted delivery was compared directly
against desktop or non-immersive formats, the worse outcomes fell on the
headset side, ``likely due to its more immersive nature and the potential
for greater sensory conflict'' (14).\footnote{The aggregate is a reportable
moderator, not a uniform verdict: Huber and colleagues (2018) found
cybersickness \textit{less} prevalent in fully immersive VR (3.3\%) than in
augmented reality (6.6\%) (Cossio et al.\ 14). The same cost surfaces as
the VR-only discomfort attrition in Barrett and colleagues' three-arm
comparison.} Three findings from this literature bear directly on the
shape of our course. Students learn as much from a desktop simulation
completed at home as from the same simulation completed in a supervised
classroom (``Equivalence'' 8--9). Headset VR in higher education remains,
by the field's own account, a prototype technology---built and tested with
students, only rarely made part of a regular course (Radianti et al.\ 22).
Headsets carry a documented cybersickness cost that ordinary monitors do
not (Cossio et al.\ 14). Our course is fully online and asynchronous,
serving undergraduate biomedical engineering students across the
University of California's campuses, each on their own machine and their
own schedule. For a course of that shape, the desktop deployment is not a
diminished substitute for headset delivery. The home-versus-classroom
evidence says our students lose nothing by working where and when they
want. The deployment reviews say the headset alternative has almost never
been made routine at any scale. The cybersickness literature says the
monitor spares them a physical cost besides. The desktop build is what
the field's own evidence recommends for the course we actually teach.

% --- M1.6 Learning together in shared virtual space (batch 4, 2026-07-20)
% Quotes PDF-verified per bank §4. Ledger: seam DOCUMENTED here, argued in
%   M0.5 (other session); Biocca/Terry & Doolittle at statement level (argued
%   use reserved for M4.3); Kreijns Pitfall 1 quoted, Pitfall 2 glossed (full
%   deployment reserved for M5.2). Additional Works Cited entries:
%   Biocca, Frank, Judee K. Burgoon, and [initials at assembly] Harms.
%     "Toward a More Robust Theory and Measure of Social Presence: Review
%     and Suggested Criteria." Presence: Teleoperators and Virtual
%     Environments, vol. 12, no. 5, 2003, pp. 456-80.
%   De Back, Tycho, et al. "Learning in Immersed Collaborative Virtual
%     Environments: Design and Implementation." [venue + pages pinned at
%     assembly], 2021.
%   Garrison, D. Randy, Terry Anderson, and Walter Archer. "Critical
%     Inquiry in a Text-Based Environment: Computer Conferencing in Higher
%     Education." The Internet and Higher Education, vol. 2, no. 2-3,
%     1999, pp. 87-105.
%   Gunawardena, Charlotte N., and Frank J. Zittle. "Social Presence as a
%     Predictor of Satisfaction within a Computer-Mediated Conferencing
%     Environment." American Journal of Distance Education, vol. 11,
%     no. 3, 1997, pp. 8-26.
%   Kreijns, Karel, Paul A. Kirschner, and Wim Jochems. "Identifying the
%     Pitfalls for Social Interaction in Computer-Supported Collaborative
%     Learning Environments: A Review of the Research." Computers in Human
%     Behavior, vol. 19, 2003, pp. 335-53.
%   Richardson, Jennifer C., and Karen Swan. "Examining Social Presence in
%     Online Courses in Relation to Students' Perceived Learning and
%     Satisfaction." [venue pinned at assembly], 2003.
%   Richardson, Jennifer C., et al. "Social Presence in Relation to
%     Students' Satisfaction and Learning in the Online Environment: A
%     Meta-Analysis." [venue pinned at assembly], 2017.
%   Short, John, Ederyn Williams, and Bruce Christie. The Social
%     Psychology of Telecommunications. Wiley, 1976.
%   Terry, Krista, and Peter Doolittle. "Re-examining Social Presence:
%     Implications for Digital Pedagogies." Technology, Instruction,
%     Cognition and Learning, vol. 11, 2019, pp. 121-39.
%   Tu, Chih-Hsiung. [title/venue pinned at assembly], 2002.
%   van der Meer, [first names + byline form pinned at assembly], et al.
%     [title pinned], Frontiers in Virtual Reality, vol. 4, 2023, 1159905.

The section's second literature studies a different presence: not the sense
of being in a place but the sense of being with someone. Social presence
entered the research vocabulary in 1976, in Short, Williams, and Christie's
social psychology of telecommunications, naming how fully a communication
medium lets its users register one another as persons---decades before
rendered worlds existed to be present in. Biocca, Burgoon, and Harms,
surveying the construct's first quarter-century, define its core as ``the
sense of being with another'' (456) and set out the criteria a robust
theory of it must meet. Their sharpest verdict falls on the field's own
instrument. The inherited scale was built to rate media channels, not to
describe interactions, and the authors press the question that follows:
``is social presence just an attitude towards a medium?''
(469).\footnote{The scale's provenance explains its orientation: it
descends from studies funded by the UK Post Office, the Department of
Transportation, and General Electric---organizations commissioning
ratings of their communication channels (Biocca et al.\ 469).} A
construct meant to describe an interaction had been operationalized as an
opinion about equipment.

What the construct predicts, in its home tradition of online education, is
substantial. In Gunawardena and Zittle's origin study of a text-only
academic conference, social presence alone accounted for roughly 60
percentage points of the variance in learner satisfaction
(18--19).\footnote{Their closing design recommendation cites Johansen,
Vallee, and Spangler's suggestion that social presence can ``be
cultured''---built up by design and practice rather than fixed by the
medium (Gunawardena and Zittle 23).} Richardson and Swan, adapting the
same scale, found social presence correlating .68 with perceived learning
(73--74)---a second result from the same research program, extending
rather than independently confirming the first.\footnote{The correlation
is strongest in class discussions (.83) and group projects (.80), weakest
in lectures, notes, and readings (.40) (Richardson and Swan
74--75).} The meta-analytic aggregate holds the pattern at scale: across
26 effect sizes, social presence correlates 0.56 with satisfaction and
0.51 with perceived learning (Richardson et al.\
410--11).\footnote{Both correlations are moderated by course length,
discipline, and by which scale was used---no two social-presence scales
are equivalent (Richardson et al.\ 410--12, 415). The outcome throughout
the meta-analysis is self-reported perceived learning, never graded
performance (411); Tu's decomposition of the construct returned
reliabilities he himself calls ``medium-low'' and an instrument offered
as ``a first step'' (Tu 43--44).} Terry and Doolittle press hardest on
what the construct has become: the Community-of-Inquiry tradition
entangled social presence with the teaching and cognitive presences
around it, and interactivity keeps being mistaken for it. The repair they
propose is to identify it ``as a psychological state, as opposed to one
that is mediated through, or as a by-product of the use of a specific
technology'' (121).\footnote{Five factors contribute to it on their
account: group cohesion, co-presence and the salience of others,
affective association, immediacy, and intimacy (Terry and Doolittle
130--31).} That relocation is this section's load-bearing result: social
presence, like spatial presence in the first literature, belongs to persons.
It also sets the construct on ground Heidegger prepared: being-with-others is
not an addition to a self-standing subject but a constitutive structure of
human existence---the world is always already one shared with others---and
the distinctive enactment of that structure is
speaking-with-one-another.\footnote{``The world of Dasein is a
\textit{with-world} [\textit{Mitwelt}]. Being-in is \textit{Being-with}
Others. Their Being-in-themselves within-the-world is \textit{Dasein-with}
[\textit{Mitdasein}]'' (BT 155; italics Heidegger's). The Marburg lectures on
Aristotle name the enactment: ``The distinctive mode of being in
being-with-one-another is in speaking-with-one-another. To set forth the
possibilities of being-with-one-another is the \gk{ἔργον} of rhetoric''
(\textit{Basic Concepts} 91--92).} A design that opens room for
speaking-with-others therefore does more than decorate a channel: it
furnishes the enactment in which others are there at all. It is something a
design can activate, not something a medium contains.

The same tradition states what activating it requires. Garrison, Anderson,
and Archer locate social presence in the communication context rather than
the medium and name teaching presence ``[t]he binding element in creating
a community of inquiry'' (96)---when text-based education fails, their
diagnosis is a deficit of teaching presence, not of bandwidth. Kreijns,
Kirschner, and Jochems give the tradition its bluntest warning, naming the
first of two pitfalls of computer-supported collaborative learning:
``taking for granted that social interaction automatically takes place
just because an environment makes it technologically possible'' (336)---
the second being the restriction of what interaction does occur to
task-focused exchanges. Providing the channel is not producing the
contact; sociability must be designed for. The findings extend into
rendered worlds. De Back, Tinga, and Louwerse's immersed collaborative
environments produced learning gains with an inverse group-size
gradient---smaller groups, larger gains (5375--76).\footnote{Cited for
collaborative-learning design and outcomes only; social presence does not
appear in the study's text.} Van der Meer and colleagues' review of 139
virtual-reality collaborative-learning systems found 78.4\% of them
monitor-based against 7.2\% head-mounted (6).\footnote{Among the nine
design features the review requires of such environments, verbal audio
communication is judged essential and support for socialization is named
a primary function (van der Meer et al.\ 9--11).} Collaborative
learning in virtual worlds, as actually practiced and studied, already
runs on ordinary screens.

These two literatures share their origin and almost nothing after it.
Short, Williams, and Christie stand at the head of both: the presence
tradition the VR-pedagogy field inherited and the social-presence
tradition of online education descend from the same 1976 book. Downstream,
the citation traffic stops.\footnote{The meta-analysis that pools the
social-presence evidence draws its studies from online courses, not
rendered worlds, and pools neither Biocca nor Kreijns (Richardson et al.\
415--16); the desktop and headset studies of the first literature measure
spatial presence, cognitive load, and learning outcomes, and the
social-presence instruments appear in none of them.} This section reads a
desktop world whose students meet inside it by voice, and so it must join
two bodies of evidence that have not met each other---the second
literature's vocabulary applied to an object the first literature built.
What this movement hands forward is the
base for that work: a construct that names being-with as its own register
of presence; evidence that it predicts satisfaction and perceived
learning; the warning that it must be designed for and convened, never
assumed; and the finding that it belongs to persons---available, in
principle, to any world that gives its students a way to be there
together.

% --- M1 LOCKED 2026-07-20 (user). All four movements approved: M1.1 rev 3
%   -> merged M1.2/3 -> M1.4 rev -> M1.6 rev 2 (footnote rebalance
%   approved). DO NOT EDIT until the newest batch of sources has been
%   ingested (archon + snapshot rebuild), indexed (corpus/index entry
%   additions), and analyzed for useful/necessary material — then a
%   supplement pass on M1 under the same bank/verification workflow.
%   Deferred to unlock/assembly: venue/name pins in header comments;
%   Works Cited consolidation; ledger final check before M4.3/M5.2
%   drafting in the main session.

% ========== END INLINED FILE: DESKTOP-M1-DRAFT-v1.tex ==========

% ===== M2 (carries \section{The World and Its Deployment})
% NB M2.3 (method of reading) intentionally absent — blocked on arc decision.
% ========== BEGIN INLINED FILE: DESKTOP-M2-DRAFT-v1.tex ==========
% ============================================================
% Part III — Desktop Section — M2 (The World and Its Deployment) — DRAFT v1
% Mode: Stage 3 guided drafting (per-paragraph protocol). Plan approved
% 2026-07-20: Batch A = M2.1 artifact (3 ¶¶); Batch B = M2.2 cohort/instrument
% (2 ¶¶); M2.3 method-of-reading BLOCKED by Open decision #1 (arc placement).
% Conventions: as DESKTOP-M0-DRAFT-v2.tex (MLA; abbreviated-title parentheticals;
%   first person for K&S work; register + bans per memory).
% NEW SOURCE: King & Salvo, "Scalable Virtual Reality Global Clinical Immersion
%   for Culturally Responsive Engineering Design Skills" (ASEE paper #49715;
%   venue/year VERIFICATION PENDING via index agent; PDF pages 1-15, no journal
%   pagination — parenthetical cites use PDF pages). Platform facts pp. 3-4;
%   results p. 5; discussion p. 9. Ingestion + vle-05 unit creation running in
%   background agent.
% AUTHOR-CONFIRM FLAGS (Batch A): (1) RESOLVED (AUTHOR-SUPPLIED 2026-07-29,
%   assembly pass): deployed-build procedure list = Bilateral Scar Revision ·
%   PeriMortem C-Section Simulation · Transcatheter Aortic Valve Replacement ·
%   Deep Brain Stimulation · Spinal Deformation Surgery · Pediatric Auricular
%   Reconstruction · Cornea Transplant · Standard Cataract Procedures
%   (Paraguay) · Standard Cataract Procedures (Vietnam) · Standard and Laser
%   Cataract Procedures (UCI); each procedure includes its supplementary
%   video materials (animated models of procedure and/or equipment;
%   interviews with physicians and medical staff). Integrated into ¶1
%   2026-07-29 — REVIEW RULINGS (user 2026-07-31, batch 3): (i) CORRECTED —
%   UCI MC + CHOC cover SEVEN domestic items; the perimortem C-section
%   SIMULATION was recorded at the UCI Medical Education Simulation Center
%   (medschool.uci.edu, Division of Educational Technology; name verified
%   from the center's page). ¶1 restructured accordingly. (ii) CONFIRMED:
%   keep "spinal deformation surgery";
%   (2) RESOLVED (AUTHOR-CONFIRMED 2026-07-24): the launching-point/portal
%   structure was PILOT-ERA ONLY (rudimentary hospital). The deployed build
%   used by the 241 students = the full multi-story hospital; each procedure
%   room contains its own media player; the student walks in and plays the
%   video in place — no teleporting. Capture-verified 2026-07-24 against the
%   three videos at D:\PhD\Dissertation\Media\VLE (frames + transcripts in
%   corpus/index/Wendt.../bridge-sources/bme-vr-env-analysis/);
%   (3) RESOLVED (AUTHOR-CONFIRMED 2026-07-24): the deployed content library
%   STILL CONTAINS the animation and overview items (pilot-paper TAVR
%   inventory: "an animation of a TAVR procedure" + Edwards SAPIEN XT valve
%   overview; SAPIEN XT overview visible in-world in the 2024 group capture)
%   — outline M2.1's "device animations" stands in the artifact inventory;
%   (4) RESOLVED (AUTHOR-CONFIRMED 2026-07-20): Vietnam/Paraguay content was
%   present in ALL THREE surveyed terms.
% M2.1 TODO (AUTHOR-DIRECTED 2026-07-24): add a description of the video
%   delivery system — footage streamed from AWS S3 through the in-world
%   player, adaptive bitrate, ~3-second buffer on play and after scrubbing;
%   player opened with E, content under Procedures/Interviews tabs; 360°
%   footage with physician-POV picture-in-picture inset + Hide-POV toggle
%   (source: MediaPlayer.mp4 how-to narration + group-session transcript;
%   YouTube-streaming rejection reason CORRECTED, AUTHOR-CONFIRMED 2026-07-24:
%   it was NOT POSSIBLE to include YouTube streaming inside the UE
%   environment — a technical impossibility, not a weighed preference).
%   Related (AUTHOR-CONFIRMED 2026-07-24): YouTube-style conveniences such as
%   increased playback speed were DELIBERATELY omitted as a pedagogical
%   decision — feeds M3's prediction context.
% M2.1 (Batch A) — plan approved 2026-07-20; DRAFTED TEXT IN REVIEW (user
% reviews A + B in one pass).
% ============================================================

\section{The World and Its Deployment}

% --- M2.1 The artifact (Batch A, 2026-07-20 — in review)

The deployed platform is a downloadable application for Windows PCs, built in
Unreal Engine, that renders a navigable virtual hospital.\footnote{The platform
facts in this part are reported in a second of my publications, co-authored
with Prof.\ King and Kadin Diec, and cited parenthetically by abbreviated title
hereafter: King, Christine E., Kadin Diec, and Dalton Salvo. ``Scalable Virtual
Reality Global Clinical Immersion for Culturally Responsive Engineering Design
Skills.'' 2025 ASEE Annual Conference \& Exposition, paper \#49715.
% CITATION STATUS: year 2025 VERIFIED (title-page ASEE copyright line);
% conference name INFERRED from Paper-ID + copyright (not printed in body).
% Kadin Diec CONFIRMED as co-author by user 2026-07-20, spelling correct.
The paper has no journal pagination; parenthetical page numbers refer to the
PDF.} Students move through it in the first person, by mouse and keyboard, on
their own monitors. Each procedure room gathers the materials for one
operation or category of operations. The room's principal material is the
operating-room footage, which was captured two ways at once: a 360-degree
stereoscopic camera records the whole room, and a head-mounted camera worn by
the operating physician records the surgeon's own first-person view
(``Scalable Virtual Reality'' 4). Each room also holds recorded interviews
with the physicians and engineers involved and supplementary lectures and
readings (4), together with animated models of the procedure or its
equipment. The filmed procedures now span multiple institutions and
countries. Seven procedures were recorded at the UCI Medical Center and the Children's
Hospital of Orange County: bilateral scar revision, transcatheter aortic
valve replacement, deep brain stimulation, spinal deformation surgery,
pediatric auricular reconstruction, cornea transplant, and standard and laser
cataract procedures. A perimortem C-section simulation was recorded at the
UCI Medical Education Simulation Center. To these the platform adds standard
cataract procedures filmed in hospitals in Vietnam and Paraguay, placing the
same operation side by side in three healthcare systems (3--4). On an ordinary monitor the footage plays as an interactive 360-degree
panorama; the same materials are mirrored to a YouTube channel and website for
students whose machines cannot run the application (4).

The world is built for company. Voice chat is real-time and spatial---students
hear the classmates near them in the virtual space---and a group can convene
inside the environment, watch a procedure together, and share various other
artifacts through the screen sharing functionality: the design ``supported
real-time voice chat and screen sharing, enabling collaborative team viewing if
students wanted to meet in groups within the virtual environment'' (``Scalable
Virtual Reality'' 4). That design choice answers the firmest requirement the
collaborative-learning literature states: in van der Meer and colleagues'
review of 139 virtual-reality collaborative-learning systems, ``efficient
communicational tools are essential for effective collaboration: verbal
(audio) communication is crucial'' (9). The voice channel, that is, is not
one feature among others; it is the medium of speaking-with-others, the mode
in which a group is a group at all.\footnote{Heidegger, reading Aristotle's
definition of the human being as the living thing that speaks:
being-with-one-another ``is not being-with-one-another in the sense of
being-situated-alongside-one-another, but rather in the sense of
\textit{being-as-speaking-with-one-another} through communicating, refuting,
confronting'' (\textit{Basic Concepts} 33; italics Heidegger's).} Every
student is assigned to a procedure-paired group at the start of the term, the
capability is announced in the course materials, and meeting in the world is
encouraged. Group interaction inside the platform is one of the course's stated
practice objectives (3). Co-watching in the virtual hospital is thus not an
incidental feature but the designed use of the environment's social
architecture.

Every element of this environment is authored. The hospital and its rooms, the
two camera vantages, the interviews, the group assignments, and the channel
over which students speak are design decisions, and each is a made source---an
\textit{arch\=e}---of the experience it occasions. The virtual hospital is, in
the dissertation's terms, a manufactured and curated rhetorical ecology: a
designed \textit{Umwelt} whose makers supply the matter and the form of the
encounter while the act remains the student's own. The inventory above is
therefore also an inventory of design levers, and the analysis will return to
each element as the survey data show what it did and did not occasion.

% --- M2.2 The cohort and the instrument (Batch B, 2026-07-20 — in review)
% Facts verified against survey-analysis/00-data-audit, 01-quant-summary,
% 02-codebook-final: N per term 102/79/60; 14 items, 4 open (Q5/Q6/Q9/Q10);
% 873 open responses; strata APP 52/51/50% vs VID 45/46/43% per term,
% chi-sq=0.02, Cramér's V=0.009 (n=231 excl. small UNK) -> term-stable;
% reliability = batch 6, 143 responses double-coded, mean Jaccard .878,
% 4 manual adjudications; codebook v1 unrevised through coding.
% METHODS-TRANSPARENCY: RESOLVED 2026-07-23 (coding review Stations 6+12) —
% disclosure ruling (d) hybrid: plain in-prose clause + substantive footnote at
% the reliability sentence. ¶2 below FINALIZED per Station 12 approval.
% Post-review facts (2026-07-23): ALL canonical passes Fable 5 (Haiku batches
% 1/2/5 re-passed + replaced; cross-model Jaccard .596/.514/.668 documented);
% 7 manual adjudications; α-MASI .82 companion; per-code agreement 50–100%.
% Records: survey-analysis/02, 07-repass-impact-memo, coding/xmodel-*.
% Minimal coding-methods appendix APPROVED — draft at
% drafts/DESKTOP-APPENDIX-CODING-METHODS-DRAFT-v1.tex.

The survey corpus comes from three offerings of the course: Winter 2025 (102
respondents), Spring 2025 (79), and Winter 2026 (60)---241 students in all.
Each term used the same 14-item end-of-course instrument, four items of which
were open-ended; those four items produced the 873 open-text responses analyzed
in this section. One wording note applies throughout: the instrument asks about
the ``VR platform,'' phrasing inherited from the pilot era, though the
deployment it surveys is the 2D-PC-primary build described above. The
instrument also recorded whether students could run the application and whether
the platform elicited immersion, and those answers divide each cohort into two
strata: students who used the desktop application (52\%, 51\%, and 50\% across
the three terms) and students who watched the videos only (45\%, 46\%, and
43\%). The split is statistically indistinguishable across terms
($\chi^2 = 0.02$, Cram\'er's $V = 0.009$, $n = 231$), and that stability
carries weight in what follows: it licenses reading differences between the
strata as structural features of the deployment rather than accidents of a
given cohort. The course survey is itself anchored in the global deployment:
its title is ``Global Healthcare Needs and Final Assessment Survey,'' and its
three content items test the global material---nearly all students answered
them correctly. The dedicated global-procedures questionnaire---two Likert
items on exploring healthcare and on seeing cross-country differences, and an
open invitation to comment on the Vietnam and Paraguay videos---was a separate
Winter 2025 administration, reported in the ASEE paper above; two of its items
share their wording with this survey's platform-use and immersion questions.
Where its results are used, they are cited to that paper.
% [M2.2 ¶1 global-questions revision APPROVED 2026-07-20 — verified against
%  raw Canvas headers (all three terms) + ASEE paper Appendix B.]

% ****** USER MANUAL-REWRITE PENDING (ruled 2026-07-31, readthrough batch 3):
% ****** the author will personally rewrite THIS PARAGRAPH ("Responses were
% ****** anonymized...") AND its coding-passes footnote ("All coding
% ****** passes..."). Treat both as provisional; do not polish further until
% ****** that rewrite lands.
Responses were anonymized before analysis: names and identifiers were dropped,
and each response carries a term-and-respondent code. From the 873 open-text
responses I derived a coded corpus. The codebook was fixed before coding
began, and its categories come from the dissertation's own analytic
apparatus. It contains roughly forty codes organized into eight families,
each family collecting the codes for one kind of signature: forms of boredom;
flights away from the platform; the texture of reported time; the station of
the actualization chain a response addresses; the involvement channels a
response names; signatures of rhetorical incorporation; student suggestions;
and an open family for codes that emerged during coding. Coders read
each response segment by segment, a segment being the clause or sentence
carrying one codable idea, and recorded a set of codes per response, with
verbatim evidence spans retained for every application in the load-bearing
families. Coding was performed by AI coders working blind under fixed written
instructions; I adjudicated every disputed application myself.\footnote{All
coding passes in the final corpus were performed by Anthropic's Claude Fable 5
model, each pass blind and independent under written coder instructions
reproducing the codebook. Three batches were initially single-coded by a
smaller model (Claude Haiku 4.5) for economy; a same-instrument comparison
measured its agreement with Fable 5 at mean Jaccard .51--.67 against the .878
same-model benchmark, and I discarded those passes and had the batches
recoded, so no Haiku coding remains in the corpus. Per-code agreement on the
double-coded subsample ranged from 50\% to 100\%, weakest on union-merged
mention codes and strongest on the load-bearing families. Every cross-pass
dispute on a load-bearing code reached me for manual ruling, and I
subsequently reviewed the apparatus station by station---every rule, boundary,
and disputed application. A dedicated two-pass recount of the
shared-involvement evidence followed the same protocol (93\% exact agreement).
All pass files, comparison reports, and adjudication records are retained.}
All 873 responses were coded at least once, and a 143-response
subsample---drawn deliberately from the term the codebook had least seen---was
coded twice by independent blind passes, yielding a mean Jaccard agreement of
.878 over full code sets, the set-overlap measure appropriate to multi-label
coding (Krippendorff's $\alpha$ with MASI distance, .82). Merging the two
passes was deterministic and asymmetric: mention-level families merged by
union, the analytically load-bearing families by intersection, with the four
residual disputes resolved by hand; a subsequent line-by-line review of the
whole apparatus added three further manual corrections. Throughout, codes name
signatures in student talk, never certified inner states.

The corpus has four standing limits.
It is a single self-reported channel: no physiological, behavioral, or
performance measure accompanies it, so claims about inner states rest on what
students chose to write. It is retrospective: students answered at the end of
the term, so the corpus preserves remembered experience, not experience in
progress. It supports no learning-outcome claims: the course's content items
sat at ceiling, and nothing here measures whether one format taught better than
another. Term-to-term comparisons carry a composition caveat: the three cohorts
differ in size and possibly in makeup, so cross-term movements are read as
trends, not tests. What the corpus does provide, at a scale the pilot could
not, is 241 students' own accounts of entering---or failing to enter---the same
virtual world.
% --- M2.3 The method of reading — BLOCKED by Open decision #1 (arc placement;
%     deferred until lab-pole O-9 processing). Do not draft before the fork is
%     chosen.

% ========== END INLINED FILE: DESKTOP-M2-DRAFT-v1.tex ==========

% ===== M3 (carries \section{The Deployment Record} — title flagged, see (d))
% ========== BEGIN INLINED FILE: DESKTOP-M3-DRAFT-v1.tex ==========
% ============================================================
% DEFERRED REVISION — LAB-KEYSTONE DEPENDENCY (ruled 2026-07-29)
%   STATUS: queued, NOT to be executed until the laboratory boredom
%   section is drafted. Author ruling: working through the lab material
%   will likely revise the hypotheses and interpretations, so these
%   spans are revised ONCE, afterward, against settled findings —
%   not now, and not twice.
%
%   WHAT IS STALE: the rendered prose below cites the laboratory
%   keystone in its SUPERSEDED form — a brain measure reading
%   boring-like against a gaze measure reading engaged. The O-9
%   reprocessing (2026-07-16, N=8) withdrew it: P1's "clinical EEG
%   ~ boring" came from one clean subject and does not replicate
%   (clinical DMN is LOW = engaging-like; boring-like in 4/8 = chance).
%   The corrected keystone is ALL-PHYSIOLOGY-ENGAGED vs
%   SELF-REPORT-BORED (clinical self-report 6.0/9).
%
%   LAYER ASSIGNMENT (author-ruled 2026-07-29, governs the rewording):
%   physiology as a whole = the OCCUPIED SURFACE (outward comportment);
%   the post-video retrospective self-report = the HOLLOW DEPTH.
%   Grounded in fcm-05 SS24-28: at the dinner party we are "totally
%   involved" and DENY boredom, recognizing it only retrospectively;
%   "behavioural engagement is not evidence against boredom; in this
%   form it is the very medium of it."
%
%   SITES IN THIS FILE: see the per-site marker(s) below.
%   COMPANION SITES: DESKTOP-M3-DRAFT-v1.tex (M3.6) and
%   DESKTOP-M4-DRAFT-v1.tex (M4.1 P6); DESKTOP-SECTION-OVERLEAF-v1.tex
%   mirrors BOTH and is GENERATED — re-flatten after the edits, then
%   recompile DESKTOP-SECTION-ASSEMBLED-v1.tex (it \input-s M3/M4, so
%   the 37-pp PDF carries the stale claim until then).
%
%   CANONICAL SOURCES WHEN THE TIME COMES:
%     corpus/index/Heidegger - The Fundamental Concepts of Metaphysics/
%       _synthesis/fcm-king-salvo-bridge.md  (S4, revised 2026-07-29)
%     tmp/Dissertation/Part_III/reanalysis/boredom-o9-physiological-findings.md
%   NOTE: the structural point at both sites SURVIVES unchanged (a
%   survey holds one channel; the laboratory holds two). Only the
%   naming of the two channels is wrong.
% ============================================================
% ============================================================
% Part III — Desktop Section — M3 (The Deployment Record) — DRAFT v1
% Mode: Stage 3 guided drafting (per-paragraph protocol). M3.1 plan approved
%   2026-07-24; batch = M3.1 "The threshold" (3 ¶¶). Band R throughout;
%   hinge discipline: exactly ONE Band-G hinge sentence closes each
%   sub-module, naming the chain node; all argument deferred to M4.
% Conventions: as DESKTOP-M0-DRAFT-v2.tex / DESKTOP-M2-DRAFT-v1.tex (MLA;
%   respondent-code parentheticals for survey verbatims; chain nodes
%   A\textsubscript{n}; register + bans per memory).
% Data provenance: all rates/counts from t3-stats.md + 05-findings.md F1/F2
%   (all-Fable corpus, refreshed 2026-07-23). Verbatims verified char-exact
%   against data/open_text.csv 2026-07-24: W25-R042/Q9 (emoticon reproduced
%   verbatim — USER RULING 2026-07-24); W26-R059/Q9.
% Predicted-vs-unpredicted totals (¶2/¶3): predicted = lag 25 + nav 21 +
%   dark 10 + must-stand-close 2 = 58; unpredicted = crash 29 + install 26 +
%   hw 13 + A1.other 32 = 100. Excluded-at-threshold n = 61 (Q7-No ∧
%   video-only) = 25.3% of 241.
% Prediction-source framing: predictions registered in the coding design
%   (methodology-application/application-memo.md §4.2, fixed 2026-07-06,
%   before coding); Wendt is NOT named in section prose (not introduced by
%   the outline's M-structure; provenance lives in the analysis files).
% M3.1 — v2.1 APPROVED 2026-07-24 (LOCKED). M3.2 v1 DRAFTED 2026-07-24 —
%   IN REVIEW. Friction appendix DRAFT v1 written same day
%   (DESKTOP-APPENDIX-FRICTION-CATALOG-DRAFT-v1.tex; M3.1 ¶2 footnote
%   points to it). The M3.1 v1 rulings that produced v2 (style canon for
%   the rest of M3):
%   (a) REGISTER: too scientific — slow down, explain statistics in clearer,
%       more humanistic language; style sample presented for confirmation.
%   (b) NUMBERS: per-100-response rates unclear → switch to per-STUDENT
%       percentages (recomputed 2026-07-24 from 03-coded-corpus × respondents
%       q8 strata APP 123 / VID 108 / UNK 10): install 17% VID vs 3% APP;
%       hw 7% vs 3%; nonuse-talk 18% vs 0%; crash 15% APP vs 8% VID; lag 11%
%       vs 8%; nav 11% vs 6%. Raw counts stay. Define the two strata plainly
%       up front; "video route" → "YouTube route" (instrument's own wording:
%       "I utilized the YouTube videos only").
%   (c) ¶1: work through each obstacle in turn, one claim per sentence; the
%       user's model sentence for the Mac quote ADOPTED ("...excluded all
%       students who used Apple products, leading them to want it to be
%       '[m]ore inclusive for macbook users'").
%   (d) ¶2: grammar "the ones"→"those"; dev-team voice "we predicted five
%       elements..."; rename categories: VIDEO buffer latency (S3/AWS
%       streaming — delivery-system description noted for M2.1), video
%       fidelity/quality (not "dark-room"), close-proximity viewing; ALSO
%       predicted (AUTHOR-CONFIRMED 2026-07-24): startup crashes (UE5
%       performance-intensive), Mac-compatibility exclusion, missing YouTube
%       conveniences (playback speed omitted deliberately — pedagogical
%       decision). The 58-vs-100 predicted/unpredicted arithmetic is RETIRED.
%       Wendt DIRECT QUOTATION to include (user directive): breakdown-
%       anticipation, Design for Dasein pp. 33-34 (pdf 35-36), archon
%       exact-1.00+bbox verified: "anticipate the forms of breakdowns";
%       Bonsiepe qtd. in Wendt: breakdown "reveals the tissue of the
%       relations needed to fulfill the tasks".
%   (e) ¶3: de-stylize; thrust + final hinge sentence KEEP; tie in
%       world-disclosedness — the world FORECLOSED rather than disclosed for
%       technical reasons, not from any design aspect of the world itself.
%   (f) NEW APPENDIX ordered: every friction point, most→least reported
%       (scope proposal pending).
% ============================================================

\section{The Deployment Findings}
% [Title changed from "The Deployment Record" 2026-07-29, user-ordered at
%  assembly — "record" ban reaches the section title.]

% --- M3.1 The threshold (v2 REWRITE 2026-07-24, per approved style sample —
%     in review. Per-student percentages; strata defined in-prose; dev-team
%     prediction voice; Wendt quotation included; world-disclosedness close.
%     Data checks this session: 18 VID install-complainers, 15 of them Q7-No
%     ("problem was final" licensed); playback-speed request attested
%     W26-R021/Q9 + S25-R037/Q9; Wendt/Bonsiepe spans archon exact-1.00.)

The survey divides the 241 students into two groups, and the division itself
is the first finding. Each student reported whether they used the desktop
application or watched only the YouTube copies of the same footage: 123 used
the application, while 108---nearly half of the total participants---chose
the answer ``I
utilized the YouTube videos only.'' The YouTube group is itself two stories.
Sixty-one of its students reported that they could not use the application
at all; for them, YouTube was not a preference but the only door left open.
The other 47 reported that they could have used the application and turned
to YouTube anyway---a flight resulting from hassle rather than a failure to
access the platform.
The complaints make the difference between the groups concrete. Among the
108 students on the YouTube route, about one in six---17\%---described
problems installing the application or getting it to run, and for most of
them the problem was final: fifteen of those eighteen students are among the
sixty-one who never got in. Among the 123 application users, only 3\%
mention installation at all, and theirs are accounts of difficulty overcome,
since these students did, in the end, get the program running. The most
common obstacle was the operating system itself: the build ran only on
Windows, which excluded all students who used Apple products, leading them
to want it to be ``[m]ore inclusive for macbook users :(''
(W25-R042/Q9).\footnote{Survey verbatims are cited by term, respondent, and
survey item: W25-R042/Q9 is Winter 2025, respondent 42, question 9. The
three terms are Winter 2025 (W25), Spring 2025 (S25), and Winter 2026
(W26); respondent numbers were assigned during anonymization; question
numbers are the survey's own item numbers. Student responses are quoted
verbatim throughout, including typos and punctuation.}
The account gate stopped others before the hospital ever appeared; one
student's entire answer to the improvement question was ``[m]ake sure the
sign-in is actually possible'' (W26-R059/Q9). A smaller number were stopped
by the machine itself, reporting hardware too weak for the application or no
Windows computer to put it on (7\% of the YouTube group, against 3\% of
application users). One of them described the strain in physical terms:
``I think the VR program needs more tweaking, especially to work on a less
powerful computer, because, at least for me, my laptop sounded like it was
about to blow up after 5 minutes'' (S25-R033/Q6).

The application users tell a different friction story, and theirs happens
inside the world. Their complaints are those the analysis was designed to
test. Given the development and testing of the Windows build, we predicted
five elements where students might encounter friction in using the virtual
space: buffering delay on the video streams, since every video is delivered
over the internet from cloud storage rather than shipped inside the
application; wayfinding, the work of locating a particular room in a
multi-story hospital; video fidelity, the sharpness of the footage once a
room's screen is playing it; the close proximity to the screen that
comfortable viewing requires; and avatar occlusion, one student's avatar
blocking another's view in a shared room. Predicting where a designed thing
will break is part of the work of design itself. The design theorist Thomas
Wendt, following Gui Bonsiepe, treats breakdowns as ``a necessary
accompaniment to design,'' each one revealing ``the tissue of the relations
needed to fulfill the tasks'' (qtd.\ in Wendt, \textit{Design for Dasein}
33--34); the designer's charge is to ``anticipate the forms of breakdowns''
rather than to imagine them away (\textit{Design for Dasein} 33--34). The
survey returned a verdict on each prediction. Buffering and latency drew
complaints from 11\% of application users (25 complaints in all);
wayfinding from another 11\% (21 complaints), often paired with
requests for maps or clearer guidance; video fidelity from 5\%, whose
complaints of blurry or hard-to-see footage name the worry in the students'
own vocabulary. The predicted proximity problem barely registers---two
complaints---and avatar occlusion never appears, for the simple reason that
most students entered the hospital alone: with no other avatars in the room,
there was no one to block the view. Beyond these five we had also
anticipated frictions of a blunter kind: crashes at startup, given how
demanding Unreal Engine 5 is of an ordinary student laptop; the exclusion of
Mac users; and the absence of conveniences students know from
YouTube---above all changing the playback speed, a feature we deliberately
withheld as a pedagogical decision. All three arrived. Crashes are the
largest single friction in the corpus (15\% of application users; 29
complaints); the Mac exclusion did its work at the threshold; and students
asked for exactly the feature we had withheld: ``I think it would be helpful
to have the option to change the playback speed of the videos''
(W26-R021/Q9).\footnote{A complete catalogue of every reported friction,
ordered from most to least reported, appears in the appendix.}
% [appendix number at assembly; catalogue draft ordered 2026-07-24]

Almost none of this concerns the world itself. Taken together, the two
friction stories name the machinery around the virtual hospital---the
installer, the operating system, the account gate, the stream, the frame
rate---and almost never its architecture, its rooms, or its footage as
designed. For the students it stopped, that machinery decided everything: a
world that cannot be reached discloses nothing, and what was lost at the
threshold was world-disclosedness itself. The virtual hospital foreclosed
rather than disclosed, and it foreclosed on technical grounds---not through
any design aspect of the world itself. In the chain's
terms the finding touches A\textsubscript{0} alone: availability is in fact
its condition of possibility, and for a quarter of the class the world closed
at the installer, before anything in it could solicit so much as a look.

% --- M3.2 A world to see, nothing to do (v1 DRAFTED 2026-07-24 — in review.
%     ANCHOR F3; per-student stats computed 2026-07-24: D2-A3 166 apps /
%     147 students (61% of 241; APP 67% / VID 59%); D3-A3A4 37 apps /
%     33 students (15%/13%); G4-INTERACT 57 apps / 53 students (22%; APP 30%
%     vs VID 15%). Verbatims verified char-exact vs open_text.csv:
%     W25-R044/Q5+Q10, W25-R056/Q9, S25-R052/Q9, W26-R054/Q9 (student typo
%     "mor einteractive" — USER RULING 2026-07-24: quote verbatim incl.
%     typo, no [sic], per Mac-emoticon precedent), W25-R025/Q6, S25-R035/Q6,
%     W25-R058/Q6 (notes-outside-the-world anchor), W26-R027/Q9 ("occuring"
%     typo verbatim; inner double quotes → MLA single). Hinge = A3→A4.
%     orekton development EXPANDED per user ruling 2026-07-24 (ground in VLE
%     specifics + student responses). v2 of ¶2+¶3 written 2026-07-24:
%     same-student W25-R044 point made explicit; "One voice in 241" deleted;
%     "actually used the VLE" wording; action-talk defined in prose.)

Students praised what the world gave them to see. Mentions of the
content---the filmed operations, the physician interviews, the interest of
watching real surgery---are the largest positive channel in the whole
corpus: 166 coded mentions from 147 students, 61\% of the class. The praise
comes from both groups at nearly the same rate (67\% of application users,
59\% of the YouTube group), which follows from the deployment itself: both
routes carried the same footage, and what students valued was what the
footage showed. One student's summary is representative: ``the VR
simulation was super interesting and the wide-view videos of the procedures
being performed was very helpful'' (W25-R044/Q5).

Talk of doing is another matter. The coding distinguished students' talk
about the environment's content---what there was to watch---from talk about
acting in the environment: performing a task, making something, answering a
question, choosing a path. Talk of that second kind appears 37 times across
the 873 responses, from 33 students, and almost every instance is a request
or a complaint rather than a report: students describing not what they did
in the virtual hospital but what they wished it had given them to do. One
student put the wish in terms of routes: ``The VR clinical immersion
platform would be better if it was more interactable and had options on
which route you want to take rather than just being playback of videos''
(W25-R056/Q9). Another's entire answer to the improvement question was a
single request: ``Make assignments based on the VR clinical immersion''
(S25-R052/Q9). A third wrote: ``Maybe add something mor einteractive like
games/checklist/checkpoints etc'' (W26-R054/Q9). Requests of this kind came
from 53 students, 22\% of the class, and they came twice as often from
application users (30\%) as from students on the YouTube route (15\%): the
students asking for something to do are, above all, the students who
actually used the VLE. Among the 37 instances of action-talk, exactly one
reports an action the virtual hospital gave rather than an action it
lacked---and it comes from the same student, W25-R044, whose praise stood
for the content channel: ``The virtual space feels like you are in it,
while just a video doesn't. It makes you more involved and interactive with
the space'' (W25-R044/Q10).

What that missing act would have interrupted is long, unstructured
watching, and students describe that experience in detail. For a student
alone in the building---as most were---the world asks three things: find
the room, open its player, and watch. Nothing in the world asks anything
back. No room poses a question, no procedure ends in a prompt, and the
course's actual work of finding an unmet clinical need in the footage
happens outside the world entirely, in notes and documents the platform
does not include: ``I wish the VR platform would accommodate
other open windows. It wasn't easy to take notes on a Google doc while
listening to interviews or watching videos on the VR platform, so I
resorted to the youtube video afterwards'' (W25-R058/Q6). What remains
inside the world is watching, and the watching is long. One student found
``a lot of space of inactivity. Videos were unnecessarily long''
(W25-R025/Q6). Another asked for help with the seeking itself: ``I wish
there was a bit more guidance in watching the clinical procedures. Some of
them are many hours long, and I'm not sure what exactly to be looking for
during that time'' (S25-R035/Q6). A third met the same difficulty in the
course's own terms: ``at times it was hard to figure out a `problem' or
`unmet need' that was occuring. I think I was looking at the wrong things
in the beginning'' (W26-R027/Q9). These students are describing one
situation. Each arrived with attention to give and a task in hand---find a
need---and the footage gives every minute of the operation the same weight.
No moment is flagged as a difficulty; no prompt says \textit{watch how the
team works around this}; nothing distinguishes the places where an unmet
need shows itself---a workaround, a delay, an awkward instrument---from the
hours of routine that surround them. Seeking of the kind the course assigns requires a determinate
object---something in the world picked out as what the seeking is
\textit{for}---and the dissertation's term for that object is the
\textit{orekton}. The students quoted above report its absence in their own
words: interest solicited, seeking underway, nothing marked as the thing
sought. In the chain's terms the finding touches the passage from
A\textsubscript{3} to A\textsubscript{4}: the world reliably solicits
attention and interest, and the acts it furnishes stop at watching. It
comes down to a lack of affordances, of meaningful potentialities
for the students to actualize.
% --- M3.3 The flight to video (v2 REVISED 2026-07-24 per user rulings —
%     in review. MAJOR THEORY RULING 2026-07-24: the flight is NOT a chain
%     stall — the chain COMPLETES by another route; object of desire (the
%     content) unchanged; doxa + emotion at A3 select the ROUTE. Hinge
%     recast to A3 (route-selection); developed analysis ¶4 added BY USER
%     ORDER (hinge-discipline relaxed here). v3 THEORY REFINEMENTS
%     2026-07-24: chain completes through SAME mechanisms always (never
%     "another route" for the chain itself — route varies only for object
%     delivery); phantasia expansion added per user — discursive-
%     deliberative orientational mode weighs both platform-phantasmata →
%     committal doxa (taking-as-true) in which emotion concretizes
%     (vocabulary verified against Part II v4 chain spine, line 231).
%     SEVEN-CAUSES ADDITION 2026-07-24 (user order): new ¶5 — voluntariness
%     cut maps the two anatomies (61 = compulsion, arche outside; 47 =
%     reasoning on habit's floor, arche within; "calculated retreat" echoes
%     Section 6's own example); habit kept strictly distinct from hexis per
%     standing ban; quotes = Rhet. I.10 1368b33-1369a6 + I.11 1370a6-8,
%     wording verified against Section 6 Live.md. "virtually natural"
%     ARCHON-VERIFIED 2026-07-24 exact-1.00+bbox vs Aristotle -
%     Rhetoric_(2014)_[Clean Copy].pdf pp.30-31 (cite I.11 1370a6-8 per
%     Section 6). ¶5 v2 rulings applied 2026-07-24: reasoning spelled out
%     (weighed steps vs return); "deliberate" not "calculated"; yes-or-no
%     STRUCTURE (also harmonized in ¶2); could-not/origin-within sentence
%     expanded with the giving-up student. ¶5 v3 2026-07-24: "comparative
%     reasoning enabled by phantasia" + DA III.11 434a5-10 footnote (verbatim
%     from Part I §1.3 / verbatim_passages, verified 2026-05-22 entry);
%     crossed-line sentence anchored in arche/compulsion/voluntary
%     vocabulary. Hinge extended with origin-within clause. Batch is 5 ¶¶.
%     M4 COORDINATION NOTE (user: recalibrate after phantasia-expansion
%     approval): the A3 route-selection analysis + seven-causes agency
%     individuation now live in M3.3 ¶¶4-5 — M4 must build on them, not
%     repeat them.
%     ANCHOR F4. Stats: B1-FALLBACK 48 apps /
%     41 students (VID 26% / APP 11%); anatomies 61 excluded vs 47 chosen
%     (Q7×Q8). Verbatims verified char-exact + strata checked 2026-07-24:
%     W25-R099/Q6 (Q7-No, excluded; "accomadate" typo verbatim),
%     W25-R017/Q9 (Q7-Yes∧VID — certified chosen flight), W25-R002/Q9
%     (Q7-Yes, APP), W25-R041/Q6 (Q7-No BUT narrates choice — soft-boundary
%     evidence, presented as such). Zeitvertreib grounded in prose per
%     M3.2 orekton precedent. YouTube-mirror deliberateness = M2.1 platform
%     fact (DEP), not re-cited.)

The 108 students who watched on YouTube divide into two groups. Sixty-one
reported that they could not run the application at all; for them the
flight was forced, and the videos were the accommodation the course had
built for exactly their case. Their talk is regret as often as complaint: ``I wish
that I could access the vr setting on macbook because it feels like a great
opportunity that I missed but the videos helped to accomadate for me''
(W25-R099/Q6). The other forty-seven students reported that they could have
used the application and watched on YouTube anyway. For them the world
stood open and they declined it. Talk of switching to the videos appears 48
times in the corpus, from 41 students, and it comes from 26\% of the
YouTube group against 11\% of application users---those last being students
who entered the world and then left it.

The students who declined the world explain themselves in the language of
effort. One, from the group that could have used the application, watched
the comparison tilt over time: ``the VR videos were cool at first, but the
youtube format ended up being simpler to use and watch'' (W25-R017/Q9). An
application user performed the same calculation from inside the world,
room by room: ``The effort going into finding a procedure to watch is way
more than just clicking on the youtube video, so it would be better if the
room had more than just the viewing feed and something to do to make that
effort worth it'' (W25-R002/Q9). A third student shows how soft the line
between the two anatomies really is. On the survey's yes-or-no item this
student answered that they could not use the platform, yet their account of
it is a decision, not a defeat: ``The software. I ended up just giving up
on it and just watching the youtube videos. It's not worth the effort to
even learn how to use it and download it all when I could just pull it up
on youtube'' (W25-R041/Q6). For some students, that is, ``could not'' did
not name a hard technical wall---the student was a Mac user and could not
run the program at all, or the hardware was too weak to run the build---but
the point at which they stopped trying: capable in principle of getting the application working, they
abandoned the download, the installation, or the troubleshooting when the
effort began to exceed what the platform seemed likely to repay. The line
between the sixty-one and the forty-seven is therefore softer than the
survey's yes-or-no structure suggests; an unknown share of the excluded students
excluded themselves. All three students quoted above are doing the same
calculation. Each step the virtual hospital requires---download, install,
sign in, launch, walk, find the room, open the player---is a step the
YouTube mirror does not, and what waited at the end of those steps was the
same footage. One student asked for the steps to be abolished outright:
``Maybe not having the place in the game just directly clickable links to
open the videos. It was weird having to walk around just to find the right
video'' (W26-R047/Q9)---the walk that makes the world a place, experienced
here as pure overhead between the student and the footage. Friction sharpened the comparison where it appeared: mentions
of switching to the videos co-occur with install complaints eleven times,
with crashes five times, with latency four. One response carries the whole
sequence from occasion to flight: ``the platform was extremely slow in
terms of accessing the videos and interviews, which eventually became
frustrating and led to watching the Youtube videos'' (W26-R023/Q9)---the
friction named, the frustration concretizing, then the switch. A failed
install or a crash
was, in case after case, the moment a student switched. But the moment is
not the mechanism. The mirror was available before any crash and after
every one---always one click away, always carrying the same footage, never
asking for a download, an account, or a walk through a building first.
Friction supplied the occasions. What made every occasion sufficient was
the comparison that always stood behind it: at any moment of difficulty
inside the world, the same footage was already available by an easier
route, so the difficulty never had to be weathered---it only had to be
noticed.

Watching on YouTube was never against the rules. The course created the
mirror channel deliberately, so that no student would be left without the
material, and for the sixty-one students whose machines could not run the
application---the Mac owners, the weak laptops, the failed installs---it
did exactly the job it was built for. But a channel built as an
accommodation for students who could not enter the world also stands as an
alternative for every student who could. Heidegger's name for the everyday
counter-measure to boredom is passing the time
(\textit{Zeitvertreib})---``a driving away of boredom that drives time
on''---the reach for whatever occupation lies nearest when the present
one fails to hold us (\textit{Fundamental Concepts}
93).\footnote{Passing the time, Heidegger insists, is never taken up on
its own account: it is ``a passing the time which lays claim upon us
specifically out of and in opposition to a particular boredom''
(\textit{Fundamental Concepts} 93). The reach for the nearest occupation
is thus itself evidence of the boredom it answers.} The flight to video has exactly that
structure: not a verdict pronounced against the virtual hospital, but a
drift toward the nearest sufficient occupation---and here the drift was
institutional, because the course itself supplied the nearer channel and
stocked it with the same content. Whenever the virtual hospital offered a
student nothing beyond what the videos already offered---the same footage,
reached through more steps---the videos won, term after term, at a stable
43--46\% of each cohort. One student stated the pattern as a plain fact
about the class: ``The VR platform is too cumbersome to get working on
your computer, most students just resort to the online videos''
(S25-R034/Q6).

The dissertation's frame makes the mechanics of that drift precise. The
flight is not a failure of the chain of desire, nor even an unusual run of
it: the chain completed through the same mechanisms by which it always
completes, and what varied was not the chain but the delivery of its
object. What these students desired---the object their coursework required
and their interest tracked---was the content itself: the operations, the
interviews, the material they needed to view. Nothing bound that object to
the virtual hospital; the same footage was available through both
distribution platforms. So the chain ran to its end. The footage solicited
interest. Between that interest and the act stood a comparison, and the
comparison is the work of \textit{phantasia}: in its
discursive-deliberative orientational mode, \textit{phantasia} holds a
\textit{phantasma} of each path before the mind---the download, the
install, the walk through the building on one side; the single click on the
other---and runs through routes not yet taken, which is what makes a
reasoned weighing of alternatives possible at all. That deliberation issued
in the committal \textit{doxa}, the taking-as-true: ``simpler to use and
watch'' (W25-R017/Q9); ``not worth the effort'' (W25-R041/Q6). In that
ratification the emotion concretized---the frustration of the threshold,
the impatience with the hassle---giving the judgment its weight. Then came
the act: the student watched, and the desire for the content was
actualized, on YouTube.

The \textit{Rhetorica}'s sorting of action by its causes names what kind of
act this was---and for whom. Aristotle divides all action by whether its
origin lies in the agent or outside: chance, nature, and compulsion move a
person from without; habit, reasoning, anger, and appetite are the agent's
own (\textit{Rhetorica} I.10, 1368b33--1369a6). The sixty-one students the
platform shut out fled by compulsion: the \textit{arch\=e} of their flight
was the machine's refusal, an external wall standing against their own
\textit{orexis}---the regret in ``a great opportunity that I missed''
(W25-R099/Q6) is the sound of a desire overridden from without. The
forty-seven who could have entered and did not fled from causes within
their own power, and their words name which. What their responses record is
comparative reasoning enabled by \textit{phantasia}: these students held
both platforms before the mind, measured them by a single standard---the
coursework's need for the content against the effort each platform
demanded---and pursued what appeared the greater good.\footnote{The
operation is the one Aristotle assigns to deliberative \textit{phantasia}:
``whether this or that shall be enacted is already a task requiring
calculation; and there must be a single standard to measure by, for that is
pursued which is greater. It follows that what acts in this way must be
able to make a unity out of several images'' (\textit{De Anima} III.11,
434a5--10).} Their retreat was a deliberate one. It rested, moreover,
on habit's floor: YouTube was for every student the
often-done thing, the channel whose use had long since become, in
Aristotle's phrase, ``virtually natural'' (\textit{Rhetorica} I.11,
1370a6--8)---so the familiar route asked nothing of them while the
unfamiliar one asked much. One student located the asymmetry in
themselves, unprompted: ``I preferred the YouTube channel for virtual
reality, over the app. I'm not a video game fan, maybe is that''
(W26-R007/Q6)---the unfamiliar route is unfamiliar precisely to the
student without the gamer's second nature. The taxonomy also says precisely what the
survey's yes-or-no structure blurred: the line between the two anatomies of
the flight is the line between an \textit{arch\=e} outside the student and
an \textit{arch\=e} within. Students like the one who ``ended up just
giving up on it'' answered that they could not use the platform, but their
own accounts locate the \textit{arch\=e} of the flight elsewhere. Had the
machine truly refused them---a build their hardware could not run, an
install that failed no matter how long they worked at it---the cause of
their flight would have been compulsion: an origin outside the agent,
external circumstance overriding the student's own desire to enter, the
involuntary cause under which the genuinely excluded students fall. What
these students describe instead is a stopping-point they chose. The
install had not finally failed; they had stopped entering the VLE, on their
own deliberation that the additional effort was not worth its return. In the
\textit{Rhetorica}'s strict sense their flight was voluntary, the act their
own, whatever the survey answer records. The world, in the end, was not refused as
a world; it was weighed as a delivery platform for an object it did not
exclusively hold, and the effort it asked was too great for a destination
reachable in one click. In the chain's terms the finding touches A\textsubscript{3}: the
\textit{phantasmata} of both platforms taken up in the
discursive-deliberative mode and ratified by a \textit{doxa} in which the
emotion concretized---the chain completing as it always completes, deciding
not only whether the act would happen, but how, and from an origin that
was, for the forty-seven, entirely their own.
% --- M3.4 Presence happens on screens (v1 DRAFTED 2026-07-24 — in review.
%     ANCHOR F7 first half; hinge = A2 + world-disclosedness. 70.7% (87/123)
%     per standing correction, NOT the outline's ≈69%. Q8 item wording
%     quoted VERBATIM from raw Canvas header (W25 Student Analysis CSV,
%     item 3427111), same all terms per 00-audit. Pilot 70% verified
%     against p3-phenom-raw p. 388 ("70% of students noted the VR version
%     elicited a greater felt sense of presence"); cite form
%     ("Phenomenological Evaluation" 388) per M0 convention; pilot survey
%     respondent N=20 — prose avoids N to sidestep the 22-enrolled/20-
%     answered mismatch. F1-PRESENCE per-student 2026-07-24: 19 students,
%     APP 13 (11%) vs VID 6 (6%). G1-PROVISION 9 students; verbatims
%     W26-R014/Q9 + W25-R020/Q9 verified char-exact. W25-R044 same-student
%     continuity made explicit (per standing preference).
%     v2 REWRITE 2026-07-24 per user rulings: per-term percentages spelled
%     out; "deployed VLE"; scale = 20 vs 123 respondents stated; "through
%     traditional displays as much as 3D VR"; ¶2 expanded analytically —
%     world-disclosedness ("place dwelt in vs scene surveyed", Part II v4
%     L235 formula) + rhetorical incorporation (comes to full actuality at
%     committed-cognition→act; presence = "half its work") + phantasia
%     subtracts-the-surround (concept-lock framing); ¶3 opens direct;
%     pilot library-headset inversion FOOTNOTED (p. 389 verified, "not a
%     single student went there to check one out" verbatim);
%     TERMINOLOGY-LOCK COMPLIANCE (memory: rdr2-secondary-integration —
%     "phantasia-load"/"data-discrepancy"-in-analyses/"image-making" ALL
%     BANNED in analytical prose; earlier claim that v4 L205-207 licenses
%     them was WRONG for analyses — those are metaphysics-section coinages;
%     approved framing = medium-supply vs phantasia gap-filling /
%     subtract-the-surround): division-of-labor prose carries the concept
%     with phantasia only; hinge recast at the phantastic motion rendering
%     A2 (world taken-as dwelt in).
%     v3 EXPANSIONS 2026-07-24 (user rulings): same-display point added
%     (¶2 — difference is DESIGN not hardware; film can world-disclose;
%     phantasia active on both routes); incorporation-debt worked through
%     in own ¶ (disclosure = world's half of exchange; act ratifies;
%     begun-and-suspended; teleology = disclosure for action as means for
%     ends); NEW ¶ — matched-numbers insight (70% headsets vs 70.7%
%     monitors → immersive power from explorable-world inclusion, not
%     stereoscopic viewpoint) + Calleja: definition ARCHON-VERIFIED
%     exact-1.00+bbox (In-Game 169, pp.168-170 chunk) + spatial/kinesthetic
%     cornerstone pair (In-Game 170, digest ig-10 line-anchored); bare
%     "incorporation" for Calleja per standing rule. Batch is now 5 ¶¶.
%     Wendt considered for this insight, Calleja fit exactly (screen-game
%     provenance).)

Each term the survey put a direct question to the students who used the
application: ``Did you find the VR clinical immersion platform to elicit a
greater degree of immersion, presence, and/or embodiment (the feeling that
you are actually present in the operating rooms/clinical environment)?''
Of the 123 students who used the application and answered, 87 said
yes---70.7\%. The majority held in every term taken alone, not merely in
the pooled total: 64\% of Winter 2025's application users answered yes,
82\% of Spring 2025's, and 67\% of Winter 2026's. The figure lands with a
history behind it. In the pilot, we had asked whether the 3D VR versions
elicited a greater degree of presence than the other formats, and 70\% of
respondents said they did (``Phenomenological Evaluation''
388)---students wearing headsets, standing inside stereoscopic footage.
The deployed VLE reproduces the pilot's number almost exactly, with no
headset involved: where the pilot's question was answered by twenty
students in a single quarter, the deployment's immersion item was answered
by 123 application users across three terms, every one of them working
through an ordinary monitor by mouse and keyboard. Presence, as the
students report it, happened through traditional displays as much as it
did through 3D VR.

The yes-answers do not stand alone; the open responses corroborate them
in the students' own words. Nineteen students, unprompted by any
yes-or-no item, wrote descriptions of feeling present in the virtual
hospital---of being in the space rather than watching it---and thirteen
of the nineteen were application users (11\% of that group, against 6\%
of the YouTube group). The fullest of these descriptions comes from the
same student, W25-R044, who supplied the corpus's single account of given
action: ``The virtual space feels like you are in it, while just a video
doesn't'' (W25-R044/Q10). The sentence repays slow reading, because in a
student's plain words it draws the dissertation's own distinction---and
draws it under conditions that isolate what makes the difference. Both
routes reach this student through the same hardware: the YouTube videos
and the virtual hospital alike arrive on an ordinary computer display.
Whatever separates being-in it from watching it therefore cannot be the
screen; it is the design---how the same footage is framed and delivered
to the student who watches. Nor is traditional video barred from
disclosing a world: a film seen on that same monitor can take its viewer
into its world, and \textit{phantasia} is active on both routes. The
difference lies in what the design gives \textit{phantasia} to work
with. The YouTube route delivers the footage as one video among
videos---a page, a player, a list of what to watch next---content framed
as content, the surround reasserted at every edge of the frame. The VLE
delivers the same footage as the interior of a place: to reach it the
student has walked a building, found a room, and stood before its
screen, and the hospital continues around the footage while it plays.
What the student's statement registers is the difference between those two
relations---a world disclosed as a place dwelt in, against a scene
surveyed from without. That difference is what the dissertation names
world-disclosedness: the student taken into the world acted in, and the
world taken into the student's awareness. On a flat screen that
disclosure is \textit{phantasia}'s victory in either case. A monitor
gives the eyes a bounded rectangle inside a room that remains fully
present around it, and it is \textit{phantasia} that subtracts the
surround, supplies what the screen withholds, and encloses the student
in the world the screen only samples. The designs differ in how much
they help: the virtual hospital hands \textit{phantasia} a surrounding
place already built---an inside to be-in---where the video page hands it
a frame to be overcome. The coded responses register that difference in
help directly: of the nineteen students who described being-in, thirteen
came from the 123 application users and six from the 108 on the YouTube
route---students given the built world reported being-in at nearly twice
the rate of students given the framed video, though both watched the
same footage on the same kind of screen. A second description extends the
being-in past place into role: ``I loved the VR videos; I feel that it
actually made me feel like an engineer trying to find and solve problems
that real doctors are facing'' (W26-R007/Q5). The being-in here is
vocational as much as spatial: the student reports becoming the engineer
the course wants them to be, the seeing already oriented toward the
problems real doctors are facing. The same student, in another answer,
preferred the YouTube route, tracing the preference to being ``not a video
game fan'' (W26-R007/Q6)---a reminder that the presence the footage
occasions and the route a student settles on are separate verdicts,
delivered here by one voice.

This disclosure is also the first face of what the dissertation calls
rhetorical incorporation. The term names the union a designed world
achieves with its inhabitant, and the actualization chain states the
condition on which that union is fully achieved: the world solicits
perception (A\textsubscript{1}), \textit{phantasia} composes the charged
image of a surrounding place (A\textsubscript{2}), the student's
committed judgment takes the world as worth engaging
(A\textsubscript{3}), and the union is fully actual only when that
judgment passes into an act the student performs in the world
(A\textsubscript{4}).
It is worth working through what that means for these students. Feeling
present is not the completion of rhetorical incorporation; it is the world's half
of an exchange. The chain is understood backward from its end---from
the desired object whose pursuit completes it in an act---so a world
that has disclosed itself has, so far, only made acting possible: it
has given the student a place in which an act could happen. Rhetorical
incorporation is finished when the student's own act, done in and with
the world, ratifies the disclosure---when being-in becomes acting-in. A
world that discloses but does not task therefore leaves the exchange
half-done. The student stands in a place that asks nothing of them; the
disclosure that invites completion remains an invitation; rhetorical
incorporation is begun and then suspended at the very threshold where it
would become actual. That is the precise sense in which such a world still owes its
students an act: the debt is written into the chain's own teleology,
where disclosure is for action as means are for ends. Other students
put the being-in more modestly---``Being able to get somewhat first
hand observations'' (W25-R089/Q5); ``I found the immersive VR helpful''
(W26-R028/Q5)---but the relation they report is the same. One more names
what the being-in held off: ``It was different for each place and being
able to make it seem like we were in the procedure and watching made it
more immersive rather than boring'' (S25-R028/Q5)---immersion, in this
student's own accounting, standing where boredom would otherwise stand.

Nine students asked the course to provide VR headsets: ``This class
should provide students with VR headsets. Even if there were a technology
fee, it would be more beneficial for learning'' (W26-R014/Q9); ``I feel
like we should have some sort of cheap VR set because the VR experience
feels off when seeing it in the laptop screen''
(W25-R020/Q9).\footnote{The request inverts the pilot's most curious
result. In the pilot, we housed several Meta Quest 2 headsets in the
school's library, and ``not a single student went there to check one
out''; as one student phrased it, obtaining a headset ``was a hassle''
(``Phenomenological Evaluation'' 389). Three years on, students ask for
headsets to be handed to them. The desire for the richer medium is real
in both cohorts; what governs it is the same convenience gradient that
routes everything else in this deployment---intensity is wanted, but only
at zero fetching cost.} What these nine notice is the division of labor
just described. Every medium supplies some share of what the senses would
receive in the non-virtual scene, and \textit{phantasia} must furnish the
rest: a headset supplies much and asks little of \textit{phantasia}; a
flat screen supplies least and asks most. The student who writes that the
experience ``feels off\,\ldots\,in the laptop screen'' is reporting how
large \textit{phantasia}'s share becomes on the desktop route---feeling,
in the moment of viewing, how much of the world it falls to their own
\textit{phantasia} to hold up. What the 70.7\% then says is that the
share was held. Seven in ten application users, on the route where the
medium supplies least and \textit{phantasia} must supply most, reported
the feeling of actually being present in the operating room.

Set the pilot's number and the deployment's side by side, and they
sharpen into a finding about where immersion comes from. The pilot's
students wore stereoscopic headsets and stood inside the footage; the
deployment's students sat before flat monitors; both cohorts reported
presence at 70\%. If the stereoscopic viewpoint were what
carried immersion, the number should have fallen when the headsets
disappeared---and it did not fall. What the two cohorts shared was not
display hardware but a designed, explorable world: the virtual
hospital, entered and navigated, on whichever display. The likeliest
reading of the matched numbers is therefore that the immersive power
\textit{phantasia} conveys derives not from the stereoscopic viewpoint
but from the inclusion of the footage in an explorable world. Calleja's
account of incorporation, built from screen games---players at
monitors, no headset anywhere---predicts exactly this. He defines
incorporation as ``the absorption of a virtual environment into
consciousness, yielding a sense of habitation, which is supported by
the systemically upheld embodiment of the player in a single location,
as represented by the avatar'' (\textit{In-Game} 169): habitation grows
from emplacement in a navigable world, not from any property of the
display. Of the six dimensions of involvement his model blends into
incorporation, the cornerstone pair without which it cannot occur is
the \textit{spatial} and the \textit{kinesthetic}---the internalized
space and the learned movement through it (\textit{In-Game} 170)---and
that pair is
precisely what the VLE adds to the footage and YouTube does not. Behind
the mechanism stands a principle of Merleau-Ponty's: ``What counts for
the orientation of the spectacle is not my body, such as it in fact
exists, as a thing in objective space, but rather my body as a system
of possible actions, a virtual body whose phenomenal `place' is defined
by its task and by its situation. My body is wherever it has something
to do'' (\textit{Phenomenology of Perception} 260).\footnote{Wendt
reaches this same passage from the design side: for him it carries
direct consequences for the design of experience, which must orient
itself to what users are there to do---their goals, conscious and
non-conscious---within the multiple stable uses any technology admits
(\textit{Design for Dasein} 126). The convergence is the point: what
the phenomenologist states of the body, the design theorist prescribes
for the world built around it---a world is designed well when it gives
the virtual body something to do. A third tradition arrives at the same
structure independently. Chalmers, with no reference to Merleau-Ponty,
builds his defense of virtual realism on the same figure: ``If I pick
up a virtual coin in Second Life, I really use my virtual body to take
possession of a virtual coin. There is nothing fictional about this''
(``The Virtual and the Real'' 317); ``virtual actions are plausibly
real actions (albeit with a virtual body), so that when one performs
virtual actions, one really is doing something'' (339); and a user's
sense of ownership over a virtual body ``need not always be an
illusion'' (333). Phenomenology states the principle---the body is
where its tasks are; design prescribes it---build worlds that give the
virtual body something to do; analytic metaphysics certifies its
object---what the virtual body does is really done. For the survey's
70.7\% the triangulation matters: the students' reports of being
present in the operating room need not be read as a flattering
illusion, since they describe a virtual body genuinely located by task
and situation. It sharpens the deployment's deficit in the same stroke,
for a virtual body whose acts are real still needs acts to perform.}
% [Chalmers cites verified char-exact vs corpus PDF (pp. 317/333/339,
%  journal pagination); PDF marked 2016, published Disputatio 2017 —
%  settle year at Works Cited. ADD Chalmers TO WORKS CITED.]
% [CITE UPGRADED 2026-07-24: direct quotation from user's edition —
%  Merleau-Ponty, Phenomenology of Perception, Landes trans. (Routledge
%  2014 [1945]), p. 260, Part Two "Space" (Wertheimer passage); wording
%  verified via pdftotext (Landes: "wherever it has something to do",
%  NOT Smith's "wherever there is something to be done"). ADD TO WORKS
%  CITED. Wendt footnote = PARAPHRASE gloss (digest's "major
%  implications" phrase did NOT verify verbatim — no quotes used).
%  PENDING OPS: archon ingest of the new PhP PDF + snapshot rebuild
%  (vector-compact) queued for next ingestion session — OP RULE.
%  CORRECTION: Chalmers "Virtual and the Real" does NOT quote this
%  passage (earlier archon hit = partial-phrase match on his own
%  "virtual body" usage); his virtual-body realism converges but does
%  not cite M-P. CONTINUATION NOTED for possible later use (PhP
%  260-261): possible actions "sketch out...a possible habitat";
%  "the reflected room conjures up a subject capable of living in it.
%  This virtual body displaces the real body" — meets Calleja's
%  "sense of habitation" almost verbatim.]
The explorable hospital gives the student's virtual body somewhere to
be---a place constituted by situation and task, a room to find and a
player to open---on whatever display the world arrives through. In the
chain's terms the finding sits at the \textit{phantastic} motion that renders
A\textsubscript{2}: from the thin \textit{aisth\=emata} a monitor
delivers, \textit{phantasia} still composes the charged
\textit{phantasma} of a surrounding world---a world taken-as dwelt
in---and it composes it from the explorable place the design supplies,
not from the apparatus on the head.
% --- M3.5 The WE: had, wanted, unactivated (v1 DRAFTED 2026-07-24 —
%     in review. ★ SPOTLIGHT MODULE. LOCKED-SPEC COMPLIANCE: mention-floor
%     phrasing throughout; NO rates (no denominator — instrument never
%     asked; ceiling restated in ¶1, M6 carries consequence); ACTIVATOR
%     stated SINGULAR (W26-R037/Q10 sole confirmation); non-activation =
%     evidenced pole (11); OPT-MANDATE 4 / OPT-CAP 13 / awareness-gap 12
%     (W25-R084 excluded per author ruling — genuine feature request);
%     gap = activation ∧ awareness; coordination-cost = the hinge, stated
%     as the survey's own arithmetic, theorizing DEFERRED (no module
%     labels in text). Disclosure: recount passes covered by M2.2 ¶2
%     footnote (option d) — no methods restatement here. Verbatims
%     verified char-exact vs open_text.csv 2026-07-24: W25-R038/Q10,
%     S25-R072/Q10, W25-R014/Q10, W25-R089/Q6, W26-R037/Q10, W26-R052/Q9,
%     W25-R060/Q10, S25-R040/Q10, W25-R071/Q9, W26-R022/Q10; S25-R043
%     instructor-session = paraphrase (abridged text only in 06 queue).
%     Platform facts = M2.1 DEP (restated in one clause, not re-described).
%     M3.6 COORDINATION NOTE: W26-R037 (the singular activator) is ALSO
%     F5's best within-respondent divergence pair AND an F1-PRESENCE row —
%     M3.6 should surface that continuity when it presents the pair.
%     Recount mechanics (115 candidates = 67 pre-refresh E-SHA ∪ 48-row
%     keyword sweep; E-SHA now 91 post-refresh; recount IMMUNE per Station
%     11) kept OUT of prose — floors are the adjudicated FINALs.
%     v2 RULINGS APPLIED 2026-07-24: ¶1 course-materials clause + the
%     AUTHOR-SUPPLIED first-person methods admission ("we did not think
%     to ask...falsely assumed...") VERBATIM per user 2026-07-24; floors
%     explained plainly (minimums; invisible non-mentioners); recount
%     sentence per user wording + four-kinds preview. ¶2 semicolon after
%     W25-R089/Q6 cite (user cited W25-R014 but the em-dash was at R089 —
%     applied there); "course design intended and which was heavily
%     encouraged" (user wording). ¶4 REBUILT: no cost/price/ledger/
%     arithmetic metaphors; "demand none of the solo uses make";
%     counts stated directly; "YouTube platform"; co-disclosedness/
%     world-disclosedness relation spelled out (students present-but-
%     alone, singly before monitors); plain closing sentence.

% ****** FIGURE TODO (user 2026-07-31, batch 4): include a picture showing
% ****** multiple people meeting inside the VLE (group co-watching scene);
% ****** candidate source: the 2024 group-session capture (Media/VLE).
\begin{center}
\fbox{\parbox{0.9\linewidth}{\centering\textbf{****** FIGURE PENDING
******}\\ Insert screenshot: multiple students meeting inside the VLE
(group co-watching scene). Candidate source: the 2024 group-session
capture.}}
\end{center}
Of everything the virtual hospital offers, one affordance has no YouTube
counterpart at all: the VLE is built for watching together. Students
hear the classmates near them in the space, can share their screens
there, and every student belongs to a procedure-paired group from the
term's first week; the designed use of the environment is that the
group convenes inside the world and watches as a group, and that
multiplayer functionality is clearly stated and encouraged in the
course materials. What the survey can say about that design is limited
in a way the other findings are not. Unfortunately, we did not think to
ask the students whether or not their group convened in the world, so
no rate of use statistics are available. Given the group based nature
of this class, we falsely assumed students would organize meetings
within the VLE of their own volition, and that false assumption led the
course and survey designers to ignore crucial questions to be asked and
course designs that could have mitigated such. What the 873 responses
hold instead are incidental mentions---remarks students happened to
volunteer about watching together or about being alone---and counts
built from such mentions can only set minimums. That one student
described co-watching means at least one group convened; it cannot mean
only one did, because students who met and never mentioned it are
invisible to the survey. Each count in what follows is therefore a
floor---the minimum established by those who happened to speak---and
never a rate. After reviewing the relevant 115 candidate responses by
what each actually reports and sorting them, they fall into four
different kinds: reports that the group meeting never happened; a
report that it did; requests that the course structure the meeting; and
requests for a co-use capability that, in nearly every case, the
platform already possessed.

The best-attested kind is absence. Eleven students report directly that
the meeting did not happen for them. One logged in and found the
building empty: ``I found that when I did log in, I was usually alone
in the platform so it didn't really make me feel more connected with my
peers'' (W25-R038/Q10). One's group met, but elsewhere: ``The virtual
space to meet up for teamwork was very smart, but my group utilized
discord a lot'' (S25-R072/Q10). One knew the capability existed and
never used it: ``Meeting up with other people in the virtual space is
an interesting idea but because I didn't actually experience it all I
can say is it seems cool in theory'' (W25-R014/Q10). A fourth asked for
``more of a community since it just seemed isolating''
(W25-R089/Q6); a sentence that reads as a complaint about the design
until it is set beside the platform facts: the community infrastructure
existed, and what this student experienced was a groupmate-less
default, not an absence in the world. The isolation reported in these
responses is real. Its cause, in every case the responses let us see,
was that the assigned group never convened where the course design
intended and which was heavily encouraged.

Against these eleven stands one confirmation that the designed use
happened. In all 873 responses, exactly one student reports co-watching
as an experienced feature of the platform: ``It is more effective that
it does put you into a VR clinical environment and a cool interactive
space. The main pro to the platform is you can watch with others''
(W26-R037/Q10). No second response confirms it. The rest of the WE talk
is wishing, and the wishes split in two tellingly different ways. Four
students ask for the meeting to be structured---not for the capability,
but for someone to carry the organizing: ``Incorporating a required
group meetup within the VR simulation could incentivize more
interactive learning'' (W26-R052/Q9); a class sheet ``where times are
set by each group to enter the VR program'' (W25-R060/Q10); ``remote
teamwork VR space weekly meetings, to ensure accountability''
(S25-R040/Q10); an instructor logged in at the same time as the
students, explaining as they watch. Thirteen more ask for co-use
capability in wording that implies it does not exist---and for twelve
of the thirteen, the request is for something the platform already has.
``I think if there was a way to interact with other users who are
online at the same time that would add to it'' (W25-R071/Q9); ``perhaps
the ability to meet up with groups in a virtual reality setting''
(W26-R022/Q10)---the platform's existing capability, requested as an
innovation. The gap the survey shows is therefore double:
activation---groups that could have met and did not---and
awareness---students who never learned the meeting was possible at all.

Read together, the counts show why the meeting so rarely happened.
Every other use of the platform asks one student to act alone:
download, install, launch, walk, click play. Watching together is the
single use that requires two or more students to act at the same
time---agree on a time, log in together, find the same room---a demand
none of the solo uses make. The counts reflect that demand directly:
one student's responses confirm a group that met in the world; eleven
report groups that never did; and seventeen ask that the organizing be
done for them, by the course, the instructor, or a posted schedule. The
same pull of convenience that sent students to YouTube worked against
the meeting as well, and it worked hardest against the one affordance
the YouTube platform cannot reproduce: watching together, in one place,
at one time. In the dissertation's terms, what stood ready on this
route was co-disclosedness---the shared counterpart of
world-disclosedness. The students who described feeling present in the
virtual hospital were each alone in it: taken into the world, but
singly, one student before one monitor.
Co-disclosedness is that same being-taken-into-the-world had together
with others who are present in it at the same time, and it is the half
of the disclosure the platform's social architecture was built to make
possible. The survey states its standing directly: every student had it
available by design, and one student's responses confirm that it was
ever had.
% --- M3.6 From friction to justification (v1 DRAFTED 2026-07-24 — in
%     review. ANCHORS F5+F6. Trend-not-test stated FIRST (M2.2 DEP).
%     Sweep numbers canonical (02 §Emergent): 16 total, 5/2/9 by term =
%     1.3/0.7/4.2 per-100; sweep-methods transparency FOOTNOTED (artifact
%     code; 4 independent blind rediscoveries; circularity caveat).
%     Friction declines: crash 21→4→4, hw 8→4→1 (t3). A4-DIV 13
%     candidates / 12 unique (dup once). Residue = W26-R016/Q6 +
%     S25-R009/Q10 per F5. W26-R037 continuity surfaced (M3.5 activator =
%     divergence pair author); W26-R022 thread surfaced (skeptic proposes
%     the existing WE capability). Second-form gloss sourced from
%     construct-crosswalk §1 keystone vocabulary (occupied surface /
%     hollow depth; two-channel disagreement), NOT Part II v4, per
%     standing caution; NO German form-names (unverified); echoes-not-
%     detections ceiling stated. Verbatims verified char-exact
%     2026-07-24: W26-R033/Q9, W26-R056/Q9 (entire answer), W26-R022/Q10,
%     S25-R022/Q9, W25-R014/Q9 ("dont" typo verbatim), W26-R016/Q6,
%     S25-R009/Q10, W26-R037/Q9+Q10 ("then" typo verbatim). Hinge =
%     the justification question; theorizing deferred, no module labels.

The deployment's last finding lies in the difference between its
terms, and it must be read modestly: the three cohorts differ in size
and possibly in makeup, so what follows is a trend to be described,
not a test that was passed. The trend begins with the frictions
falling. Crash complaints numbered 21 in Winter 2025 and 4 in each
term after; hardware complaints fell from 8 to 4 to 1. By its third
term, the platform largely worked. What rose as the frictions fell was
a different kind of talk---talk that questioned the platform's value
rather than its function. A keyword search, run with the same
fixed list of search terms across all 873 responses so that every term's
cohort was screened identically, finds 16 responses of this kind: five,
two, and nine across
the three terms, which is 1.3\%, 0.7\%, and then 4.2\% of each term's
responses---the largest share by far in the final and smallest
cohort.\footnote{The counts come from a
keyword sweep applied uniformly across all three terms rather than from
the coded label for this theme, because that label was invented mid-run
by one batch's coders and its all-Winter-2026 distribution reflects
coder vocabulary, not the terms. The sweep found 16 responses (two
further keyword matches were inspected and discounted as false
positives). The theme itself is not an artifact: independent blind
coding passes rediscovered it in Winter 2025 under other names. One
caveat stands---the sweep's vocabulary is partly seeded from Winter
2026 phrasings, so the exact size of the rise is soft, though its
direction is robust.} The late voices are not reporting failure. ``I
think the youtube videos were enough to fully grasp the experience''
(W26-R033/Q9). The register is not confined to the final term; its fullest
argument came in the middle one: ``the vr learning platform wasn't
accessible, and i think if youtube videos do the job, they do the job.
some things don't need to be improved on; take the paperclip for example''
(S25-R039/Q6)---sufficiency argued by analogy to a perfected design.
``Less dependency on the VR aspect altogether''
(W26-R056/Q9)---that is one student's entire answer to the improvement
question. A third goes past sufficiency to a failure of imagination about
any use at all: ``I do not think this one should ever be mandated in the
future. It should be an extra option. I cannot imagine a time when a VR
space would be more helpful'' (W26-R055/Q10). Traditional learning, a
fourth wrote, is ``the least gimmicky'' way (W26-R022/Q10). One voice adds
a complaint of another kind entirely---not that the world is unnecessary
but that it fails to look serious: ``The VR clinical immersion lobby is
modelled poorly and it reminds me of a children's game. I didn't find it
to be very academic? The initial impression really turned me off of using
it and the actual videos were disorienting'' (W26-R024/Q9)---the game-like
look that makes the world enterable, undermining at first sight its
authority as a clinic. The complaint has changed kind:
in the first term students told us the world did not work; by the
third, with the world working, they asked why it should exist.
% ****** USER MANUAL-REWRITE PENDING (ruled 2026-07-31, batch 4): the author
% ****** will personally rewrite the keyword-sweep footnote above ("The
% ****** counts come from a keyword sweep..."). Treat it as provisional; do
% ****** not polish it further until that rewrite lands.

A small set of responses sharpens that question: in one and the same
answer, these students praise the platform and report that something
about it felt empty. Thirteen responses carry the pattern. Two of
them---submitted in different terms, under different respondent
codes---are word-for-word identical, so the analysis counts that text
once, which leaves twelve distinct responses that praise the platform
while finding it hollow. Read closely, most of them resolve
into ordinary complaints, because the hollowness has a named culprit.
``The actual idea of having the VR clinic is amazing, the videos and
images that come with it not so much'' (S25-R022/Q9)---the culprit is
the footage. ``The VR is a cool idea but the quality is so low I don't
gain from the atmosphere and I dont feel immersed'' (W25-R014/Q9)---the
culprit is the quality.\footnote{The perceived lack of quality has two main
sources. First, the program does not detect and automatically apply the
native resolution of the student's laptop; the student must set it manually
in the graphics settings menu. Students are provided with onboarding videos
that state this clearly and walk through how to adjust the settings to
optimize quality and performance for their specific hardware, but it seems
many never actually do so. Second, the quality can suffer from limited
internet bandwidth. Each operation's footage comprises two or three
simultaneous video streams: the 180- or 360-degree stereoscopic view of the
operating room, the physician's point-of-view video, and, in some cases,
recordings from surgical equipment such as the systems used to conduct eye
surgeries. The playback system adapts to each student's bandwidth, serving
lower-resolution footage when bandwidth is limited, but playing the streams
at 4K resolution requires well over 100\,Mb/s. The same videos uploaded to
YouTube must be baked together into a single stream, so only one video is
technically playing; and YouTube's compression and streaming systems, being
substantially more advanced than those I was capable of developing, will
always deliver better quality for the same bandwidth. Where my system
shines is in the ability to relocate the point-of-view videos while they
play, or to hide them entirely, letting students place them where they see
fit for a better viewing experience or clear them away to focus on the
panoramic view of the operating room.} Complaints of this shape are
repairable in principle: fix the named thing and the praise stands alone. Two
responses, though, name nothing. ``The virtual reality environment,
while it allows for more immersion, wasn't the highlight of the course,
and could be amended for something else'' (W26-R016/Q6). ``It is a cool
concept, but I feel like at the moment it is not doing much to help
regarding this class or the assignment'' (S25-R009/Q10). Nothing is
broken in these responses; something is missing, and the students
cannot say what. The sharpest case is a pair of answers by one student
---the same student whose response gave the survey its one confirmed
co-watching. Asked about educational use, this student praised the
platform in the survey's fullest presence language: it ``does put you
into a VR clinical environment,'' and ``you can watch with others''
(W26-R037/Q10). Asked what to improve, the same student wrote that the
platform ``is really laggy at times and is not significantly more
immersive then a YouTube video'' (W26-R037/Q9). One voice, both sides:
the world entered and shared, and the world flattened to a comparison
it barely survives.

What this register can testify to has a stated ceiling. In the
dissertation's account of boredom, which will be analyzed in the next
section, there is a form in which nothing determinate is
wrong---the occupation is intact, even pleasant, while the whole of it
quietly empties: an occupied surface over a hollow depth.\footnote{The
form is the second of the three Heidegger distinguishes in \textit{The
Fundamental Concepts of Metaphysics}: becoming bored \textit{with} an
occupation, as against being bored \textit{by} a determinate thing
(\textit{Fundamental Concepts} \S\S24--28). The full apparatus---the
three forms, their signatures, and what each requires of
evidence---is set out in the next section.} That form is detectable
only where two channels of evidence can disagree about the same
%
% >>> [STALE-1] REVISION QUEUED (see REVIEW NOTE at top of file).
% >>> Two clauses in this sentence and the next carry the SUPERSEDED
% >>> keystone: "physiological record vs behavioral one" and "a brain
% >>> measure against a gaze measure". Under the corrected layers both
% >>> are the SAME layer (surface); the depth is the retrospective
% >>> self-report. The structural point survives -- only the channel
% >>> names change. Do not revise until the lab section is drafted.
%
experience, as the laboratory's physiological record can disagree with
its behavioral one. A survey holds one channel: what a student
noticed, and chose to write, about their own experience. The
laboratory could set a brain measure against a gaze measure and let
the two contradict each other; a survey response has nothing set
against it that could contradict it. The two responses quoted above
that praise the platform while naming no culprit---W26-R016's and
S25-R009's---therefore show the same shape as this form of boredom:
commendation of an occupation that works, standing over an emptiness
the student cannot name. The shape matches; the survey cannot confirm
what lies beneath it. This section reads them as echoes of the form,
never as detections. What the survey can show directly is where the
arc of complaint ends. One further fact belongs to that ending: the
student who called the platform gimmicky went on, in the same
response, to propose its remedy---``perhaps the ability to meet up
with groups in a virtual reality setting'' (W26-R022/Q10)---the
platform's existing, unactivated capability. The survey's most
skeptical voice, answering a question about where virtual environments
could serve learning at all, proposes the one affordance this platform
offers that YouTube cannot: meeting inside the world, together. The
arc ends at a question no repair can answer: not whether the world
works, but what it is for---the justification the world itself must
supply.

% ========== END INLINED FILE: DESKTOP-M3-DRAFT-v1.tex ==========

% ===== M4 — ASSEMBLY-PROVISIONAL head (outline title de-stalled; user to rule)
\section{The Analysis: Boredom, Justification, and the WE}
% ========== BEGIN INLINED FILE: DESKTOP-M4-DRAFT-v1.tex ==========
% ============================================================
% DEFERRED REVISION — LAB-KEYSTONE DEPENDENCY (ruled 2026-07-29)
%   STATUS: queued, NOT to be executed until the laboratory boredom
%   section is drafted. Author ruling: working through the lab material
%   will likely revise the hypotheses and interpretations, so these
%   spans are revised ONCE, afterward, against settled findings —
%   not now, and not twice.
%
%   WHAT IS STALE: the rendered prose below cites the laboratory
%   keystone in its SUPERSEDED form — a brain measure reading
%   boring-like against a gaze measure reading engaged. The O-9
%   reprocessing (2026-07-16, N=8) withdrew it: P1's "clinical EEG
%   ~ boring" came from one clean subject and does not replicate
%   (clinical DMN is LOW = engaging-like; boring-like in 4/8 = chance).
%   The corrected keystone is ALL-PHYSIOLOGY-ENGAGED vs
%   SELF-REPORT-BORED (clinical self-report 6.0/9).
%
%   LAYER ASSIGNMENT (author-ruled 2026-07-29, governs the rewording):
%   physiology as a whole = the OCCUPIED SURFACE (outward comportment);
%   the post-video retrospective self-report = the HOLLOW DEPTH.
%   Grounded in fcm-05 SS24-28: at the dinner party we are "totally
%   involved" and DENY boredom, recognizing it only retrospectively;
%   "behavioural engagement is not evidence against boredom; in this
%   form it is the very medium of it."
%
%   SITES IN THIS FILE: see the per-site marker(s) below.
%   COMPANION SITES: DESKTOP-M3-DRAFT-v1.tex (M3.6) and
%   DESKTOP-M4-DRAFT-v1.tex (M4.1 P6); DESKTOP-SECTION-OVERLEAF-v1.tex
%   mirrors BOTH and is GENERATED — re-flatten after the edits, then
%   recompile DESKTOP-SECTION-ASSEMBLED-v1.tex (it \input-s M3/M4, so
%   the 37-pp PDF carries the stale claim until then).
%
%   CANONICAL SOURCES WHEN THE TIME COMES:
%     corpus/index/Heidegger - The Fundamental Concepts of Metaphysics/
%       _synthesis/fcm-king-salvo-bridge.md  (S4, revised 2026-07-29)
%     tmp/Dissertation/Part_III/reanalysis/boredom-o9-physiological-findings.md
%   NOTE: the structural point at both sites SURVIVES unchanged (a
%   survey holds one channel; the laboratory holds two). Only the
%   naming of the two channels is wrong.
% ============================================================
% ============================================================
% Part III — Desktop Section — M4 (Phenomenological Analysis: The Stalled
%   Chain and the Answered Question) — DRAFT v1
% Mode: Stage 3 guided drafting (per-paragraph protocol). M4.1 architecture
%   (7 ¶¶, 3 batches) approved 2026-07-27; batch 1 = ¶¶1–3 DRAFTED
%   2026-07-27 — IN REVIEW. Band G throughout; style canon = approved M3.1
%   sample (humanistic, slowed statistics, one claim per sentence, quotes
%   worked one at a time). Conventions as DESKTOP-M3-DRAFT-v1.tex (MLA;
%   respondent-code parentheticals; chain nodes A\textsubscript{n}; no
%   module labels in rendered text; no metatextual forward-refs; per-student
%   percentages w/ % symbol; no price/cost/ledger metaphors; student typos
%   verbatim, no [sic]; "YouTube platform/route"; terminology lock
%   [medium-supply vs phantasia gap-filling only]; "being-in" hyphenated).
%
% SOLICITATION RULING (2026-07-27, user-confirmed — gate CLOSED; full text
%   in HANDOFF-M4-GATE-PRESENTED-2026-07-25.md §5): solicitation in the
%   course deployment splits by agent and moment. (1) The recruiting hail
%   is INSTITUTIONAL and pre-perceptual — the professor's stated pitch at
%   the course's outset, in language, before the world renders; its content
%   = the VLE-only affordance cluster (co-watching the assigned procedure's
%   footage in shared virtual space + communication, toward group analysis).
%   (2) The world's own solicitation work is SUSTAINING — holding attention,
%   charging the clinical image, keeping use alive under mandate — and goes
%   silent only where soliciting would have to become TASKING (A3→A4).
%   Part II untouched (A1 = the Sonny event's own stationing: predatory-
%   solicitation mimicry atop the generic leisure entry-motive). The
%   outline's "solicited at A2/A3" = compression, unfolded in ¶¶1–2.
%   Leergelassenheit gloss CORRECTED: the TASK is withheld, not
%   solicitation; the emptiness is REGIONAL, stationed at the last link.
%   Hingehaltenheit sharpened: mandated duration = entry-under-mandate-
%   without-leisure-motive from inside. Drafting precision: the pitch
%   promised a place FOR the (course-side) task, not a task inside the
%   world — why the silence could hide under a successful pitch.
%
% USER OBSERVATION FOR LAB BOREDOM SECTION (2026-07-31, batch 5, at the
%   third-form ¶): the user believes they have an example of a TRANSITION
%   from FIRST form to THIRD form — a student who begins with high-arousal
%   boredom, fighting it, then surrenders and becomes completely docile
%   once he accepted (consciously or not) that he could not escape the
%   boredom. Carry into the lab section (cf. P2's high-arousal searching
%   gaze vs low-arousal "zombie stare" distinction).
% AUTHOR THEORY CAUTION (2026-07-27): boredom is ROUTE-INDIFFERENT at the
%   content level — footage that bores, bores on YouTube and in the VLE
%   alike; the routes affect how fast and to what degree, not whether.
%   Boredom therefore cannot ground the route-differential (that work
%   belongs to the WE, M4.3); its diagnostic value here = the watching-only
%   situation both routes stage. The WE-reduces-boredom effect is
%   SPECULATIVE ONLY (no reliable data) — if voiced at all (M4.3/M5
%   territory), it must be marked as speculation. User 2026-07-27: full
%   three-forms exposition PROCEEDS as planned (¶¶4–5); removable or
%   editable at review.
%
% INHERITANCE (build on, never re-derive): M3.3 ¶¶4–5 (A3 route-selection;
%   seven-causes individuation) · M3.4 (incorporation-debt; matched-numbers
%   insight) · M3.5 (floors 1/11/4/13, awareness-gap 12; coordination-
%   demand theorizing deferred to M4.3) · M3.6 (promissory footnote: the
%   three-forms exposition due in this section — to be paid at ¶¶4–5).
%   ¶2's "three of every five" inherits M3.2's D2 per-student stat (147/241
%   = 61%); presence-from-those-who-entered inherits F7. W26-R037 fragments
%   re-quoted deliberately for continuity (standing preference: make
%   same-student threads loud), char-exact verified 2026-07-24 per the
%   M3.6 ledger ("then" typo verbatim).
%
% DECISION POINTS RESOLVED 2026-07-27: pitch = author-confirmed course
%   fact, first-person plural ("we presented"); published trace footnoted =
%   ASEE 2025 "Scalable Virtual Reality" (King, Diec & Salvo; cite by PDF
%   page — conference paper, no journal pagination; vle-05 L1). VERIFIED
%   this session: span archon exact-1.00+bbox (vle-05 ledger) AND pdftotext
%   page-pinned to PDF p. 4 ("supported real-time voice chat and screen
%   sharing, enabling collaborative team viewing if students wanted to meet
%   in groups within the virtual environment"). OPEN AT REVIEW: the pitch's
%   delivery vehicle (syllabus vs first-session announcement; whether each
%   term) not yet author-specified — ¶1 kept to "from the outset"; confirm
%   or tighten wording. EXPANSION CANDIDATE held in reserve: vle-05 L4 —
%   the authors' own published boredom line ("engagement alerts will be
%   integrated to address student feedback on boredom," PDF pp. 9–10) —
%   fits ¶2 or ¶7 if wanted. WORKS CITED ADD QUEUED: King, Christine E.,
%   Kadin Diec, and Dalton Salvo. "Scalable Virtual Reality Global Clinical
%   Immersion for Culturally Responsive Engineering Design Skills." ASEE
%   Annual Conference & Exposition, Paper ID #49715, ASEE, 2025.
% Batch-2 verification ledger (2026-07-27, index-entry-first then PDF —
%   corpus/metaphysics/Heidegger…Fundamental Concepts…pdf; OCR text layer
%   w/ minor artifacts [stray pipes]; entry pdf-page loci run 1–2 HIGH in
%   the early chapters — true pdf pages pinned below; BOOK pages cited in
%   prose per the entry's triple loci):
%   "that which holds us in limbo and yet leaves us empty" p.87 (pdf 55) ✓
%   "the tasteless station of some lonely minor railway" + "It is four
%     hours until the next train arrives" + "catch ourselves looking at
%     our watch yet again" p.93 (pdf 58) ✓
%   "helpless gesture" p.97 (pdf 60) ✓
%   "There is nothing at all to be found that might have been boring
%     [about this evening]" + "I was bored after all this evening" p.110
%     (pdf 66) ✓
%   "what is boring does not come from outside: it arises from out of
%     Dasein itself" p.128 (pdf 76) ✓
%   "it is boring for one" + "'it is boring for one' to walk through the
%     streets of a large city on a Sunday afternoon" p.135 (pdf 79) ✓
%   "compelled to listen" p.136 (pdf 80) ✓
%   "the while becomes long" p.152 (pdf 88) ✓
%   Leergelassenheit / Hingehaltenheit / Zeitvertreib ✓ (translators'
%     brackets + glossary).
%   GERMAN FORM-NAMES (Gelangweiltwerden von etwas / Sichlangweilen bei
%     etwas / es ist einem langweilig): NOT PRINTED in the McNeill–Walker
%     translation (GA-side apparatus only) → per M3.6 precedent the prose
%     uses ENGLISH form-names + the translation's own "it is boring for
%     one." WORKS CITED ADD QUEUED: Heidegger, The Fundamental Concepts
%     of Metaphysics: World, Finitude, Solitude, trans. McNeill & Walker,
%     Indiana UP, 1995.
% Batch-2 v2 RULINGS APPLIED (user review 2026-07-27, "P4/P5"):
%   (1) opening reworded per user: "one of the most sustained,
%       philosophical, and structural analyses it has received";
%   (2)–(5) + P5: exposition RESTRUCTURED ¶¶4–5 → ¶¶4–7, one ¶ per form,
%       each with direct definition + Heidegger scene EXPLAINED (not
%       merely quoted) + demonstrative example + evidence requirement
%       worked against the record. "Channel" defined plainly up front
%       (¶4). "Named culprits/frictions/exits" unpacked as the form's
%       anatomy — culprits (quality/length/lag), occasions (crash/
%       buffering/wayfinding), counter-move = the flight to the same
%       footage elsewhere — w/ verified quotes. Second form: structure
%       spelled out + streaming-evening example + lab two-channel
%       demonstration (reference only) + record's echo-shape w/ codes.
%       Third form: Sunday-afternoon scene EXPLAINED (works because
%       nothing in it is boring) + student-version example + "powerless"/
%       "compelled to listen" + stillness/felt-time evidence + no-instance
%       disclaimer. MODULE NOW 9 ¶¶ (¶¶8–9 = application + diagnosis).
%   Batch-2 v2 additional verifications (2026-07-27): "all passing the
%   time is powerless against this boredom" p.135 (pdf 79) ✓ ·
%   W25-R017/Q9 "the youtube format ended up being simpler to use and
%   watch" char-exact vs open_text.csv ("youtube" lowercase verbatim) ✓ ·
%   W25-R014/Q9 re-confirmed char-exact vs open_text.csv ("dont" typo
%   verbatim) ✓.
% ============================================================
% --- M4.1 Completion, not bandwidth (¶¶1–7 APPROVED 2026-07-27 (¶¶4–7
%     approved w/ two ordered changes, both APPLIED same day): (a) RULE
%     RE-AFFIRMED — never open a sentence with "And" (three violations
%     fixed: ¶3 [approved-text touch, wording only], ¶4, ¶5); (b) P6
%     ruling — W26-R016/Q6 + S25-R009/Q10 quoted IN FOOTNOTE at the
%     culprit-less-responses clause w/ insight explanation (verbatims per
%     2026-07-24 ledger). ¶4 = apparatus + channel term; ¶5 = first form;
%     ¶6 = second form; ¶7 = third form (FCM §§19–23/24–28/29–36;
%     vocabulary from the FCM entry + crosswalk §1/§4, NOT Part II v4;
%     English form-names per verification finding). ¶¶8–9 APPROVED
%     2026-07-27 — **M4.1 COMPLETE (9 ¶¶, all approved)**. ¶8 = the two
%     moments at the deployment's
%     address (Leergelassenheit/Hingehaltenheit named; at-hand emptiness;
%     station-refusal analogy → hospital-refuses-itself-as-clinic;
%     S25-R035/Q6 orekton-withheld continuity quote; mandate-held close);
%     ¶9 = mis-attunement diagnosis (learner/ecology/task) + energeia
%     atelēs + regional emptiness + route-indifference guard + hinge to
%     the for-the-sake-of-which (M4.2's demand, unnamed). MODULE FULLY
%     DRAFTED pending ¶¶8–9 review.
%     Batch-3 verification ledger (2026-07-27): "To leave empty means to
%     be something at hand that offers nothing" p.103 (pdf 63) ✓ · "leave
%     us empty precisely because they are at hand" p.102 (pdf 63) ✓ ·
%     "The station at hand refuses itself to us as a station and leaves
%     us empty because the train that belongs to it has not yet arrived"
%     p.103 (pdf 63) ✓ · S25-R035/Q6 char-exact vs open_text.csv ✓
%     (continuity re-quote; also in M3.2 ledger) · energeia atelēs =
%     Part I canonical (spelling per §1.5 v3; gloss "incomplete
%     actuality, actuality whose end lies beyond it" reused verbatim from
%     §1.1 v3; locus Physics III.1, 201a10 per §1.1 fn; Metaph. IX.6
%     1048b30–36 available via §1.5). PHILOSOPHICAL GUARD: energeia
%     atelēs attached to interest-underway-toward-doing, NEVER to
%     watching/seeing as such (seeing = Aristotle's own example of
%     COMPLETE energeia, Metaph. IX.6 — a strict examiner would catch
%     the conflation).)

A world built for recreation never has to say why it should be entered.
Its player arrives already seeking something---an evening's diversion, a
story, a game worth playing---and the world reserves its arts of
solicitation for the encounters the player did not seek. A course world
begins from the opposite condition. The student arrives because the
course requires it, and so the solicitation that recruits them cannot
come from inside the world at all. It comes earlier, and it comes in
language. From the outset we presented the environment as the place
where a procedure group could watch its assigned footage together and
talk while watching---group analysis conducted inside a shared clinical
space, a way of working the course offered nowhere else.\footnote{The
platform's published description names the affordance directly: the
environment ``supported real-time voice chat and screen sharing,
enabling collaborative team viewing if students wanted to meet in
groups within the virtual environment'' (``Scalable Virtual
Reality'' 4).} That
announcement is the deployment's act of solicitation. It names the
reason to enter this world rather than the YouTube platform that
carries the same footage. What remains for the rendered world itself is
the quieter half of the work: not arresting attention but keeping it,
across the weeks of a term, in a world the student was assigned rather
than drawn to.

Judged by what it set out to sustain, that quieter campaign succeeded.
Three of every five students wrote something affirming the footage's
interest, the largest positive register in the collected data;
presence language came from those who entered the world; attention,
once gathered, was held. The deepest complaint in the data is not that
the world failed to draw anyone. The campaign goes silent at exactly one
point. Interest that has been brought to commitment wants something to
complete itself in, and completing it would require the world to stop
soliciting and begin tasking---to station an \textit{orekton}, a
determinate object for desire and action: something to find, perform,
or submit from inside the hospital itself. The world contains no such
station. There is no procedure to attempt in its rooms, no
needs-finding to capture where the footage plays, no deliverable the
world can receive. In the chain's terms, the deployed world reliably
carries a student to committed interest at A\textsubscript{3}; a
crossing to A\textsubscript{4} always follows, because the chain runs
always, and the question is which potentiality that crossing
actualizes. The only potentiality this world furnishes it is watching. The hail itself made this
silence hard to see. What we announced was a place for the group's
analysis, not an analysis the world would conduct; the task always
lived on the course's side, in deliverables the world never housed. A
pitch can succeed---students entered---while the world it recruited for
gives the interest it sustains nothing beyond watching to complete itself
in. A world that can only be watched has a boredom half-life.

None of this makes boredom a property of the footage. The same virtual
hospital drew the survey's fullest presence praise and its flattest
comparison, and its sharpest witness is a single student---the survey's
one confirmed co-watcher---who wrote in one answer that the platform
``does put you into a VR clinical environment'' where ``you can watch
with others'' (W26-R037/Q10), and in the next that it ``is not
significantly more immersive then a YouTube video'' (W26-R037/Q9).
Nothing about the hospital changed between those two answers; the
relation did. What holds for the footage holds for the route as well. A
student who tires of a long surgical recording will tire of it in the
virtual hospital and in a browser tab alike; I take the two routes to
differ in how soon that tiring arrives and in how far it goes, not in
whether it comes. Boredom, then, cannot be what separates the two
routes. Whatever it diagnoses, it diagnoses the watching itself---the
situation of being given something to see and nothing beyond watching to
do, which the deployment and the YouTube platform stage alike. Saying what that
boredom is, rather than merely that it occurred, requires distinctions
of structure rather than degree. In the dissertation's account, boredom
takes three forms; the forms differ in what they attach to, and they
differ just as sharply in what evidence can show them.

Heidegger's lecture course of 1929--30, published as \textit{The
Fundamental Concepts of Metaphysics}, gives boredom one of the most
sustained, philosophical, and structural analyses it has received, and
the analysis begins from a definition small enough to carry: what
bores is ``that which holds us in limbo and yet leaves us empty''
(\textit{Fundamental Concepts} 87). Boredom, on this account, is
always two things at once. Something holds us: we cannot simply leave,
because the situation keeps us in it. At the same time, what holds us
offers nothing, so that we are left empty while being held. Against this
double grip we deploy a counter-move Heidegger calls passing the time
(\textit{Zeitvertreib}): occupations taken up not for their own sake
but to keep the boredom off. The three forms he distinguishes are
three depths to which this machinery can reach. Setting them out
requires one term of method. Call a channel any independent way an
experience shows itself in evidence: what a person says or writes
about their experience is one channel; what their visible behavior
shows is a second; what an instrument reads off the body is a third.
The forms of boredom differ in how many channels it takes to see them,
and for the collected data that difference matters as much as the
forms themselves.

In the first and shallowest form, we become bored \textit{by}
something. A definite thing is the culprit---a book, a lecture, a
wait---and we could name it if asked, because the boredom stays
attached to its object. Heidegger's scene is a missed connection at
``the tasteless station of some lonely minor railway'': ``It is four
hours until the next train arrives'' (93). The wait holds us---we
cannot leave without losing the train---and the station offers
nothing, and so we fight, in the open: we read the timetable, count
trees, draw figures in the sand, and ``catch ourselves looking at our
watch yet again'' (93). The repeated glance at the watch is what
Heidegger calls a ``helpless gesture'' (97): it shows that the passing
of time is failing against the wait. Everything about this form is
visible, and that is its evidential mark. Because the culprit can be
named and the fight can be seen, a single channel suffices to attest
it. A student who writes ``[t]he VR is a cool idea but the quality is so
low I don't gain from the atmosphere and I dont feel immersed''
(W25-R014/Q9) has named the culprit inside the complaint itself; no
second instrument is needed to confirm what the sentence already
carries. The collected data holds this form's whole anatomy in
that one channel. The culprits are named: the footage's quality, its
length, the lag. The occasions are named: the crash at startup, the
buffering pause, the room that could not be found. The counter-move is
named too, because in this deployment the readiest way
to pass the time spent waiting on a struggling world was to leave it
for the same footage elsewhere: ``the youtube format ended up being
simpler to use and watch'' (W25-R017/Q9). Culprit, fight, and exit,
all written down in so many words.

The second form is stranger: we are bored \textit{with} something, and
no culprit can be found. Heidegger's scene is a dinner party. The food
is good, the talk is lively, nothing drags; asked at the table, we
would say, sincerely, that we are having a fine evening, and looking
back on it we can confirm that ``[t]here is nothing at all to be found
that might have been boring'' (110). Only afterward, at home, does the
recognition arrive: ``I was bored after all this evening'' (110).
Nothing was wrong with the evening, and that is the finding itself.
Where the first form's emptiness comes from a thing that refuses us,
here ``what is boring does not come from outside: it arises from out
of Dasein itself'' (128): the emptiness forms in us, underneath a
surface that remains genuinely, pleasantly occupied. The form has a
scene any student would recognize: an evening of episode after
episode, each one actually watched, none of them bad, which leaves one
hollowed rather than restored---nothing to blame, and yet the whole of
it empty. This structure---an occupied surface over a hollow depth---
is what makes the second form hard to evidence. The surface, asked,
would deny boredom and mean it; the depth cannot be asked. One channel
can therefore never show this form: a written answer reports the
surface, and the surface really is occupied. It becomes visible only
where two channels can disagree about one and the same experience, as
%
% >>> [STALE-2] REVISION QUEUED (see REVIEW NOTE at top of file).
% >>> "a brain measure can read boring-like on the very footage a gaze
% >>> measure reads as engaged" is the SUPERSEDED keystone -- it does
% >>> not replicate at N=8. Corrected: every physiological surface
% >>> reads the clinical footage ENGAGED while the retrospective
% >>> self-report rates it bored. The structural point survives.
% >>> Do not revise until the lab section is drafted.
%
in the laboratory arm of this study, where a brain measure can read
boring-like on the very footage a gaze measure reads as engaged. A
course survey holds one channel. It can catch, at most, the form's
echo---praise of an occupation that works, standing over an emptiness
the student cannot name, the shape of the survey's two culprit-less
responses\footnote{``The virtual reality environment, while it allows
for more immersion, wasn't the highlight of the course, and could be
amended for something else'' (W26-R016/Q6); ``It is a cool concept,
but I feel like at the moment it is not doing much to help regarding
this class or the assignment'' (S25-R009/Q10). The insight the pair
provides is structural. Each response commends the platform---the
first even grants it the immersion it promises---and each reports an
emptiness without naming anything broken. The hollowness attaches to
no defect but to the whole (``wasn't the highlight,'' ``not doing much
to help''), and in the second response it attaches to what the
environment is \textit{for}: its help toward the class and the
assignment. That is the second form's silhouette in survey prose: the
occupied surface speaks well of the occupation while the emptiness
shows through only as a want the student cannot locate. With one
channel it remains a silhouette; nothing in either response can
confirm what stands beneath it.}---and that is why the deployment's
divergence register was read as echoes, never as detections.

The third form is the rarest and deepest, and it strips away even the
person who is bored. Its name is impersonal: ``it is boring for one''
(135)---not for me, with my roles and plans and reasons, but for one,
for anyone at all. No thing is the culprit here and no occasion is
required, because what goes flat is not this or that offering but
beings as a whole: everything at once, oneself included, becomes
indifferent. Heidegger's example is chosen for its very lack of
incident: it is boring for one ``to walk through the streets of a
large city on a Sunday afternoon'' (135). The example works because
nothing in it is boring. The shops are shut, but one wanted nothing
from them; there is no wait, no missed train, no evening to see
through; the afternoon lacks nothing in particular---and precisely
there, with nothing particular refusing us, everything refuses at
once. A student's version is near to hand: a Sunday on which the
assignment, the shows, the group chat, the whole of one's own room
stand available, and none of it---not even oneself---holds any pull,
and looking for something to do feels pointless before the looking
starts. That last detail is the form's signature. In the first form we
fight the boredom; in the second we fail to notice it; here we do not
even attempt the counter-move, because ``all passing the time is
powerless against this boredom'' (135)---one is instead ``compelled to
listen'' to what it discloses (136). Its evidence is accordingly
stillness and time: a comportment that has stopped searching for
occupation, and a while that ``becomes long'' (152), felt duration
stretching far past the clock. A retrospective course survey can
capture neither, and this analysis claims no instance of the form.
Where anything of its shape appears in this study's evidence---in the
laboratory's stillest episodes and dilated felt-time reports---it is
read under an announced caution, as analog rather than instance.

Read with this apparatus, the collected data sorts almost
entirely into the first form---honestly so, since a survey is a
single channel and the first form is the one a single channel can
hold. What the apparatus adds is the structure beneath the sorted
complaints. Heidegger names the two moments that constitute every
form: being left empty (\textit{Leergelassenheit}) and being held in
limbo (\textit{Hingehaltenheit}). His most exact observation about the
first is that emptiness is not absence: ``To leave empty means to be
something at hand that offers nothing'' (103); things ``leave us empty
precisely because they are at hand'' (102). The virtual hospital is at
hand in every sense that matters---rendered, explorable, stocked with
footage, praised in the survey's fullest presence language---and
precisely as such it leaves empty, because what it puts at hand offers
everything to see and nothing to complete. The station's refusal even
has an address: ``The station at hand refuses itself to us as a
station and leaves us empty because the train that belongs to it has
not yet arrived'' (103). The virtual hospital refuses itself as a
clinic in exactly this way. A clinic is a place where procedures are
performed and needs are found; this one holds only their images, and
its train never arrives---no task pulls in to make the place what its
rooms announce. One student stated the resulting emptiness with
precision: the procedures are ``many hours long, and I'm not sure what
exactly to be looking for during that time'' (S25-R035/Q6)---at hand,
and offering nothing determinate to look \textit{for}. The second
moment completes the situation. At the station we are held by an
interval with a terminus: four hours, then the train. The deployment's
students are held by mandate, with no terminus inside the world
itself---assigned footage, required hours, a term's duration, and, in
the small, the buffering pauses that hold a viewer between intention
and image. The deployment thus realizes both of Heidegger's moments at
once: its students are held---by mandate, by assigned footage, by the
buffering pauses---before a world whose acts stop at watching. The two
moments meet at exactly one place, and the location is the evaluation:
everything upstream of that point works, so the emptiness students
report belongs to the single link the deployment left unbuilt.

The diagnosis follows, and it is a diagnosis of a situation, not of a
platform. Boredom, as the collected data presents it, marks a
mis-attunement among three terms: a learner who arrives under mandate,
with no leisure motive for the world to ride on; an ecology that does
its sustaining work, as the data shows it doing; and a task that has
no address inside the world. The interest the world sustains is not a
contemplative looking that completes itself in every moment of its
exercise. It is interest underway toward doing---the students'
requests for assignments, for
interaction, for something to find make its direction explicit---and
motion of that kind is, in Aristotle's terms, \textit{energeia
atelēs}: incomplete actuality, actuality whose end lies beyond it
(\textit{Physics} III.1, 201a10). A world that sustains the motion
while housing no terminus does not merely disappoint the interest; it
holds it, unfinished, for as long as the mandate runs. The
deployment's shortfall, then, was never bandwidth---presence survived
the monitor---but completion. This is the precise sense in which the
deployed world's emptiness is regional. Nothing upstream is withheld:
the world is entered, rendered, watched, even dwelt in; the
being-left-empty is stationed at the chain's last link, where
soliciting would have to become tasking and the world falls silent.
The boredom the data attests is therefore not a verdict on the
footage, and not a difference between the two routes; it is the felt
shape of the watching situation both routes stage, and, unlike a crash
or a buffer, it names nothing a repair could answer. The arc of the
data shows as much: the frictions fell term by term, and what rose in
their place was the question of what the world is for. An emptiness
stationed where action should begin asks for one thing only---a
for-the-sake-of-which, a reason this world, rather than a video, is
where the work of the course happens.

% --- M4.2 Veri(dis)similitude's second face (plan approved w/ P1 rulings
%     2026-07-27: (1) opening = plain restatement of the revised
%     veri(dis)similitude definition, NO stylistic metalanguage ("the
%     deployment's twist" etc. banned); (2) hypothesis notation (H1/H3)
%     NEVER in rendered text (extends the module-label rule); (3) MA-thesis
%     poles RESTATED not quoted (decision point resolved by ruling 1).
%     4-¶ architecture: ¶1 condition restated + turned inward · ¶2 cohort
%     re-derives it + the un-redeemed answer · ¶3 materials audit
%     (commodity features → YouTube; place strongest form, insufficient
%     alone) · ¶4 the material video cannot approximate (hands M4.3 its
%     premise; de-facto scoping stays in M4.3). ¶¶1–2 DRAFTED 2026-07-27 —
%     in review. Verbatims reused from verified ledger: W26-R033/Q9;
%     W26-R022/Q10 ("the least gimmicky" + meet-up proposal, same
%     response per M3.6). Part II two-pole sentence ("too familiar cannot
%     move / too strange cannot be entered") reused as internal formula,
%     no cite. Nine-campus facts = ¶3 (M0.2 + author-confirmed memory).
%     ¶¶1–2 v2 RULINGS APPLIED 2026-07-27: ¶1 close EXPANDED per user
%     (subtract-the-footage comparison of the two remainders + the
%     film/theater demonstrative example); ¶2 opening = user formula
%     ("As the collected data demonstrates"); "RECORD" BAN RE-AFFIRMED
%     (standing framing rule) → GLOBAL SWEEP executed across M4.1
%     approved text + M4.2 (16 prose instances → collected data/data/
%     survey/students' language; wording-only touches to approved ¶¶,
%     And-precedent). NB the flagged "remove first sentence" items were
%     chat-presentation labels, never in this file; future batch
%     presentations carry bare ¶-numbers only.
%     ¶¶1–2 APPROVED 2026-07-27 w/ one ordered change APPLIED: the term
%     veri(dis)similitude stated DIRECTLY (named at the definition, ¶1
%     opener, verisimilar/dissimilar poles matched to the term; and at
%     the turn "Veri(dis)similitude accordingly takes a second and
%     sharper form"). ¶¶3–4 DRAFTED 2026-07-27 — in review (materials
%     audit + the unapproximable material; W25-R014/Q6 quality-trade
%     quote char-exact vs open_text.csv, inner ""immersion"" → MLA
%     ¶4 v2 APPROVED 2026-07-27 — **M4.2 COMPLETE (4 ¶¶, all approved)**.
%     single quotes, extends the W25-R014 thread [Q9 in ¶5 of M4.1];
%     ASEE fallback span "with interactive 360-degree panning" verified
%     char-exact via pdftotext at PDF p. 4, cite ("Scalable Virtual
%     Reality" 4); nine-campus facts per M0.2 + author-confirmed memory;
%     hands M4.3 its premise, de-facto scoping reserved to M4.3.
%     ¶3 APPROVED 2026-07-27. ¶4 v2 per ruling 2026-07-27: opening two
%     sentences REWRITTEN slow+concrete ("rival route"/"rides either
%     route" compression replaced by worked examples — course-side
%     assignment [worksheet/quiz/unmet-needs prompt] works under either
%     delivery because administered outside the footage; in-world
%     needs-task = unbuilt possibility; browser panning = drag-the-view,
%     lacks corridor/doorway/position); "approximation" → "imitation"
%     diction unified (one-word touch to the approved close). IN REVIEW.)

A rendered world, if it is to do anything to the person who enters it,
must satisfy a double condition: \textit{veri(dis)similitude}. It must
be verisimilar enough---familiar enough in its look, its workings, its
ways of being moved through---that a person can enter it and dwell in
it at all; and it must be dissimilar enough---different enough from
what everyday life already offers---that the entering changes
something. A world too familiar struggles to move; a world too strange
opposes being entered. For the deployed virtual hospital, the second half
of this condition is measured against an unusual standard. Every
procedure the hospital holds also plays on the YouTube platform, one
click away, free of installers, sign-ins, and crashes. A student
deciding whether to enter the world is therefore not weighing it
against everyday reality; they are weighing it against the video of
itself. Veri(dis)similitude accordingly takes a second and sharper form.
Ordinarily a rendered world earns its dissimilarity against everyday
life, out of possibilities everyday life does not hold. The deployed hospital
cannot earn its difference there, because its substance---the
procedures---exists identically outside it. Subtract the footage from
both routes and compare what remains. On the YouTube platform the
remainder is a familiar player and its conveniences: the progress bar,
the subtitles, the speed control. In the virtual hospital the
remainder is everything the footage plays inside of: the walk down the
corridor, the room found and entered, the screen on its wall, and the
possibility of another student standing at the same screen. The world
justifies its existence with that remainder or not at all. The demand
is not unique to virtual worlds. A film that can be streamed at home
identically must give a viewer a reason to go to the theater, and the
theater's answer has never been the film, which is the same in both
places; the answer is the hall, the ecology, the environment, the occasion,
the others in the seats, and the ambience. So with the deployed world: to justify its existence it must be
different enough from the video of itself, and the difference can be
made only of what surrounds the footage---never of the footage, which
both routes share.

As the collected data demonstrates, the cohort arrived at this same
condition from inside it. As the platform stabilized and the
complaints about function fell away, the talk that rose in their place
questioned existence rather than function: ``I think the youtube
videos were enough to fully grasp the experience'' (W26-R033/Q9). A
question of that kind, put to a world whose footage plays elsewhere,
is the difference-condition stated in a student's words. One student
stated not only the question but the condition itself, in full: ``There
is no difference between virtually moving through a museum and seeing
photos with text underneath for the purposes of learning''; a virtual
environment earns its place only if it ``can provide an experience that
is unable to be achieved without it, and allows for learning that is done
at one's own pace and with elements that rely on the immersion potential
and interactivity''---otherwise ``there is no purpose in using it over
other alternatives'' (W26-R059/Q10). This is the student the sign-in gate
stopped at the threshold; the voice the door turned away is also the
survey's clearest theorist of what would justify walking through it. The
question was not unanswered. It had been answered before the first term began,
in the announcement that recruited the class: the difference was
company---a procedure group watching together, inside the world,
talking while it watched. What the data then shows is that
difference going largely unconvened. Students watched alone, found
what a world minus its difference amounts to, and asked why it
existed. The survey's most skeptical voice---the student who judged
traditional learning ``the least gimmicky'' way---proposed, in the
same response, ``perhaps the ability to meet up with groups in a
virtual reality setting'' (W26-R022/Q10): the difference itself,
proposed back to the course that had announced it. The question the
deployment ends on is an answer it began with, delivered and never
redeemed.

What could the difference be made of? Not the commodity features of
delivery, which the students' own comparisons settle. The YouTube
platform streams the footage at higher quality for the same bandwidth,
carries subtitles, and changes playback speed at will; one student
weighed the trade directly, having ``gained more from the youtube
videos because I could watch in much higher quality which made up for
the slight reduction in `immersion'{}'' (W25-R014/Q6). On every
feature a video player can have, the player built for videos wins.
Three materials remain that no player has: a task, a place, and
company. A task the deployed world did not contain; that absence is
the diagnosis already made. Place it did contain, and the place
deserves its full weight. This class has never had a room. Its
students enroll across nine university campuses whose calendars do not
even align; there is no lecture hall, no laboratory section, no
building, no campus they share. The virtual hospital is the only room
the class has ever had---not a supplement to a classroom but the sole
site at which this cohort could be co-present in any sense at all. The
place also did quiet cognitive work for at least one student: ``I found
the online course format very useful--especially the VR classroom space.
I thought it was really cool to navigate, and helped compartmentalize
individual findings'' (S25-R048/Q5)---the building's rooms functioning as
an organizer for the student's own analysis, a use no video list affords.
Yet the collected data shows place alone not sufficing. Presence was had
on the desktop route and attested plainly, and the flight to the
YouTube platform happened anyway; a place entered alone, with nothing
to do in it, reduces in use toward the very thing it competes
with---a scene to survey, reached the harder way.

Two of the three materials are not wholly beyond the YouTube
platform's reach. Consider the task first, in the course's own
terms.\footnote{Course structure and graded instruments per the course
syllabus, BME 179: Biomedical Engineering Design: Addressing Unmet
Clinical Needs (University of California, Irvine); the same syllabus
has been used across the offerings surveyed.} The syllabus grades
three instruments: weekly check-in quizzes keyed to each week's
lecture, an individual innovation notebook in which unmet needs are
recorded and evaluated, and a team final report. The weekly schedule
sends every group into the virtual environment---making a login and
beginning to explore in the second week, then working through the
procedure rooms and physician interviews in the weeks that
follow---but each graded deliverable is completed and submitted
outside the world, on the course page. A check-in quiz answers to the
lecture, not the footage; the innovation notebook is written wherever
the student writes; the final report is drafted with the team, in a
document. Work of that kind runs as well under a YouTube link as
inside the virtual hospital, because the course administers it and the
student completes it outside the footage---no part of it needs the
world. A world could go further: it could make
the finding of needs something done in its rooms, captured where the
footage plays. The deployment built no such task, and until one is
built, the task-material does not separate the two ways of watching.
The place can be partly imitated as well. The footage itself is
360-degree film, and the deployment's own fallback delivered the same
videos ``with interactive 360-degree panning'' (``Scalable Virtual
Reality'' 4): a student in a browser can drag the view around the
operating room and look wherever they choose. What the browser cannot
give is everything around that view---the corridor, the doorway, a
position in the room---so the pannable video is a thin, surveyable
shadow of the place; but it is a shadow of it. The meeting admits no
imitation at all. Watching together in a shared rendered room,
talking while the procedure plays, is not a feature a video player
lacks in degree; it is a kind of thing no video player has. For most
groups of students in this class there is no library to meet at
instead, no dorm lounge, no lab bench; the alternative to meeting in
the virtual hospital is, as one student noted, social media platforms
like Discord or not meeting at all.\footnote{``The virtual space to
meet up for teamwork was very smart, but my group utilized discord a
lot'' (S25-R072/Q10). The course itself maintained a Discord channel
for the whole class in every term surveyed, so the nearest alternative
meeting place carried institutional standing of its own.} If the world's difference from the video of itself is to be
made of anything, it must be made of the meeting.

% --- M4.3 The WE as the desktop's ontological differentiator (★ SPOTLIGHT.
%     Plan approved 2026-07-27 incl. the ONE-SENTENCE marked-speculative
%     WE–boredom acknowledgement in ¶5. 5-¶ architecture: ¶1 dual-being +
%     toward-whom ground · ¶2 route audit + de-facto scope · ¶3 being-with
%     mode + honesty clause · ¶4 deferred theorizing + floors · ¶5 climax.
%     ¶¶1–2 DRAFTED 2026-07-27 — in review.
%     Batch-1 verifications: "multiplayer functionality, voice chat, and
%     customizable avatars" = AUTHORS' OWN prose ("our learning
%     environments"), pdftotext-pinned PDF p. 10 = journal p. 390 →
%     first-person rendering, cite ("Phenomenological Evaluation" 390);
%     p. 394 multiplayer hit = appendix instrument text, NOT used.
%     Toward-whom = Part I canonical (toward-which/toward-whom intentional
%     structure; thymotic action "always in relation to another"; Mitsein
%     slot) — RESTATED without cite per no-re-derivation + ambiguous
%     Bekker ranges across Part I files (1378a31–b5 vs a30–b9; NB
%     FourTypes Type-2 sentence carries a PROVISIONAL marker — nothing
%     quoted from it). Burke NAMED for consubstantiation, no quote —
%     ADD A Rhetoric of Motives to Works Cited at assembly.
%     M1.6 inheritance (cite forms per its locked header): Biocca et al.;
%     Terry and Doolittle — argued use lands in ¶3.
%     ¶¶3–5 APPROVED 2026-07-27 — **M4.3 COMPLETE (5 ¶¶, all approved)**.
%     Batch-2 verifications:
%     Biocca "the sense of being with another" (456) + Terry & Doolittle
%     relocation (121, paraphrase-cite) + Kreijns Pitfall 1 (336,
%     paraphrase-cite; full deployment reserved for M5.2) = M1 bank
%     PDF-verified reuses, cite forms per M1.6; W25-R038/Q10 "when I did
%     log in, I was usually alone in the platform...connected with my
%     peers" char-exact vs open_text.csv this session; W26-R037/Q10
%     fragment = ledger reuse (continuity). Floors per M3.5 locked spec
%     (1/11/4/13, awareness 12) — counts only, NO rates. ¶5 carries the
%     approved ONE-SENTENCE marked-speculative WE–boredom acknowledgement
%     + answers M3.4's open question (exclusive binding = the others) +
%     un-redeemed-hail close ("temporal economy" = outline-approved
%     phrasing).)

Rhetorical incorporation is one in substance and dual in being. A
world becomes a student's own in two ways at once: as a place they
dwell in---world-disclosedness, the register every presence finding so
far has measured---and as a place that is theirs together with
others---co-disclosedness, the same world had as ours. The two are not
a feature and its supplement; they are two aspects of one
phenomenon, and the chain itself grounds the second. Desire's thymotic
species is irreducibly social: one is angry at someone, stands on
honour before someone; the emotions that concretize judgment into act
are articulated toward or before determinate others---anger toward the
slighter, shame before the audience that matters. An act, when it
finally comes, is rarely aimed at no one; the chain's last link
carries a toward-whom as part of its structure. What Burke names
consubstantiation---persons made, through acting together, of one
substance---names the same ground from rhetoric's side: to act with
another is to be with them in a way no shared spectacle alone
provides. A world that supports action will support it toward and with
others, or it will support less than action's full structure.

The three routes hold this dual structure very differently. The
published headset version maximized world-disclosedness: the encircling
display, the fullest supply of the operating room a medium could give.
Co-disclosedness it did not carry. Students entered alone, and we said so in
print, naming among the version's limits ``the sheer lack of features
found in traditional games like multiplayer functionality, voice chat,
and customizable avatars'' (``Phenomenological Evaluation'' 390). The
YouTube platform possesses neither in kind: it is a player, not a
world; there is nothing in it to be in, together or otherwise. The
desktop route inverts the headset's profile. Its world-disclosedness
rests on the thinnest supply: the monitor gives less, and \textit{phantasia} furnishes
more of the world. But as the program stands
deployed, it is the only route on which co-disclosedness is available
at all. The scope of that claim is historical, not metaphysical:
proximity voice chat could be built into a headset version---nothing in the
hardware forbids it---but development moved to the desktop build, and
the social possibilities were built only there. The affordance is not a
consolation native to screens; it is a developed asset this platform
alone carries.

What the meeting changes is not the platform's feature list but the
being of the world. A room entered alone is a viewing station; the
same room entered with another student is ours before anything is said
in it---a place jointly attended, where the judgment forming over the
footage can be spoken, answered, and made common. The research on
learning at a distance has a name for this register of experience and
a hard-won result about it. Social presence---defined at its core as
``the sense of being with another'' (Biocca et al.\ 456)---is not a
property a medium contains but a state persons achieve, something a
course can activate or let lie (Terry and Doolittle 121). Those
findings were made about text forums and video calls, not about
rendered worlds, but the insight is applicable all the same:
co-disclosedness, being-with-others is a second presence, of its own
kind, that the desktop hospital can host---and the first kind, presence
in a place, cannot substitute for it.

A world, though, can only hold a meeting open; it cannot convene one.
Footage is delivered to whoever presses play, and a place is there for
whoever walks in, but being-with requires that others arrive. Arrival
is exactly what the deployment's shape made unlikely. Convening a
group asks what solo watching never asks: schedules aligned across
campuses that share no calendar, an invitation risked on classmates
never met, proactive organization, and commitment to actually meet
with each other. What that commitment can come to in practice, one
student recounts in full: ``They reached out to me in order to begin the
group project, and I responded, and I was ghosted until the final few
weeks of the quarter where I found myself very behind in the group's work
with no knowledge that they were organizing without me'' (W26-R036/Q6).
The account concerns the group's work generally rather than the virtual
hospital, but it shows the social ground the meeting had to stand on: an
invitation extended, answered, and dropped. The collaborative-learning
literature warns of precisely this: interaction cannot be assumed
merely because an environment makes it possible (Kreijns et al.\ 336).
The collected data shows the warning borne out. In 873 responses there
is one confirmed convening---the student who wrote that ``you can
watch with others'' (W26-R037/Q10). Eleven students wrote of being
alone in the world; ``when I did log in, I was usually alone in the
platform,'' one of them explained, ``so it didn't really make me feel
more connected with my peers'' (W25-R038/Q10). Four asked the course
to require the meeting; thirteen asked for the capability, twelve of
them in wording implying it did not exist. The affordance that answers
the world's justification-question was unknown to part of the class
that wanted it, unconvened by nearly all of it, and lived by one voice
in three terms. The world held the WE open; the course's temporal
shape and lack of designated tasks let it lapse.

The deployment's deepest loss now has its exact name. It is not
presence---presence the desktop world had, measured and attested. It
is the WE left unconvened: a co-disclosedness the ecology held open
for three terms and was predominantly unfulfilled. The arc that
began this analysis closes on itself. The difference we announced at
the course's outset was the meeting; the audit of materials arrived
at the meeting; the survey's most skeptical voice, unprompted,
reinvented the meeting. The announcement was right, and what it named
went largely unlived. The question the presence analysis left
open---what could bind a world to a student exclusively, when its
footage plays anywhere---has the same answer: its others. The virtual
hospital is irreplaceable exactly insofar as it is shared, and
replaceable exactly insofar as it is not. Whether a convened WE would
also blunt the boredom this analysis diagnosed---company giving
sustained interest a completion of another kind---the collected data
cannot say; the suggestion is speculative, and it is offered as such.
What the data can say is narrower and firmer: the one thing this
world alone could be was the one thing it was almost never used as.

% --- M4.4 What the desktop case gives back to the frame (SOFT-BLOCK
%     WAIVED by user 2026-07-27 — the Part II closing differential is not
%     yet written. ***REVISION-PASS REQUIRED: revisit this sub-module
%     once the Part II closing differential exists; harmonize the
%     empirical register here against the differential's final claims
%     (division of labor AUTHOR-CONFIRMED 2026-07-16: Part II = media-
%     type/interaction-mechanics/design differential incl. phantasia's
%     near-leveling; Part III = diagnosis of the experienced difference).***
%     2 ¶¶ per outline: frame-return (matched numbers 70%/70.7% INHERITED
%     from M3.4, not re-derived; phantasia-leveling echo AUTHOR-APPROVED
%     2026-07-16; "Route sets the load; the task sets the dwelling" =
%     cross-case §2(vii) motto, first rendered appearance) + one
%     gesture-scale ¶ on what only a headset deployment could test
%     (gesture only per standing rule). Terminology lock held ("phantasia
%     furnishes more" — no image-making/phantasia-load). ¶¶1–2 APPROVED
%     2026-07-28 — **M4.4 COMPLETE. M4 MODULE COMPLETE (M4.1 9¶¶ ·
%     M4.2 4¶¶ · M4.3 5¶¶ · M4.4 2¶¶ — all approved). M4.4 revision-pass
%     pending Part II closing differential (see marker above).**)

Two findings from this deployment reach beyond it. The first is the
pair of matched numbers: the pilot's 70\% under headsets, the
deployment's 70.7\% before ordinary monitors. The supply could hardly
differ more---an encircling display against a screen---yet
the felt result barely moved. Parity of that kind is what the leveling
power of \textit{phantasia} looks like in data: where the medium
supplies less, \textit{phantasia} furnishes more, and the felt world
arrives nearly whole either way. The second finding is what no route
could level: the act that none of them contained, and the meeting that
only one of them carried. Route sets the load; the task sets the
dwelling. The first half of the motto names the chain's opening
stations: the route---headset, desktop, browser---decides how much of
the world reaches the student through the senses and how much
\textit{phantasia} must furnish from within (A\textsubscript{1} into
A\textsubscript{2}), and the matched numbers show \textit{phantasia}
equal to any load a route set it. The second half names the chain's
end. Dwelling is world-disclosedness sustained over time: not the
momentary feeling of being placed in a world, but the settled being-in
of a student who returns to the world, moves through it toward
something, and has their doings answered by it. What converts a place
merely felt into a place dwelt in is the act---the passage from
committed judgment (A\textsubscript{3}) into an actualized potentiality
(A\textsubscript{4}). That passage occurred on every route, because
watching is itself the actualization of a potentiality the world
furnishes; the deployment's students actualized it by the hour. The
design question is therefore never whether students will actualize a
potentiality---they always will---but which potentialities the world
makes available, which one each student actualizes, and why that one
rather than a richer one. At the scale of 241 students, the deployment
attests that this is where worlds are won or lost: at the passage from
A\textsubscript{3} to A\textsubscript{4}, in the worth of the
potentialities a world furnishes to the commitments it has solicited.

% [¶ DELETED by user order 2026-07-31 (batch 5, 5.20): the closing
%  headset-deployment gesture paragraph ("What this case could not
%  test...") removed in full. M4.4 is now 1 ¶.]

% ========== END INLINED FILE: DESKTOP-M4-DRAFT-v1.tex ==========

% ===== M5 — ASSEMBLY-PROVISIONAL head (outline's own title)
\section{Design: The Levers as \textit{Archai}}
% ========== BEGIN INLINED FILE: DESKTOP-M5-DRAFT-v1.tex ==========
% ============================================================
% Part III — Desktop Section — M5 (Design — The Levers as Archai) —
%   DRAFT v1
% Mode: Stage 3 guided drafting (per-paragraph protocol). Plan approved
%   2026-07-28. Band G; style canon = approved M3.1 sample; conventions
%   as DESKTOP-M3/M4 drafts (MLA; respondent-code parentheticals; chain
%   nodes A\textsubscript{n}; no module labels or hypothesis notation in
%   rendered text; no metatextual forward-refs; no stylistic
%   metalanguage; NEVER sentence-initial "And"; NEVER "record" as name
%   for evidence; % symbol per-student percentages; student typos
%   verbatim, no [sic]; "YouTube platform/route"; terminology lock;
%   "being-in" hyphenated).
%
% VOCABULARY RULING (user 2026-07-28): keep technical terms that do
%   analytical work — coordination game, equilibrium, activation
%   energy — REMOVE all cost/price/budget/spend metaphor language per
%   the standing ban; chronos = a DESIGN MATERIAL (allotted/reserved,
%   never spent/budgeted). User: "we can and will revise as we go."
%
% Architecture (approved): M5.1 (2 ¶¶ levers-as-archai) → M5.2 (5 ¶¶,
%   ★ the expanded WE lever: equilibrium · synchrony constraint ·
%   rung 0 + low ladder · ladder top w/ ignition logic + bia guard ·
%   motto) → M5.3 (1 ¶ motto/mirror w/ capture-economy breath).
% INHERITANCE: M4.3 (hold-open-vs-convene; floors); M3.5 (locked spec);
%   M0.2 (systemwide asynchrony); M1.6 locked bank (Kreijns Pitfall 2
%   full deployment HERE; De Back inverse group-size; Garrison binding-
%   element; van der Meer verbal-audio-essential); M3.3 ¶5 (bia/hexis
%   guard). Anchors: cross-case §3.3; deployment-brief §5; F7 correction
%   + recount.
% ============================================================
% --- M5.1 The levers, ranked, as archai (¶¶1–2 DRAFTED 2026-07-28 — in
%     review. ARCHĒ LOCUS PINNED: Part I §1.1-v3 footnote — six senses at
%     Metaphysics V.1 (Δ.1) 1012b34–1013a23, chain uses the FOURTH
%     ("that from which...a thing first arises, and from which the
%     movement or the change naturally first proceeds," 1013a7–10, Ross
%     wording as quoted in Part I; outline's "Part I §3" pointer was
%     off — it is §1.1). Lever content per deployment-brief §5 +
%     cross-case §3.3; threshold "quarter of the class" inherits M3.1;
%     orekton lever inherits S25-R035 thread (no re-quote); pacing lever
%     kept AFTERMATH-FREE (depletion/carryover = lab-pole evidence,
%     reference only — phrased as design-outcome claim instead).
%     EXPANSION CANDIDATE held: ASEE published planned fixes
%     (MacOS/Linux, reduced storage, mobile — vle-05 §5) could cite into
%     the threshold lever; page unpinned, add on user order.
%     P2 RULINGS APPLIED 2026-07-28 — ¶2 REWRITTEN as ¶¶2–3:
%     (1) plain-named lever list up front: THE DOOR (grounded in M3.1's
%     own "only door left open" image) · THE TASK (chain A4) · THE FIND
%     (needs-finding + orekton) · THE MEETING (M4.3's own term) · THE
%     WHILE (FCM Weile — "the same while that boredom stretches");
%     (2) "threshold" dropped for door language, made concrete;
%     (3) task lever slowed: outside-completable assignments → in-world
%     work → reason-to-enter → justification answer;
%     (4) find lever recast as OBJECT OF DESIRE (not of the look):
%     mandate-context curation — name what a needs-finder watches for,
%     caption/transcribe, deliverable-linked finds, finds that reward a
%     second pair of eyes (desirable object doubles as occasion to meet
%     → feeds the meeting lever);
%     (5) "co-pole" removed — POLE BAN surfaced and recorded
%     ([[feedback-no-pole-language]]); meeting lever grounded in
%     co-disclosedness. RETROACTIVE SWEEP of M4.3 executed same day
%     (6 instances: register/two sides/world-disclosedness maximized/
%     "Co-disclosedness it did not carry"/neither in kind/rests on the
%     thinnest supply). OPEN FLAG: M3.6 approved prose has ONE "both
%     poles" in the divergence-pair sense ("One voice, both poles") —
%     user ruling needed before touching the M3 file.
%     P2/P3 v2 RULINGS APPLIED 2026-07-28 — ¶¶2–3 REWRITTEN as ¶¶2–4
%     (M5.1 now 4 ¶¶): (P2.1) lever names ITALICIZED throughout;
%     (P2.2) third-term question FOOTNOTED w/ W26-R033/Q9 + W26-R056/Q9
%     (ledger reuses); (P2.3) TASK-LEVER DOCTRINE CORRECTED (AUTHOR):
%     in-world requirement = compulsion = "pedagogically prosaic," never
%     the justification; the world earns existence by being BETTER —
%     rhetorical incorporation as the mechanism of pedagogical
% TASK-LEVER QUOTATION PASS 2026-07-28 (user-approved set; Makransky
%   genetics HELD OUT by user ruling). Backup 20260728T155206.
%   Insertion 1 (¶3): Dalgarno & Lee 2010 [VoR year] — affordance
%     sentence + trivial-fact warning, BOTH PDF-verified char-exact,
%     TRUE p. 25 (entry loci ran 1 high; form-feed count = 23 pp =
%     10–32 exact); W26-R033/Q9 fragment re-verified vs CSV line 787.
%   Insertion 2 (new ¶ after ¶4 close): Cook et al. 2011 JAMA meta —
%     609 studies / 35,226 trainees / "question the need" quote
%     PDF-verified TRUE p. 987; scoped honestly (vs.-no-intervention,
%     medium-agnostic, "without ranking routes"). Dubovi, Levy & Dagan
%     2017 — PILL-VR desktop clinical sim; gains vs lecture held at
%     5-month retest (23); 9/10-vs-9/10 dissociation SCOPED to n=10
%     recorded subsample (23–24); control + presence-not-unique quotes
%     PDF-verified char-exact TRUE p. 25 (entry said 26 — 1 high;
%     form-feed count = 12 pp = 16–27 exact).
%   Works Cited adds queued: Dalgarno & Lee 2010 (short title iff 2012
%     followup also enters), Cook et al. 2011, Dubovi/Levy/Dagan 2017.
%   M5.1 = 7 ¶¶ after pass. APPROVED 2026-07-28 — PASS COMPLETE.
%     superiority. FUTURE ITEM BANKED: expanded task-lever discussion w/
%     secondary pedagogical literature on proven interactive-learning-
%     environment efficacy — placement TBD (raise at M5.3). (P3.1)
%     FIND-LEVER DOCTRINE CORRECTED (AUTHOR): mandated student DOES
%     bring a governing desire — pass → degree → life = the
%     FOR-THE-SAKE-OF-WHICH (Heidegger tie rendered); routes vary
%     beneath it; YouTube-gravitation rational under route-parity;
%     design aim = ENGINEERED AGENCY — VLE better past parity, YouTube
%     stays a real option (quick refresh, phone), free choice keeps
%     landing on the world; Part II curation formula reused as internal
%     echo ("handed a world, never a word"; author's "hack their soul"
%     doctrine — soul-line reserved for M5.3's ethics beat). (P3.2)
%     second-pair-of-eyes EXPLAINED (multiple activity streams, assigned
%     vantage points, reconcile-table-with-room prompts) — moved to the
%     meeting portion as the find×meeting bridge. (P3.3) WHILE-LEVER
%     CORRECTED (AUTHOR): procedures must NOT be shortened — edits
%     misrepresent the operation and can delete the very corners where
%     unmet needs sit; levers = SEGMENTATION (weekly assigned portions +
%     follow-up) and/or INTERVAL INTERACTIVITY (prepared sequences that
%     interrupt the drag and redirect attention); "the while is not
%     shortened; it is articulated."
%     v4 RULINGS APPLIED 2026-07-28 (P2/P3 second round) — M5.1 now
%     6 ¶¶ (¶2 list+door · ¶3 task-obstacle · ¶4 task-on-the-chain ·
%     ¶5 find · ¶6 meeting+while [¶6 = prior ¶4, no rulings, carried]):
%     (P2.1) footnote quote REPLACED w/ unused W25-R035/Q10 (student-
%     designed in-world group task; char-exact vs CSV 2026-07-28);
%     (P2.2/2.3) "question the third term asked" REFRAMED as the
%     OBSTACLE the deployment never solved; "pedagogically prosaic" +
%     "world justified by compulsion is not justified at all" REMOVED
%     (author: Plato's cave is still a world) → replaced by (a) Aristotle
%     seven-causes compulsion-painful doctrine — paraphrase in prose +
%     verbatim footnote "So all acts of concentration, strong effort,
%     and strain are necessarily painful; they all involve compulsion
%     and force, unless we are accustomed to them, in which case it is
%     custom that makes them pleasant" VERIFIED char-exact vs
%     rhetorical_ontology Rhetoric_(2014) Clean Copy pdf p. 32, cite
%     I.11 1370a12–14 [marginal "15" lands at the following clause —
%     range pinned conservatively]; (b) G&G candle-lighting grounding
%     (USER-ORDERED cross-part reference): start-button-as-first-lesson,
%     grip/trigger, "exploration over instruction" = Part II v4's own
%     phrase reused internally, paraphrase only, no quote; (c)
%     choice-architecture doctrine: assignments completable outside
%     (YouTube + course page) OR everything in one place in-world;
%     illusion-of-agency clause rendered ("more arranged than open...
%     out of conscious recognition"); one-roof convenience.
%     (P2.4) task-on-the-chain ¶ added: chain links specified
%     (perception → phantasia's rendering → committal judgment w/
%     emotion concretizing → act; vocabulary per M3.3-v3 canon);
%     INCORPORATION-ACHIEVED correction (design cannot manufacture it —
%     achievement is the student's; design makes it achievable);
%     "in every way, shape, and form" close per user wording.
%     (P3.1) for-the-sake-of-which ITALICIZED; (P3.2) definition = "the
%     end from which every decision, conscious or otherwise, is
%     perceived, evaluated, understood, and enacted" (user wording);
%     (P3.3) "Routes are what vary beneath it" → concrete pursuit-ways
%     sentence; (P3.4) "same grade" → "same end"; (P3.5) YouTube-open
%     cases = no-strong-computer + quick-verify (user wording).
%     M3.6 "One voice, both poles" FIXED → "both sides" (backup
%     .backups/20260728T091044/; pole ban reaches it per user ruling).
%     v5 RULINGS APPLIED 2026-07-28 (third round): "Aristotle explains
%     why"; "the fumbling the lighting process takes"; ¶4 swap — USER'S
%     CANONICAL FORMULA "Rhetorical incorporation is a mode of
%     being-in-the-virtual-world: it is something curated, a
%     harmonization between the student and the virtual world within
%     which they dwell" (user dictated "user"; rendered "student" for
%     register consistency — FLAG for confirmation); ¶5 curating-the-
%     find sentence BROKEN UP (parity condition → past-parity elements,
%     one sentence each: prompts/target, captions/citable, deliverable-
%     linked finds, the waiting meeting, one-more-reason close). TWO NEW
%     STANDING RULES MEMORIZED: chain-node walkthrough = fundamental
%     analytical form ([[feedback-chain-node-analysis-fundamental]] —
%     user "Note to System"); bare "incorporation" = Calleja ONLY
%     ([[feedback-rhetorical-incorporation-naming]] — swept, 1 fix).)

What remains is design, and the diagnosis dictates its form. A design
lever, in this analysis, is an \textit{archē}: Aristotle's fourth
sense of the word names ``[t]hat from which \ldots\ a thing first
arises, and from which the movement or the change naturally first
proceeds'' (\textit{Metaphysics} V.1, 1013a7--10)---the productive
source from which a change first comes. To call something a design lever is to make a two-part claim: that one
of the deployment's documented failures---the quarter of the class
stopped at the installer, the forty-seven students who chose YouTube
over the world, the group meetings that never convened---has a definite
productive source in the design itself, and that design can reach that
source and change it. The five levers that follow are ranked by the
station of the actualization chain each one works on: the way into the
world (A\textsubscript{0}), the act the world houses (A\textsubscript{3}
to A\textsubscript{4}), the object the seeking is for, the convening of
the WE, and the shape of the watching's time. Ranked this way, the
levers recapitulate the diagnosis: each names a place where the
deployment fell silent, and the productive source design must work on
there.

Five levers carry that claim, and each can be given a plain name: the
\textit{door}, the \textit{task}, the \textit{find}, the
\textit{meeting}, and the \textit{while}. The \textit{door} is the way
into the world---the download, the installation, the sign-in, the
machine all of it must run on. A quarter of the class never got past
it; for those students the world disclosed nothing, because it never
opened. Everything design does inside a world waits on this lever,
since a world is only as good as its way in: a browser-deliverable
build, Mac support, and entry without setup re-admit the students the
\textit{door} turned away. The students proposed the same remedy
themselves: ``Provide a cloud version of the VR. Meaning, provide the VR
setup and ready to go on a cloud machine provided by UCI. This will help
for inclusivity in the VR'' (S25-R059/Q9)---entry without local
installation, asked for in the name of exactly the students the
\textit{door} excluded.

The \textit{task} is what there is to do once inside, and its absence
is the obstacle the deployment never solved. As deployed, every
assignment the course gave could be completed without entering the
world; the watching was mandated, but nothing had to be done where the
watching happened---and students asked for exactly what was
missing.\footnote{One proposed it in detail: ``Maybe an interactive
surgical platform where a group can work together to dissect/do
surgery to further understand what a surgery consists of \ldots\ it
has steps so a group can virtually perform the surgery''
(W25-R035/Q10).} The field's founding account of what
three-dimensional learning environments are for named this obstacle in
advance. What such environments distinctively offer, Dalgarno and Lee
argue, are tasks: ``well-designed 3-D VLEs can enable learning tasks
that are not possible or not as effective in 2-D environments''
(``Learning Affordances'' 25)---which is what the student's surgical
platform proposes, almost
to the letter. The same paper issues the warning this deployment went
on to demonstrate: comparisons that hold the learning design constant
across a 2-D and a 3-D environment ``would be likely to only
demonstrate the trivial fact that if the unique affordances of 3-D
VLEs are not harnessed within the learning design, there will be
minimal unique learning benefits'' (25). The deployment ran that
comparison without meaning to---the same design, watch the footage, on
both routes---and the cohort returned the predicted verdict in its own
words: ``the youtube videos were enough'' (W26-R033/Q9). A second
proposal closes a loop the analysis left open. The student whose praise
stood over an emptiness they could not name---the divergence register's
culprit-less voice---returned, in another answer, naming the missing
thing exactly: ``I think there should be reflections in the VR
environment. While being immersed in the procedure is fascinating, having
reflections afterwards would allow us to look for the main takeaways from
these procedures and allow us to think critically about clinical needs''
(W26-R016/Q9). What the second-form echo could not locate, the same
student's design request supplies: an act of reflection, housed where the
watching happens. One way to
remove the obstacle is requirement: make
the work possible only inside the world, and the rooms will fill. A
compulsory world is still a world. What compulsion changes is how the
world is entered and felt, and Aristotle explains why: what is done under
compulsion runs contrary to nature, and for that reason it is
painful---felt as force unless custom has long since softened
it.\footnote{``So all acts of concentration, strong effort, and strain
are necessarily painful; they all involve compulsion and force, unless
we are accustomed to them, in which case it is custom that makes them
pleasant'' (\textit{Rhetorica} I.11, 1370a12--14).} The dissertation's
earlier analysis of a tutorial world, \textit{Gnomes \& Goblins},
shows the alternative in miniature. That experience also needed a
beginning---no world starts without one---and it designed its start
button as a first lesson: to enter, the player must light a candle,
and the fumbling the lighting process takes introduces the grip, the trigger,
and the world's whole philosophy of exploration over instruction. The
requirement did not disappear; it was made to teach. The deployment's
version of the same design is choice. Let the assignments remain
completable outside---the YouTube platform and the course page stay
open---and build the environment so that everything can also be done
in one place, inside the world, where the footage, the finds, the
notes, and the group already live. A student who enters on those terms
enters by their own act. The force Aristotle warns of drops out of the
equation---or, where the choice is more arranged than open, out of
conscious recognition, which for the experience amounts to the same---
and the world gains a convenience no other route offers: the course's
whole work gathered under one roof.

What the \textit{task} should be built toward, the actualization chain
specifies link by link. The chain traces how a world becomes an act:
perception takes the world in; \textit{phantasia} renders what was
taken into a determinate image; judgment commits to what the image
presents, and emotion concretizes in that commitment; desire, so
formed, completes itself in an act. A task built on the chain gives
every link its material inside the world: footage worth the taking-in;
prompts that give the forming image something determinate to resolve
on; judgments the deliverable asks the student to commit to and
defend; and an act---the need noted, captured, submitted---housed in
the same rooms where the seeing happened, so that nothing in the
world's workings sends the student elsewhere to complete what the
world began. Design cannot manufacture rhetorical incorporation.
Rhetorical incorporation is a mode of
\textit{being-in-the-virtual-world}: it is something curated, a
harmonization between the student and the virtual world within which
they dwell; what design does is make it achievable, clearing the way
at every link.
Build the task so, and the world's existence stops needing an
argument: in every way, shape, and form that matters to the student's
own end, the world is perceived as better than any alternative.

The task lever's pedagogical warrant is already established at
meta-analytic scale. Across health-professions education, training
that puts the learner's hands on the procedure---simulated, not
real---has been tested against no training at all in 609 studies
enrolling 35,226 trainees, with large pooled effects on knowledge and
skills; the evidence is so uniform that the reviewers close by
doubting the comparison is worth running again: ``with only 4\% of the
outcomes failing to show a benefit, we question the need for further
studies comparing simulation with no intervention'' (Cook et al.\
987). That meta-analysis is medium-agnostic---mannequins and cadavers
count alongside screens---so it grounds the task's warrant without
ranking routes; the finer work belongs to the closest analogue this
deployment has in the literature. Dubovi, Levy, and Dagan built the
kind of world the student's proposal asks for: a non-immersive desktop
simulation of a clinical procedure---medication administration, in a
virtual hospital ward---where nursing students assess the patient,
prepare the dose, administer it, and monitor what follows. Students
who trained in that world scored higher on the study's knowledge
measures than a matched cohort taught the same material by lecture, and
the advantage was still undiminished when the students were retested
five months later (23). The study's sharpest
finding is the watching-versus-doing distinction made empirical: of
ten students whose sessions were recorded, nine answered the procedure
worksheet correctly, and nine nonetheless selected the wrong
medication on their first attempt inside the world---errors that fell
to near zero by the second scenario (23--24). Knowing the procedure on
paper and being able to perform it are different attainments, and only
a world with a task in it exposes the difference and then closes it.
The authors' design principle is the task lever's own conclusion,
reached by measurement: ``ensuring that participants gain sufficient
control of the environment is an important condition for learning,''
for ``presence is not a unique attribute of the environment, but it is
induced by the fidelity of representation and the high degree of
interaction and user control'' (25). That last sentence reaches past
the task to the door and the rooms: nothing about a desktop world
forecloses presence; fidelity---to a lesser extent---and interactive
control and agency---to a much greater extent---produce it.

The \textit{find} is the object of desire, and here the diagnosis must
be exact, because the student under mandate is not desireless. Every
student arrives carrying a governing desire that organizes the rest:
to pass the course, for the sake of the degree, for the sake of the
life the degree serves. In the vocabulary this dissertation has
carried from Heidegger, that is the student's
\textit{for-the-sake-of-which}---the end from which every decision,
conscious or otherwise, is perceived, evaluated, understood, and
enacted. What varies beneath that end is the way it is pursued: which
platform to open, where to watch, whether to watch alone. Watching on
the YouTube platform and watching in the virtual hospital are, to that
governing desire, two routes to the same end, and if the routes look
equally effective, the more convenient one wins every time. Curating the \textit{find} therefore means shaping the world so
that the student's own weighing lands inside it. The first condition
is parity: the build must at least equal the YouTube platform in video
quality and in everyday conveniences, because no curated object
survives a worse player. Past parity come the elements only the world
can add. Prompts name, for each procedure, what a needs-finder watches
for, so the looking has a target. Captions and transcripts mark the
moments that matter, so the target can be found and cited. Each find a student captures should feed directly into the group's
graded deliverable---an unmet need identified in the footage becomes an
entry the innovation notebook and the final report are built from---so
that the desire for the find is not a new desire the course must create
but the student's existing desire for the grade, extended one step down
into the world. The meeting waits inside, worth the
\textit{door}'s inconvenience. Each element gives the governing desire
one more reason to choose this room. The choice
must remain real. The YouTube platform stays open for the quick second
viewing, for the student without a computer strong enough to run the
environment, and for the one who wants to jump back in quickly to
verify something about an operation; the aim is not to close the other
route but to make the world the one a free choice keeps landing on.
Agency itself becomes the pedagogical instrument: the student moved by
being handed a world, never a word.

The \textit{meeting} is co-disclosedness convened. The proximity voice
chat affordance exists and every student already belongs to a procedure
group; what design must produce is the convening itself---whatever
brings two students into the same room at the same hour. It is the one
lever whose product no video can imitate, and the \textit{find} can be
shaped to invite it. A procedure fills the frame with more than one
stream of activity---the surgeon's hands, the instruments passing, the
movements of the team around the table---and no single watcher attends
to all of it at once. A prompt that assigns two vantage points, or a
find that asks what happened at the table to be reconciled with what
happened around it, makes a second pair of eyes worth more than a
second viewing: the desirable object doubled into an occasion to meet.
The \textit{while} is the temporal frame---the same while that boredom
stretches. The length complaints ran through every term, but the cut
they imply would be the wrong lever. A procedure edited for pace is no
longer the procedure, and what an edit removes may be exactly the
unremarkable corner of the operating room where an unmet need sits;
the footage must stay whole. The watching does not have to. Segment
the procedures into portions assigned week by week, and set interactive
sequences at intervals inside the
watching---prepared moments that interrupt the drag, catch the
attention as it wanders, and redirect it to what matters. The students
asked for this articulation in nearly those terms. One tied the length
directly to lost focus: ``The long videos of the procedures in the VR
could probably be improved. I didn't really understand what was going on,
and the videos felt too long. If they were split into multiple parts it
would probably be easier to focus on the procedure'' (S25-R017/Q6).
Another proposed the weekly portioning with its diagnosis attached: ``I
think a good think would be to assign certain VR on video to watch weekly
so that all members could stay on track. because the main issues was
people were just waiting too last minute to watch the videos''
(S25-R037/Q6)---segmentation not only against the drag of the while but
against the procrastination a term of total freedom invites. The while is
not shortened; it is articulated. Session length, sequence, and rest
belong to the levers, not the logistics.

% AUTHOR FACT (2026-07-31, batch 5): an instructor-run Discord channel
%   existed for the whole class in ALL surveyed terms (also in the S26
%   syllabus). Context for the meeting-lever remedies: the meeting's
%   nearest competitor channel had institutional standing of its own.
% --- M5.2 Collapsing the activation energy of the WE (★ expanded lever.
%     5-¶ architecture per approved plan: ¶1 equilibrium · ¶2 synchrony
%     constraint · ¶3 rung 0 + low ladder · ¶4 ladder top w/ ignition
%     logic + bia guard · ¶5 motto. Vocabulary ruling in force
%     (technical terms kept; chronos = design material, allotted/
%     reserved). ¶¶1–2 DRAFTED 2026-07-28 — in review.
%     Batch verifications (2026-07-28): W25-R060/Q10 (class google
%     sheet of group entry times) + S25-R079/Q6 (scheduled online mixer
%     early in course; NB row also attests a personality test used for
%     grouping — course fact, unused) char-exact vs open_text.csv ·
%     Kreijns Pitfall 2 "the tendency to restrict social interaction to
%     educational interventions aimed at cognitive processes" verified
%     vs corpus/pedagogy PDF, journal p. 336 (pdf p. 2; same page as
%     Pitfall 1, cite form per M1.6) — deploys in batch 2 ·
%     chronos = Part I §1.1 established term (co-implicated w/ kinēsis),
%     term-use only, no quote. Floors/equilibrium facts inherited from
%     M4.3/M3.5 — no re-derivation; "micro-synchrony" named in ¶2 (term
%     feeds M7.1's minimum-dose hand-forward).
%     P1/P2 RULINGS APPLIED 2026-07-28: (P1.1) "a decade of platform
%     design" REMOVED (misread as the author's own 5-year platform
%     work) → "the solo path's frictions have long since been engineered
%     away by the platforms that host it"; (P1.2) "one student confirmed
%     convening in the three terms the class has run" + FLOOR-HONESTY
%     FOOTNOTE (likely more than one convened; survey built without a
%     direct did-your-group-meet item — foresight admission consistent
%     w/ M3.5 ¶1's; count = floor, never a measure); (P2.1) ¶2 opener
%     rewritten concrete ("Every remedy for the meeting draws on one
%     material: shared time") — NEW STANDING RULE memorized:
%     [[feedback-no-vague-paragraph-frames]] (no vague paragraph
%     openers/closers; check first/last sentences in isolation).
%     ¶¶1–2 APPROVED 2026-07-28. ¶¶3–5 DRAFTED 2026-07-28 — in review.
%     Batch-2 verifications: W25-R071/Q9 ledger reuse (M3.5) ·
%     W25-R060/Q10 + S25-R079/Q6 char-exact (this session) · Kreijns
%     Pitfall 2 char-exact p. 336 (named authors in prose per M1.6
%     form) · De Back et al. 5375–76 + van der Meer et al. 9–11 = M1.6
%     locked-bank claims, parenthetical et-al forms · W26-R037/Q10
%     "main pro" fragment RE-VERIFIED char-exact this session — TRUE
%     WORDING IS "The main pro to the platform is you can watch with
%     others" (analysis docs' "the main pro is you can watch with
%     others" = paraphrase; draft quotes the true span w/ [t] bracket) ·
%     Aristotle custom-clause re-quoted from the M5.1 ¶3 footnote span
%     (I.11 1370a12–14, same verified sentence). Ignition/bia guard per
%     M3.3 ¶5 seven-causes inherit (archē-outside vocabulary); "graded
%     once if ignition needs it"; motto rephrased per vocabulary ruling
%     ("reserve the least chronos," never "spend the minimum").
%     v2 RULINGS APPLIED 2026-07-28 — ¶¶3–5 REWRITTEN as ¶¶3–4 (M5.2 now
%     4 ¶¶): LADDER/RUNG METAPHOR REMOVED (implied hierarchy; replaced
%     by plain ordering "by how much shared time each asks" +
%     first/second/third/fourth). (P3.1) AUTHOR FACT BANKED: the course
%     page ALREADY DESCRIBES the co-watching functionality → "syllabus
%     line" dropped; awareness remedy sharpened to description-tried-
%     and-failed → DEMONSTRATION. (P3.2) traces explained concretely
%     (trace shows the group uses the space; cannot bring two students
%     in at the same hour; no sign of yesterday's visit creates today's
%     meeting). (P3.3 + P4.1) AUTHOR DESIGN REPLACES system-proposed
%     default hour (imposing = moot + removes agency): SCHEDULING
%     ASSIGNMENT — each group settles its weekly time via a
%     group-scheduling tool (user's example: whenavailable.com — kept
%     generic in prose) + MINUTES FOLLOW-UP assignment (brief document:
%     what watched/found/settled; occurrence graded, timeslot never).
%     (P4.2) dyad spelled out ("a pair of groupmates…two students, one
%     procedure, one shared hour—or the full procedure group").
%     (P4.3) bia failure-mode EXPLAINED plainly (term-after-term
%     minutes-only meeting = origin outside agent, felt as force,
%     compliance not practice; opposite sequence spelled out w/
%     W26-R037 main-pro experience). (P5.1) ¶5 DELETED; goal folded
%     into ¶4 close: "an available affordance turned into a consciously
%     selected practice" (user wording).)

Of the five levers, the \textit{meeting} asks the most of design,
because everything needed for it already existed and it lapsed anyway.
Co-watching is a coordination game. Two students who would each get
more from watching together can each still watch alone, and the solo
path asks nothing of them: no schedules aligned, no invitation risked,
no hour kept. The meeting asks all three, and it asks them of all
students in the same group simultaneously. A coordination game of that shape has two stable
outcomes---everyone convenes, or no one does---and which outcome a
cohort lands on is decided by the default. Left to itself, the default
decides for solitude: the solo path's frictions have long since been
engineered away by the platforms that host it, while the meeting's
demands stand untouched. What the collected data showed is therefore
an equilibrium, not a verdict on the students: co-watching solved in
principle---the affordance built, the groups assigned---and unsolved
in practice, one student confirmed convening in the three terms the
class has run.\footnote{It is likely that more than one student
convened in the space with their group. The survey, however, was built
without the foresight of a question that asked directly whether the
assigned group had used the environment to meet, so the collected data
cannot say how many groups actually did. The count is a floor---the
number of responses that happen to confirm a convening---never a
measure of how many occurred.} The lever's name
is exact. What must collapse is the activation energy of the meeting:
the demands that stand between an assigned group and its first shared
hour.

Every remedy for the meeting draws on one material: shared time. The
course was built, deliberately and for good reason, to need none of
it. Its students enroll across nine campuses, quarters set against
semesters, with no timetable in common; a synchronous version of the
course was never in the design space, and an in-person one still
less. Whole-cohort
synchrony is therefore closed to design. What remains open is small:
the overlap of two schedules, a procedure group's single shared
hour---micro-synchrony, the dyad-sized coincidence the systemwide
constraint still permits. Time, in this frame, is a design
material. The deployment allotted all of it to access---any student,
any hour, any place---and that allotment was right; the error was only
that nothing was held back. The meeting needs a small reserve of
shared time, and the deployment reserved none. The design problem is
exactly that narrow: not to restore the classroom hour the course
simply could not implement, but to reserve, inside a schedule built on
total freedom, the one small overlap the WE requires.
% RESOLVED 2026-07-31 (user, close-out): semester-UC question verified by
% the author — the batch-5 facts stand (open to Berkeley/Merced; no
% semester-campus enrollees thus far; rendered in M0 ¶6).

The remedies can be ordered by how much shared time each asks of the
group, and the first asks none: demonstrate the capability. Twelve of
the thirteen students who asked for a way to watch together asked in
wording that implied the platform had none---``if there was a way to
interact with other users who are online at the same time''
(W25-R071/Q9)---and this despite the course page, which states plainly
that groups can meet and talk inside the environment. The capability
was described in writing and asked for as if absent; what description
did not accomplish, demonstration can---an onboarding walkthrough, a
live demonstration in the first week's module, the capability seen in
use rather than read about. The second remedy also asks no shared
time: let the world carry traces of the group---who entered today,
which rooms a groupmate visited. A trace shows a student that their
group uses the space; what it cannot do is bring two students into the
space at the same hour. Watching together by proximity voice chat happens
only when both students are present at once, and no sign of
yesterday's visit creates today's meeting; traces point toward the
meeting without producing it. The third remedy is an assignment, and
it produces the meeting's hour while leaving the choice of hour
entirely with the group. Early in the term, each procedure group uses
a group-scheduling tool to settle the weekly time its members can
share; the time is the group's own finding, not the course's
imposition. The students suggested this remedy directly themselves:
one proposed ``setting up a class google sheet where times are set by
each group to enter the VR program/watch the videos'' (W25-R060/Q10);
another, a ``scheduled online mixer or networking session early in the
course'' to make finding groupmates easier (S25-R079/Q6). The second
proposal reaches past scheduling into affiliation, and the
collaborative-learning literature says it must: Kreijns, Kirschner,
and Jochems name as their second pitfall ``the tendency to restrict
social interaction to educational interventions aimed at cognitive
processes'' (336)---a meeting made only of assignments starves the
social ground the meeting stands on. The students also supplied a reason
the world's own meeting-form can carry that social ground better than a
video call: ``I think a virtual group will be nice, students could make
their own charactors, students who don't know each other may shy to speak
in reality, but if there is a virtual image for them, they might have
less stress talking to others'' (W26-R001/Q10). Two more voices
corroborate the mechanism from the camera's side, asking for ``group
communications that do not wish to do video camera meetings''
(S25-R002/Q10) and noting that ``some of us don’t wanna have their
cameras on'' (S25-R074/Q10). The avatar lowers the exposure the
invitation risks; for strangers assigned to one another across nine
campuses, that reduction works on exactly the activation energy this
lever must collapse.

The fourth remedy closes the design with accountability, still without
imposing a timeslot. Each group's weekly meeting carries a brief
follow-up assignment: minutes, submitted after the session---what the
group watched, what it found, what it settled. The occurrence is
graded; the timeslot never is; the hour remains the group's own
settlement from the scheduling assignment. Small meetings are the
right form. In immersed collaborative environments, learning gains run
inverse to group size---smaller groups, larger gains (De Back et al.\
5375--76)---and reviews of monitor-based collaborative worlds judge
voice communication essential (van der Meer et al.\ 9--11);
groupmates chatting through proximity voice chat---two students, one
procedure, one shared hour---or the full procedure group at once, is
exactly the form the evidence favors. A requirement stands at the start of this
design, and Aristotle's qualification, quoted above, explains why it
can stand there without poisoning what follows: the forced is painful
``unless we are accustomed'' to it, and then ``it is custom that
makes'' it pleasant. The assignment's work is to establish and naturalize the custom---to
bring each group through its first meetings, while the demands are
still unfamiliar and the worth still unexperienced. What it curates
is, in Aristotle's terms, a habit: among the seven causes of action,
the one that stands nearest to nature, for ``as soon as a thing has
become habitual, it is virtually natural'' (\textit{Rhetorica} I.11,
1370a6--8). If that were all
the design ever did---if, term after term, students met only because
minutes were due---the meeting would have remained what the
seven-causes analysis names compulsion: an act whose origin stays
outside the agent, done under force and felt as force, yielding
compliance rather than practice. The design intends the opposite
sequence. The requirement produces the first meeting; the first
meeting lets the students experience what the survey's one confirmed
convener reported---``[t]he main pro to the platform is you can watch
with others'' (W26-R037/Q10)---and from then on the group meets
because it wants to, the cause now inside the students rather than in
the requirement. The goal, throughout, is exactly that movement: an
available affordance turned into a consciously selected practice. That
movement is also a discovery the actualization chain makes diagnosable.
Requirement and practice can look identical from the outside---the same
group, the same room, the same hour---and the chain names where they
differ: in the \textit{arch\=e} of the act, its origin lying in the
requirement or in the students' own desire. A course that watches only
attendance cannot tell compliance from practice; a course that asks
where the act's origin now lies can see, term by term, whether the
affordance is still being forced or has begun to be chosen.

% --- M5.3 The motto and the mirror (1 ¶ per approved plan. DRAFTED
%     2026-07-28 — in review. Verifications: Mark capture-economy claim
%     = paraphrase, cite (Mark 155–65) per crosswalk §3 catalog-backed
%     range — VERIFY Mark is in Works Cited at assembly (M0/M1 may
%     carry); Wendt "facilitate profitable mistakes" char-exact vs
%     corpus/new_media Design for Dasein PDF p. 162 (pdf 164; OCR noise
%     in layer, wording confirmed) — "profitable" appears ONLY inside
%     Wendt's quoted phrase, not our own metaphor; Wendt principle 10
%     ("advocate for preferable worlds; refuse dark patterns") = entry
%     manual's distillation, deployed as paraphrase. Soul-line = Part II
%     v4's "reaching into and reshaping the soul" reused as internal
%     formula (the author's "hack their soul" doctrine landing at the
%     ethics beat as reserved). "Settled ways" not "habits" (hexis/habit
%     ban). M5.2's for-the-sake-of-which + capture/restart mirror from
%     cross-case §3.3.
%     v2 RULINGS APPLIED 2026-07-28: (1) attention-economy passage
%     EXPANDED w/ Mark catalog verbatims — "who you are, how you feel,
%     what you do when and where" + "to capture your attention" (155,
%     item #28); "occur spontaneously and automatically without
%     cognitive processing" + "when you are low in attentional resources
%     or in a bored or a rote state" + "you can't help but respond to
%     them" (159–60, items #34–35; all MARK-VERBATIM register).
%     (2) NO-STALL DOCTRINE (user): the chain NEVER stalls — "stalled
%     chain" lines DELETED; capture rendered as full chain walkthrough
%     (no determinate object → pre-charged object supplied → perception
%     seized → act completes past the cognitive orientations, "without
%     cognitive processing" = Mark's own warrant); "Nothing in the chain
%     is broken; the chain runs. It is run deliberately past the very
%     links where a person's own ends would get a vote." Doctrine added
%     to [[feedback-chain-node-analysis-fundamental]] + "stall" added to
%     sweeps; M5.1 ¶1 "a stall" → "a flight" (approved-text touch).
%     (3) DARK-PATTERNS USE CORRECTED (user suspicion CONFIRMED vs PDF:
%     Wendt lists dark patterns as one nefarious effect — "persuasive
%     design attempting to directly influence behavior--dark patterns
%     attempting to facilitate profitable mistakes" — NOT a name for
%     attention capture): capture-equation dropped; Wendt now supplies
%     the designer's norm, "advocating for preferable worlds" VERIFIED
%     char-exact p. 162 (pdf 164). (4) M5.2 ¶4 HABIT INSERTION:
%     Aristotle's habit curated — "as soon as a thing has become
%     habitual, it is virtually natural" VERIFIED char-exact pdf p. 31
%     (Rhetoric I.11 1370a6–8, M3.3's cite; habit = seven-causes ethos,
%     kept distinct from hexis per standing ban). (5) soul-line = "To
%     design this way is to hack the soul." (user wording).
%     v3 2026-07-28: "video alone cannot" (user). ¶ APPROVED —
%     **M5.3 COMPLETE. M5 MODULE COMPLETE (M5.1 6¶¶ · M5.2 4¶¶ ·
%     M5.3 1¶ — all approved). M5.2 habit insertion approved.**)

The five levers share one motto: design against the question ``why not
just YouTube?'' Every element of the world must do something video
alone cannot---give a task, hold a place, convene a meeting---or the
element loses, rationally, to the platform that plays the same footage
with less friction. A craft this effective deserves candor about what it
is, because the attention economy practices the same craft in mirror
image, and the actualization chain shows exactly where the mirror
sits. Gloria Mark's account of attention capture is precise about the
mechanism. The algorithms learn ``who you are, how you feel, what you
do when and where'' and use it ``to capture your attention'' (Mark
155). Their instruments are targeted solicitations built to elicit
what Mark calls lower-order emotions---reactions that ``occur
spontaneously and automatically without cognitive processing''---and
they are timed for the moments ``when you are low in attentional
resources or in a bored or a rote state,'' when ``you can't help but
respond to them'' (159--60). Read through the chain, the design is
exact. A bored viewer still desires---the governing end never
sleeps---but the moment offers no determinate object to want. The
application supplies one, pre-charged: perception is seized by the
notification; the image arrives already pleasurable; and the act
completes before the cognitive orientations have done their
work---no deliberative weighing of the object, no committal judgment
about its worth. Nothing in the chain is broken by such design; the
chain runs. The pedagogical craft uses the same knowledge in the
opposite direction: it curates a worthy object of desire, informs the
judgment that weighs it, and gives committed desire an act to complete
in the world---every link furnished, none bypassed, the completion
serving the student's own \textit{for-the-sake-of-which} rather than a
platform's. Wendt states the designer's side of the same boundary:
what designers make acts on the world long after they are done with
it, the nefarious end of the spectrum runs from persuasive design to
dark patterns, and the calling that stands against it is
``advocating for preferable worlds'' (\textit{Design for Dasein} 162).
To design this way is to hack the soul. That is no reason to refuse
the craft; it is the reason the craft bears the responsibility it
does, and must answer for the dispositions it lays down---the settled
ways of looking, judging, and meeting that outlast the course. A world
that teaches by being dwelt in shapes its dwellers, and the designer's
obligation is to shape them toward their own good---openly aimed, and
defensible to the very students it moves.

% ========== END INLINED FILE: DESKTOP-M5-DRAFT-v1.tex ==========

% ===== M6 — ASSEMBLY-PROVISIONAL head
\section{Limitations}
% ========== BEGIN INLINED FILE: DESKTOP-M6-DRAFT-v1.tex ==========
% ============================================================
% Part III — Desktop Section — M6 (Limitations) — DRAFT v1
% Mode: Stage 3 guided drafting. Function per outline: consequences-
%   reprise, NOT inventory — the full statement of the standing ceilings
%   lives at M2.2; M6 draws what the ceilings mean for the claims just
%   made. 1 ¶ (M6.1). Conventions as DESKTOP-M3/M4/M5 drafts; ALL
%   standing sweeps in force (And / record / pole / ladder / stall /
%   cost-metaphors / bare-incorporation / vague-frames).
% ANCHORS: F8 + cross-case §4 (limits ledger). DEP: M2.2 full statement.
%   Outline's "what each ceiling costs the argument" rendered as "what
%   each ceiling means for the claims" per vocabulary ruling.
% Ceilings mapped to claims: single channel → form-2 echoes only (M4.1
%   ¶6, M3.6); retrospective → felt time mute (M4.1 ¶7 guard held);
%   content ceiling → no learning-outcome claims (scopes M5's
%   better-than-alternatives to AIM, testable only by the redesigned
%   deployment); term composition → arcs = trends not tests (M2.2 DEP
%   lock); category caution → structural grammar, never Grundstimmung
%   detection. ¶1 DRAFTED 2026-07-28; 3-point revision (form-2 claim
%   shape spelled out / content-ceiling unpacked to 3 items / apparatus
%   as categories-not-findings) APPROVED 2026-07-28. M6 COMPLETE.
% ============================================================

Five ceilings bound everything this section has claimed, and each has
a consequence worth stating plainly. The survey is a single channel of
evidence, so the second form of boredom appears in this analysis only
as echo, never as detection. Wherever a response was read as engaged
on its surface and empty underneath---a student conceding the world's
immersion while reporting that nothing in it was the highlight---the
surface came from the student's words, but the emptiness came from
Heidegger's frame; no second measure of these students exists to
confirm or overturn such a reading, so every claim of that shape
stands or falls with the frame. The instrument is
retrospective, so felt time is nearly mute in the collected data;
where the analysis has spoken of the while stretching, it has spoken
from the form's grammar and from the laboratory's separate evidence,
not from these students' reports. The survey's three content
items---the questions testing whether students grasped the procedures
themselves---drew correct answers from nearly every student in every
cohort; scores bunched at the top of a scale leave no variation to
analyze, so those items cannot show whether one format taught better
than another, and no claim in this section concerns learning outcomes.
Effectiveness,
wherever the word has appeared, is scoped to what the chain-completion
evidence can carry---presence attested, interest sustained, acts
housed or unhoused---and the design section's ``better than any
alternative'' names an aim whose test would be the redesigned
deployment itself, not a result already in hand. The three cohorts
differ in size and likely in composition, so the arcs---frictions
falling, the value-question rising---are patterns observed but not yet
definitive. Heidegger's account of boredom, finally, supplied this section
with categories, never with findings: the three forms, the paired
moments of being left empty and being held in limbo, and the
two-channel rule functioned as a scheme for sorting what students
wrote---which complaints name a culprit and which do not---rather than
as instruments that measured an inner state. A fundamental attunement
(\textit{Grundstimmung}) is, on Heidegger's own account, something no
retrospective questionnaire answer could establish; no reading in this
section claims to have detected one, and none should be read as if it
had. What survives these
ceilings is the section's actual yield: where the world fell silent,
what its difference from the video of itself must be made of, and how
design can convene what the collected data shows was wanted.

% ========== END INLINED FILE: DESKTOP-M6-DRAFT-v1.tex ==========

% ===== M7 — ASSEMBLY-PROVISIONAL head (outline's own title)
\section{Implications and Close}
% ========== BEGIN INLINED FILE: DESKTOP-M7-DRAFT-v1.tex ==========
% ============================================================
% Part III — Desktop Section — M7 (Implications and Close) — DRAFT v1
% Mode: Stage 3 guided drafting. Plan APPROVED 2026-07-28:
%   M7.1 = 2 ¶¶ for the empirical program (gesture scale per
%   future-work rule — gestures, not plans); M7.2 = 1 ¶ for the
%   dissertation, SECTION SCALE ONLY per user ruling (O-11 deferred;
%   if O-11 resolves toward in-section conclusion, M7.2 is the
%   EXPANSION SITE — do not treat as dissertation-level close).
% RULINGS BANKED 2026-07-28:
%   (1) M7.2 register = section-scale close + O-11 expansion marker
%       (recommendation approved).
%   (2) Task-lever expansion = KEY QUOTATIONS woven into M5.1's
%       task-lever treatment (NOT a lit review, NOT an M1 supplement
%       pass) — queued as post-M7 pass on M5.1.
% STALE OUTLINE POINTERS corrected in rendering: "cost-ranked remedies
%   ladder" → ordered by the shared time each asks (cost-metaphor ban +
%   ladder removal); M7.2's "stall, flight, and unconvened WE" →
%   silence at the act / furnished flight / unconvened WE (no-stall
%   doctrine — the chain runs always).
% HARMONIZATION FLAG (assembly): the three-state proposal has a
%   gesture-scale rendering in III-5-Design-DRAFT-v1.tex ("restless
%   search, occupied-but-hollow, withdrawn"; surface/depth disagreement
%   as diagnostic; sleep-exit terminal). M7.1 ¶1 hands forward WITHOUT
%   reusing that wording; reconcile the two gestures once arc decision
%   (#1) resolves which section carries the full statement.
% LIT-SLOT (verified 2026-07-16, both sides): FCM-phenomenology ×
%   instrumentation = UNBUILT combination — 0/15 hard-channel units
%   carry FCM bridge edges; none of 34 VR-pedagogy units combines the
%   phenomenological boredom apparatus with instrumentation; nearest =
%   EEG-workload "boredom"-band labels (Makransky, Terkildsen & Mayer
%   2019; Parong & Mayer 2021) + single self-report items; DA-10
%   corroborates at network level (no citation traffic joins halves).
% ALL standing sweeps in force (And / record / pole / ladder / stall /
%   cost-metaphors / bare-incorporation / vague-frames).
% VERIFICATION LEDGER (index-entry-first per standing rule):
%   ped-sec-10 (Makransky, Terkildsen & Mayer 2019): EEG workload from
%     9-ch Advanced Brain Monitoring system, bands boredom/optimal/
%     overload, entry loci pp. 226–27 + 229–30 → cite 229–30.
%   ped-sec-13 (Parong & Mayer 2021; online-first 2020 — cite VoR
%     2021): boredom = instrument's own label for lowest workload band
%     (FAITHFUL, 234–36); IVR sig. higher boredom-band proportion +
%     lower EEG high-engagement (FAITHFUL, 234–36); presence d=1.74,
%     enjoyment d=1.20 both IVR>PPT sig. (FAITHFUL, 233–36); transfer
%     PPT>IVR sig. d=−0.71, retention marginal p=.058 (FAITHFUL, 233)
%     → "learned less by the study's transfer measure" (retention held
%     to marginal — NOT claimed). Band thresholds NOT printed (entry
%     shows ≤.40 vs <.40 discrepancy — avoid until PDF-pinned).
%   MLA: Parong & Mayer short title required IF ped-sec-08 (2018) also
%     enters Works Cited — normalize at assembly. Works Cited adds
%     queued: Makransky/Terkildsen/Mayer 2019 + Parong/Mayer 2021.
% M7.1 COMPLETE 2026-07-28 (2 ¶¶): ¶2 approved first pass; ¶1 v2
%   approved (concrete opener + Parong/Makransky specifics — the
%   field's own instruments already produced the surface/depth
%   divergence, read only as misallocated workload).
% M7.2 (1 ¶) — SECTION-SCALE CLOSE per user ruling 2026-07-28.
%   O-11 EXPANSION MARKER: if O-11 resolves toward an in-section
%   dissertation-level conclusion, THIS PARAGRAPH is the expansion
%   site; as drafted it closes the section's own argument and hands
%   the frame back as a finding — do not read it as the dissertation's
%   conclusion. DRAFTED 2026-07-28; v2 rulings applied ("demonstrated"
%   not "prove" — too absolute; silence clause anchored [assigned
%   analysis completable entirely outside; inside, only watching
%   asked]; GREATEST-DESIRE DOCTRINE: never "better route" — the chain
%   shows the person following the greatest desire, stronger outweighs
%   weaker [banked in feedback-chain-node-analysis-fundamental]; rest
%   approved) — revised spans APPROVED 2026-07-28. M7.2 COMPLETE.
% M7 MODULE COMPLETE 2026-07-28 (M7.1 2 ¶¶ · M7.2 1 ¶) — DESKTOP
%   SECTION DRAFTING COMPLETE M0–M7.
% ============================================================

% --- M7.1 For the empirical program (2 ¶¶, gesture scale) ---

The virtual-learning literature already measures something it calls
boredom, and the shape of its instrument shows exactly what is
missing. In the field's most instrumented media-comparison
experiments, boredom is the name of the lowest band of an EEG workload
metric: whatever falls beneath the scale's lower threshold is
classified as ``boredom,'' so that the word means nothing more than
workload's absence---the bottom of a single axis (Makransky,
Terkildsen, and Mayer 229--30; Parong and Mayer, ``Cognitive and
Affective Processes'' 234--36). One of those experiments shows why the
label is not enough. In Parong and Mayer's media comparison, the
students in the headset reported the highest presence and the highest
enjoyment in the study, with large effects on both; the same students
were classified into the boredom band at a significantly higher rate
than the slideshow group, and they learned less by the study's
transfer measure (``Cognitive and Affective Processes'' 233--36). The
self-report channel found those students full; the physiological
channel found them empty; the outcome measure sided with the empty
reading. That is the divergence this section could only hear as an
echo in a student's prose, produced there by independent instruments
in a single published experiment---and the study can read it only as
misallocated workload, because boredom-as-low-workload has no
structure: it cannot distinguish the student searching for something
to attend to from the student occupied but held by nothing from the
student withdrawn from the world altogether. Heidegger's account
distinguishes exactly these, and no published study joins it to an
instrumented reading of learners inside a virtual world. The
combination is unbuilt, and the laboratory study this Part turns to
next was designed inside the frame that makes it buildable: a surface
reading of what occupies the student run concurrently with a depth
reading of what holds them, the disagreement between the channels
serving as the diagnostic.

The unconvened WE leaves the empirical program a narrower question,
and it comes with its own apparatus attached: what is the least shared
time that convenes the WE? The design section's four remedies already
stand in the order the question requires, by how much shared time each
asks of a group---the demonstration asks none, the traces ask none,
the scheduling assignment asks one hour of the group's own choosing,
the minutes make that hour accountable---so the redesigned deployment
is itself the instrument: introduce the remedies in that order, term
by term, and watch for the increment at which students begin to meet.
The collected data has already fixed the baseline any such observation
starts from---thirteen students asking for the meeting, one confirmed
convening---and the question's answer would be a finding about
shared time as a design material: how small a reserve of it suffices.
A further question makes itself known for the first time.

% --- M7.2 For the dissertation (1 ¶, section-scale close) ---

The desktop clinic completes a test the dissertation's frame has been
running since Part I. Part I built the actualization chain from
perception to act. Part II built rhetorical incorporation---one in
substance, dual in being as world-disclosedness and
co-disclosedness---together with the demand of veri(dis)similitude, and
demonstrated the frame against the richest of authored worlds:
story-driven, hand-crafted, entered by choice. This section has now
tested it against the sparest world the dissertation examines: a
virtual hospital of plain rooms and procedure footage, equipped with
proximity voice chat, deployed under course mandate to 241 students.
The vocabulary held. Everything those
students met was authored---the rooms, the footage, the hail spoken
before the world first rendered---and none of that authorship
determined an outcome. The chain ran always; each student's acts
completed where that student's own judgment routed them, most often in
the alternative the course itself had furnished. What the frame
contributes is the ability to say precisely what stood in the gap
between total authorship and undetermined outcome. Without it, the
survey's answers read as a feature backlog and a usage shortfall.
Within it, they resolve into a three-part diagnosis: a world that fell
silent at the chain's final link, since the act the course
assigned---analysis of the procedure---could be completed entirely
outside the world, and inside it nothing beyond watching was ever
asked; a flight that was never a malfunction, since the chain merely
shows each student following their greatest desire---both routes
served the governing end, and the greater desire lay with the more
convenient route; and a hail whose one unmatchable promise, the
meeting, went unconvened for want of a small reserve of shared time. What
the analysis demonstrates, at the scale of 241 students, is the claim
the section opened with. A virtual world on a flat screen can place its
students: seven of every ten application users attested presence on an
ordinary monitor---the pilot's headset rate, with no headset. Presence
alone does not justify a world: every video in the virtual hospital
also played one click away, students weighed the world against the
video of itself, and wherever the world furnished nothing beyond
watching, they completed the chain through the more convenient route.
What could have justified the world is the one capability the copy
cannot reproduce: watching together, in the only room this class has
ever had. The world held that capability open for three terms; one
voice in 873 responses confirms it was ever lived. A wholly authored
world, in short, attunes its students without determining them---and
this deployment has demonstrated both halves of that sentence at
cohort scale. That is what the dissertation carries forward from this
section: not a premise defended, but a finding held.

% ========== END INLINED FILE: DESKTOP-M7-DRAFT-v1.tex ==========

\clearpage

% ===== Appendices =====
% Appendix A — the reproduced pilot article (placeholder: the article PDF is
% bound in at dissertation assembly; the label resolves M0's \ref).
\section*{Appendix A: The Published Pilot Study}\label{app:vr-pilot}
\textit{[The pilot article --- King and Salvo, ``Phenomenological Evaluation
of an Undergraduate Clinical Needs Finding Skills Through a Virtual Reality
Clinical Immersion Platform,'' Biomedical Engineering Education, vol.~4,
2024, pp.~381--97 --- is reproduced here in the final dissertation
assembly.]}

\clearpage
% ========== BEGIN INLINED FILE: DESKTOP-APPENDIX-FRICTION-CATALOG-DRAFT-v1.tex ==========
% ============================================================
% Part III — Desktop Section — APPENDIX: Friction Catalogue — DRAFT v1
% Ordered 2026-07-24 (user directive, M3.1 review): every reported friction,
%   most to least reported. Referenced by the M3.1 ¶2 footnote.
% Data: A1 friction family, all-Fable corpus (refreshed 2026-07-23);
%   counts = coded applications (t3-stats); percentages = unique STUDENTS
%   mentioning the friction, of app users n=123 / YouTube group n=108
%   (strata from Q8; computed 2026-07-24 from 03-coded-corpus × respondents).
% Exemplars: curated bank (04-exemplar-bank) or verbatim evidence spans (03).
% Residual grouping (A1.other): DESCRIPTIVE grouping for presentation only —
%   not codes; the 32 underlying applications are listed in 03 under A1.other.
% Verify all totals against the corpus at assembly (as coding appendix).
% Appendix number + \label at assembly.
% ============================================================

\section*{Appendix: Catalogue of Reported Frictions}

Every friction students reported is collected here, ordered from most to
least reported. A complaint is counted once per coded application under the
friction family of the codebook; the percentage columns count students, not
complaints---the share of the 123 application users, and of the 108 students
on the YouTube route, who raised the friction at least once. Two of the
frictions we predicted during development appear at the bottom of the table
precisely because the survey did not bear them out.

\begin{table}[h]
\centering
\small
\begin{tabular}{lccl}
\hline
Friction & Complaints & App users / YouTube & Example from the survey \\
\hline
Crashes and freezes & 29 & 15\% / 8\% &
  ``I had issues with crashes'' (W25-R064/Q6) \\
Install and compatibility & 26 & 3\% / 17\% &
  ``[m]ore inclusive for macbook users :('' (W25-R042/Q9) \\
Buffering and latency & 25 & 11\% / 8\% &
  ``maybe improving the lags, I would glitch sometimes'' (W25-R095/Q9) \\
Wayfinding and navigation & 21 & 11\% / 6\% &
  ``eliminate having to go up stairs and stuff?'' (W25-R084/Q9) \\
Hardware access & 13 & 3\% / 7\% &
  ``the file being too large, my computer could not handle it'' (W25-R060/Q9) \\
Video fidelity & 10 & 5\% / 3\% &
  ``I feel like it's very blurry in my experience'' (W25-R062/Q9) \\
Viewing proximity & 2 & --- &
  ``surgical screen far in the back which made it hard to see'' (S25-R031/Q9) \\
Avatar occlusion & 0 & --- &
  (predicted; unattested---most students entered alone) \\
\hline
\end{tabular}
\caption{Reported frictions by frequency. Counts are coded applications;
percentages are students mentioning the friction at least once.}
\end{table}

A further 32 complaints name frictions outside these categories and are
carried in the codebook as a residual group. Grouped descriptively, they
comprise: difficulty understanding the footage or the task it served (5,
e.g.\ ``hard to tell what is happening in them,'' W25-R067/Q6); interface
usability (5, e.g.\ ``The VR platform is not very user friendly,''
W26-R039/Q9); distraction and overstimulation (4, e.g.\ ``more distracting
than anything,'' W25-R092/Q9); motion discomfort and disorientation (4,
e.g.\ ``I get motion sickness when watching the VR videos,'' W25-R054/Q9);
residual crash, freeze, and loading complaints phrased outside the crash
category (4); missing or inaudible audio on the surgical footage (3, e.g.\
``Most surgery videos are silent,'' W25-R039/Q9); passivity and filler time
(2, e.g.\ ``a lot of filler space/period of inactivity,'' W25-R025/Q9); and
five single mentions: wayfinding signage, account security, solitude in the
platform (``I was usually alone in the platform,'' W25-R038/Q10), setup
abandonment, and general doubt of benefit. Not counted as frictions are the
19 students on the YouTube route (18\% of that group) who simply reported
never using the platform at all.

% ========== END INLINED FILE: DESKTOP-APPENDIX-FRICTION-CATALOG-DRAFT-v1.tex ==========

\clearpage
% ========== BEGIN INLINED FILE: DESKTOP-APPENDIX-CODING-METHODS-DRAFT-v1.tex ==========
% ============================================================
% Part III — Desktop Section — APPENDIX: Survey-Coding Methods — DRAFT v1
% Approved as "minimal appendix" at coding review Station 12 (2026-07-23).
% Scope ruling: codebook family table + per-code agreement table + merge rules
% + provenance note. 2–3 pages. Numbering/placement decided at assembly.
% Sources: methodology-application/codebook-v1.md · survey-analysis/02-codebook-final.md
% · coding/agreement-stats.txt (annotated 2026-07-23) · scripts/merge-adjudicate.py.
% All figures reflect the post-2026-07-23 all-Fable corpus.
% ============================================================

\section*{Appendix: Survey-Coding Methods}

This appendix documents the coding instrument summarized in the methods
discussion: the codebook's family structure, the reliability results in full,
and the merge rules. The complete codebook, all pass files, cross-model
comparison reports, and adjudication records are retained in the project
archive.

\subsection*{A. Codebook families}

The codebook was fixed before coding began and required no structural revision
through coding. Counts are total applications in the final corpus of 873
responses.

\begin{center}
\begin{tabular}{llrl}
\hline
Family & Scope & Applications & Principal codes \\
\hline
A & Boredom-form signatures & 194 & determinate-culprit frictions (9 sub-tags), \\
  &                          &     & diffuse dissatisfaction, monotony, \\
  &                          &     & divergence flags, withdrawal \\
B & Flight patterns & 50 & video fallback, minimal compliance \\
C & Temporal texture & 14 & time-drag, time-vanishing, objective length \\
D & Design-node codes (evaluation chain) & 256 & perceptual conditions, image quality, \\
  &                          &     & interest/meaning, action affordances, return-use \\
E & Involvement channels & 259 & kinesthetic, spatial, shared, narrative, \\
  &                          &     & affective, ludic \\
F & Incorporation-positive signatures & 30 & presence, world-language, carry-over \\
G & Suggestion taxonomy & 399 & provision, optimization, content, \\
  &                          &     & interactivity, guidance, use-cases \\
H & Emergent (open) & 38 & coder-invented, dispositioned post-run \\
X & Housekeeping & 269 & non-use, contentless, off-topic \\
\hline
\end{tabular}
\end{center}

% NOTE: family totals VERIFIED against t3-stats at assembly 2026-07-29.
% A-family corrected 155 -> 194 (A1.* subtotal 158 [other 32 + crash 29 +
% install 26 + lag 25 + nav 21 + hw 13 + dark 10 + prox 2] + A2 10 + A3 4 +
% A4 13 + A5 9). B 50 / C 14 / D 256 / E 259 / F 30 / G 399 / H 38 / X 269
% all match t3-stats sums exactly. TRANSLATED (3) and MASK (1) are corpus
% annotations, not family codes — correctly outside the table.

Two structural rules governed application. Multi-coding was permitted across
families, with best fit chosen per codable idea within a family; and codes
name signatures in student talk, never certified inner states.

\subsection*{B. Reliability}

Reliability was estimated on a 143-response subsample (16.4\% of the corpus),
drawn from the term the codebook had least seen and coded by two independent
blind passes. Mean Jaccard agreement over full code sets was .878;
Krippendorff's $\alpha$ with MASI set-distance over the same pairs was .82.
Per-code agreement (both passes / one pass only):

\begin{center}
\begin{tabular}{lrrrr}
\hline
Code & Both & Pass 1 only & Pass 2 only & Agreement \\
\hline
G6-USECASE & 29 & 0 & 0 & 100\% \\
D2-A3 & 30 & 0 & 12 & 71\% \\
X-OFFTOPIC & 25 & 2 & 0 & 93\% \\
X-EMPTY & 15 & 1 & 0 & 94\% \\
E-SHA & 11 & 1 & 2 & 79\% \\
G3-CONTENT & 9 & 0 & 0 & 100\% \\
E-NAR & 8 & 6 & 1 & 53\% \\
A1.other & 7 & 1 & 1 & 78\% \\
E-LUD & 6 & 2 & 0 & 75\% \\
G4-INTERACT & 6 & 0 & 2 & 75\% \\
D3-A3A4 & 6 & 1 & 0 & 86\% \\
A1.lag & 6 & 0 & 0 & 100\% \\
B1-FALLBACK & 6 & 0 & 0 & 100\% \\
G2-OPTIMIZE & 5 & 0 & 0 & 100\% \\
G5-GUIDE & 4 & 0 & 0 & 100\% \\
F1-PRESENCE & 4 & 0 & 0 & 100\% \\
E-AFF & 3 & 1 & 1 & 60\% \\
E-SPA & 3 & 1 & 0 & 75\% \\
D4-LOOP & 3 & 1 & 0 & 75\% \\
A2-DIFFUSE & 3 & 0 & 0 & 100\% \\
A1.install & 3 & 0 & 0 & 100\% \\
D1-A2 & 2 & 1 & 0 & 67\% \\
A1.nav & 2 & 0 & 0 & 100\% \\
D0-A0A1 & 2 & 0 & 0 & 100\% \\
C3-LENGTH & 2 & 0 & 0 & 100\% \\
A5-WITHDRAWN & 1 & 0 & 1 & 50\% \\
A1.crash & 1 & 0 & 0 & 100\% \\
\hline
\end{tabular}
\end{center}

% Rows with zero both-pass instances (H-NEW inventions, F2-WORLD, B2-MINIMAL)
% omitted for space; full table in coding/agreement-stats.txt.

Agreement is weakest on the mention-level channel codes, which the merge
handles by union, and strongest on the load-bearing families, which the merge
handles by intersection: the design concentrates coder disagreement exactly
where it is least consequential.

\subsection*{C. Merge rules}

Merging the two passes was deterministic under four rules. (1) Mention-level
families (design nodes, involvement channels, suggestions, incorporation
signatures, housekeeping, emergent codes) merged by union: for a mention, an
omission by one coder is likelier than an invention by the other. (2) The
analytically load-bearing families (boredom-form, flight, temporal) merged by
intersection, with a one-pass-only code otherwise queued for manual
adjudication; a keyword-warrant path for auto-acceptance existed but never
fired---every load-bearing dispute reached manual ruling. (3) Friction
sub-tags: all text-warranted sub-tags kept; the generic sub-tag dropped
whenever a specific one survived. (4) Divergence flags kept at either-pass
sensitivity, with precision enforced by manual review of the resulting
register.

\subsection*{D. Provenance}

All coding passes in the final corpus were performed by Anthropic's Claude
Fable 5 model, blind and independent, under fixed written instructions.
Three batches initially single-coded by a smaller model were discarded and
recoded after a measured cross-model comparison (mean Jaccard .51--.67 against
the .878 same-model benchmark); no coding from the smaller model remains. The
author adjudicated all seven manual corrections and reviewed the apparatus
rule by rule. A dedicated two-pass recount of the shared-involvement evidence
followed the same protocol (93\% exact single-bin agreement, eight boundary
disputes author-adjudicated).

% ========== END INLINED FILE: DESKTOP-APPENDIX-CODING-METHODS-DRAFT-v1.tex ==========

\clearpage

% ===== Works Cited =====
% Compiled + verified at assembly 2026-07-29. Sources of verification:
%   [IDX] = corpus/index entry citation JSON (VR Pedagogy Secondary units)
%   [PDF] = verified against the source PDF this session
%   [M1]  = M1 draft header bank (its PDF-verified ledger)
%   [REF] = pilot article's own reference list (PDF)
% Author-name format follows the M1 bank (all authors listed where banked).
% ****** UNVERIFIED markers per standing missing-source protocol.
% NOTE (harmonize at dissertation assembly): the Aristotle and Heidegger
%   entries must match Part I's canonical Works Cited entries — the forms
%   below are drafted from the corpus PDFs and flagged for that check.
% YEAR DECISIONS carried: Chalmers = 2017 (Disputatio VoR; PDF marked 2016)
%   [PDF]; Mayer/Makransky/Parong = 2022 (online VoR; print issue 2023)
%   [M1+IDX]; De Back = 2021 (online; print issue 2023) [IDX].

\begin{workscited}

\item Aristotle. \textit{The Complete Works of Aristotle: The Revised
Oxford Translation}. Edited by Jonathan Barnes, 2 vols., Princeton UP,
1984. % [PDF — corpus (2014) digital ed. of Barnes; HARMONIZE with Part I]

\item Barrett, Robin C.\ A., et al. ``Comparing Virtual Reality,
Desktop-Based 3D, and 2D Versions of a Category Learning Experiment.''
\textit{PLOS ONE}, vol.~17, no.~10, 2022, e0275119. % [M1+IDX]

\item Biocca, Frank, Judee K.\ Burgoon, and Chad Harms. ``Toward a More
Robust Theory and Measure of Social Presence: Review and Suggested
Criteria.'' \textit{Presence: Teleoperators and Virtual Environments},
vol.~12, no.~5, 2003, pp.~456--80. % [M1+IDX — Harms initials pin resolved]

% ****** PAGINATION VERIFICATION REQUIRED on both Boellstorff entries ******
% (user directive 2026-07-31): 2023 = online version-of-record, paginated
% 1-10, vol/issue not yet assigned in the PDF ("15481425, 0" footer = vol 0);
% 2009 = page range read from the eScholarship PDF's printed "Page N"
% footers — confirm against the printed AnthroNotes 30.1 issue.
\item Boellstorff, Tom. ``Toward Anthropologies of the Metaverse.''
\textit{American Ethnologist}, 2023, pp.~1--10,
https://doi.org/10.1111/amet.13228. (****** UNVERIFIED: online
version-of-record pagination; confirm final vol./no./pages.)

\item ---. ``Virtual Worlds and Futures of Anthropology.''
\textit{AnthroNotes}, vol.~30, no.~1, 2009, pp.~1--5. (****** verify page
range against the printed issue.)

\item Burke, Kenneth. \textit{A Rhetoric of Motives}. 1950. U of
California P, 1969. % [PDF imprint page]

\item Calleja, Gordon. \textit{In-Game: From Immersion to Incorporation}.
MIT P, 2011. % [M1]

\item Chalmers, David J. ``The Virtual and the Real.''
\textit{Disputatio}, vol.~9, no.~46, 2017, pp.~309--52.
% [PDF BIBLID — year settled 2017; corpus PDF filename says 2016]

\item Chow, Meyrick. ``Determinants of Presence in 3D Virtual Worlds: A
Structural Equation Modelling Analysis.'' \textit{Australasian Journal of
Educational Technology}, vol.~32, no.~1, 2016, pp.~1--18.
% [IDX — venue pin resolved]

\item Cook, David A., et al. ``Technology-Enhanced Simulation for Health
Professions Education: A Systematic Review and Meta-Analysis.''
\textit{JAMA}, vol.~306, no.~9, 2011, pp.~978--88. % [M1+IDX]

\item Cossio, Samantha, et al. ``Cybersickness and Discomfort from
Head-Mounted Displays Delivering Fully Immersive Virtual Reality: A
Systematic Review.'' \textit{Nurse Education in Practice}, vol.~85, 2025,
104376. % [PDF byline — first-name pin resolved]

\item Dalgarno, Barney, and Mark J.\ W.\ Lee. ``What Are the Learning
Affordances of 3-D Virtual Environments?'' \textit{British Journal of
Educational Technology}, vol.~41, no.~1, 2010, pp.~10--32. % [M1+IDX]

\item ---. ``Exploring the Relationship between Afforded Learning Tasks
and Learning Benefits in 3D Virtual Learning Environments.''
\textit{Future Challenges, Sustainable Futures: Proceedings ascilite
Wellington 2012}, ascilite, 2012, pp.~236--45.
% [M1 — pages inferred from proceedings range, per M1 header note]

\item De Back, Tycho T., Angelica M.\ Tinga, and Max M.\ Louwerse.
``Learning in Immersed Collaborative Virtual Environments: Design and
Implementation.'' \textit{Interactive Learning Environments}, vol.~31,
no.~8, 2021, pp.~5364--82. % [IDX — online 2021; print issue 2023]

\item Dubovi, Ilana, Sharona T.\ Levy, and Efrat Dagan. ``Now I Know How!
The Learning Process of Medication Administration among Nursing Students
with Non-Immersive Desktop Virtual Reality Simulation.''
\textit{Computers \& Education}, vol.~113, 2017, pp.~16--27.
% [PDF byline — first names verified]

\item Fowler, Chris. ``Virtual Reality and Learning: Where Is the
Pedagogy?'' \textit{British Journal of Educational Technology}, vol.~46,
no.~2, 2015, pp.~412--22. % [M1+IDX]

\item Garrison, D.\ Randy, Terry Anderson, and Walter Archer. ``Critical
Inquiry in a Text-Based Environment: Computer Conferencing in Higher
Education.'' \textit{The Internet and Higher Education}, vol.~2,
no.~2--3, 1999, pp.~87--105. % [M1+IDX]

\item Guilford, W.\ H., M.\ Kotche, and R.\ H.\ Schmedlen. ``A Survey of
Clinical Immersion Experiences in Biomedical Engineering.''
\textit{Biomedical Engineering Education}, vol.~3, 2023, pp.~113--22.
% [REF — pilot article ref 4; initials as printed there]

\item Gunawardena, Charlotte N., and Frank J.\ Zittle. ``Social Presence
as a Predictor of Satisfaction within a Computer-Mediated Conferencing
Environment.'' \textit{American Journal of Distance Education}, vol.~11,
no.~3, 1997, pp.~8--26. % [M1+IDX]

\item Heidegger, Martin. \textit{Basic Concepts of Aristotelian Philosophy}.
Translated by Robert D. Metcalf and Mark B. Tanzer, Indiana UP, 2009.
% [M1 2.12 footnote; harmonize vs Part I at dissertation assembly]

\item ---. \textit{Being and Time}. Translated by John Macquarrie and Edward
Robinson, Blackwell, 1962.
% [M1 2.12 footnote; harmonize vs Part I at dissertation assembly]

\item ---. \textit{The Fundamental Concepts of Metaphysics:
World, Finitude, Solitude}. Translated by William McNeill and Nicholas
Walker, Indiana UP, 1995. % [queued add; trans./press per M4 ledger]

\item Heim, Michael. \textit{Virtual Realism}. Oxford UP, 1998.
% [PDF title page]

\item King, Christine E., Kadin Diec, and Dalton Salvo. ``Scalable
Virtual Reality Global Clinical Immersion for Culturally Responsive
Engineering Design Skills.'' \textit{2025 ASEE Annual Conference \&
Exposition}, ASEE, 2025, paper 49715.
% [queued add; year verified vs title-page copyright line, conference
%  name inferred from Paper ID + copyright — per M2 header note]

\item King, Christine E., and Dalton Salvo. ``Phenomenological Evaluation
of an Undergraduate Clinical Needs Finding Skills Through a Virtual
Reality Clinical Immersion Platform.'' \textit{Biomedical Engineering
Education}, vol.~4, 2024, pp.~381--97,
https://doi.org/10.1007/s43683-024-00139-5. % [M0 footnote cite]

\item Kononowicz, Andrzej A., et al. ``Virtual Patient Simulations in
Health Professions Education: Systematic Review and Meta-Analysis by the
Digital Health Education Collaboration.'' \textit{Journal of Medical
Internet Research}, vol.~21, no.~7, 2019, e14676. % [M1+IDX]

\item Kreijns, Karel, Paul A.\ Kirschner, and Wim Jochems. ``Identifying
the Pitfalls for Social Interaction in Computer-Supported Collaborative
Learning Environments: A Review of the Research.'' \textit{Computers in
Human Behavior}, vol.~19, no.~3, 2003, pp.~335--53. % [M1; issue no. IDX]

\item Lombard, Matthew, and Theresa Ditton. ``At the Heart of It All: The
Concept of Presence.'' \textit{Journal of Computer-Mediated
Communication}, vol.~3, no.~2, 1997. % [M1+IDX — unpaginated journal]

\item Makransky, Guido, and Lasse Lilleholt. ``A Structural Equation
Modeling Investigation of the Emotional Value of Immersive Virtual
Reality in Education.'' \textit{Educational Technology Research and
Development}, vol.~66, no.~5, 2018, pp.~1141--64.
% [IDX — named in M0 presence-definition footnote; OPTIONAL entry, user
%  may cut if name-mentions are not to enter Works Cited]

\item Makransky, Guido, and Gustav B{\o}g Petersen. ``Investigating the
Process of Learning with Desktop Virtual Reality: A Structural Equation
Modeling Approach.'' \textit{Computers \& Education}, vol.~134, 2019,
pp.~15--30. % [M1+IDX]

\item Makransky, Guido, Richard E.\ Mayer, Nicola Veitch, Michelle Hood,
Karl B.\ Christensen, and Helle Gadegaard. ``Equivalence of Using a
Desktop Virtual Reality Science Simulation at Home and in Class.''
\textit{PLOS ONE}, vol.~14, no.~4, 2019, e0214944.
% [IDX initials; first names ****** UNVERIFIED beyond Makransky/Mayer —
%  confirm against PDF byline at final typeset]

\item Makransky, Guido, Thomas S.\ Terkildsen, and Richard E.\ Mayer.
``Adding Immersive Virtual Reality to a Science Lab Simulation Causes
More Presence but Less Learning.'' \textit{Learning and Instruction},
vol.~60, 2019, pp.~225--36. % [M1+IDX]

\item Mark, Gloria. \textit{Attention Span: ****** UNVERIFIED subtitle}.
****** UNVERIFIED publisher, 2023.
% [corpus holds quotation-notes PDF only, no imprint page; commonly
%  reported as Hanover Square Press with subtitle "A Groundbreaking Way
%  to Restore Balance, Happiness and Productivity" — VERIFY against the
%  physical book before final typeset]

\item Mayer, Richard E., Guido Makransky, and Jocelyn Parong. ``The
Promise and Pitfalls of Learning in Immersive Virtual Reality.''
\textit{International Journal of Human-Computer Interaction}, vol.~39,
no.~11, 2022, pp.~2229--38. % [M1+IDX — online VoR 2022; print 2023]

\item Merchant, Zahira, et al. ``The Learner Characteristics, Features of
Desktop 3D Virtual Reality Environments, and College Chemistry
Instruction: A Structural Equation Modeling Analysis.''
\textit{Computers \& Education}, vol.~59, no.~2, 2012, pp.~551--68.
% [M1; issue no. IDX]

\item Merleau-Ponty, Maurice. \textit{Phenomenology of Perception}. 1945.
Translated by Donald A.\ Landes, Routledge, 2014.
% [queued add; edition per M3.4 ledger — user's own edition]

\item Mikropoulos, T.\ A., and A.\ Natsis. ``Educational Virtual
Environments: A Ten-Year Review of Empirical Research (1999--2009).''
\textit{Computers \& Education}, vol.~56, 2011, pp.~769--80.
% [M1 — initials kept until verified; cited qtd.-in-Fowler, so this
%  entry is OPTIONAL under strict MLA indirect-source practice]

\item Moravec, K.\ R., E.\ L.\ Lothamer, A.\ Hoene, P.\ M.\ Wagoner,
D.\ J.\ Beckman, and C.\ J.\ Goergen. ``Clinical Immersion of
Undergraduate Biomedical Engineering Students: Best Practices for
Short-Term Programs.'' \textit{Biomedical Engineering Education},
vol.~3, 2023, pp.~217--23. % [REF — pilot article ref 6]

\item Parong, Jocelyn, and Richard E.\ Mayer. ``Cognitive and Affective
Processes for Learning Science in Immersive Virtual Reality.''
\textit{Journal of Computer Assisted Learning}, vol.~37, no.~1, 2021,
pp.~226--41. % [IDX — pages pin resolved; VoR year 2021]

\item ---. ``Learning Science in Immersive Virtual Reality.''
\textit{Journal of Educational Psychology}, vol.~110, no.~6, 2018,
pp.~785--97.
% [M1 bank values; ****** UNVERIFIED: corpus PDF is the advance-online
%  version without final pagination — confirm 110(6):785–97 against the
%  version of record before final typeset]

\item Radianti, Jaziar, et al. ``A Systematic Review of Immersive Virtual
Reality Applications for Higher Education: Design Elements, Lessons
Learned, and Research Agenda.'' \textit{Computers \& Education},
vol.~147, 2020, 103778. % [M1; subtitle pin resolved IDX]

\item Richardson, Jennifer C., and Karen Swan. ``Examining Social
Presence in Online Courses in Relation to Students' Perceived Learning
and Satisfaction.'' \textit{Journal of Asynchronous Learning Networks},
vol.~7, no.~1, 2003, pp.~68--88. % [IDX — venue pin resolved]

\item Richardson, Jennifer C., Yukiko Maeda, Jing Lv, and Secil Caskurlu.
``Social Presence in Relation to Students' Satisfaction and Learning in
the Online Environment: A Meta-Analysis.'' \textit{Computers in Human
Behavior}, vol.~71, 2017, pp.~402--17. % [IDX — venue pin resolved]

\item Short, John, Ederyn Williams, and Bruce Christie. \textit{The
Social Psychology of Telecommunications}. Wiley, 1976. % [M1]

\item Terry, Krista P., and Peter E.\ Doolittle. ``Re-examining Social
Presence: Implications for Digital Pedagogies.'' \textit{Technology,
Instruction, Cognition and Learning}, vol.~11, 2019, pp.~121--39.
% [M1+IDX]

\item Tu, Chih-Hsiung. ``The Measurement of Social Presence in an Online
Learning Environment.'' \textit{International Journal on E-Learning},
****** UNVERIFIED volume/issue, 2002, pp.~34--45.
% [IDX carries the same UNVERIFIED marker; commonly vol. 1, no. 2 —
%  confirm against the article PDF before final typeset]

\item van der Meer, Nesse, Vivian van der Werf, Willem-Paul Brinkman,
and Marcus Specht. ``Virtual Reality and Collaborative Learning: A
Systematic Literature Review.'' \textit{Frontiers in Virtual Reality},
vol.~4, 2023, 1159905. % [IDX — names/title pin resolved]

\item Wendt, Thomas. \textit{Design for Dasein: Understanding the Design
of Experiences}. Self-published, 2015.
% [PDF title/imprint pages — subtitle plural "Experiences" per title
%  page; ISBN 978-1506166537, © 2015 Thomas Wendt, no press named]

\item Witmer, Bob G., and Michael J.\ Singer. ``Measuring Presence in
Virtual Environments: A Presence Questionnaire.'' \textit{Presence:
Teleoperators and Virtual Environments}, vol.~7, no.~3, 1998,
pp.~225--40. % [M1+IDX]

\end{workscited}

\end{document}
